% Options for packages loaded elsewhere
% Options for packages loaded elsewhere
\PassOptionsToPackage{unicode}{hyperref}
\PassOptionsToPackage{hyphens}{url}
\PassOptionsToPackage{dvipsnames,svgnames,x11names}{xcolor}
%
\documentclass[
  letterpaper,
  DIV=11,
  numbers=noendperiod]{scrreprt}
\usepackage{xcolor}
\usepackage{amsmath,amssymb}
\setcounter{secnumdepth}{5}
\usepackage{iftex}
\ifPDFTeX
  \usepackage[T1]{fontenc}
  \usepackage[utf8]{inputenc}
  \usepackage{textcomp} % provide euro and other symbols
\else % if luatex or xetex
  \usepackage{unicode-math} % this also loads fontspec
  \defaultfontfeatures{Scale=MatchLowercase}
  \defaultfontfeatures[\rmfamily]{Ligatures=TeX,Scale=1}
\fi
\usepackage{lmodern}
\ifPDFTeX\else
  % xetex/luatex font selection
\fi
% Use upquote if available, for straight quotes in verbatim environments
\IfFileExists{upquote.sty}{\usepackage{upquote}}{}
\IfFileExists{microtype.sty}{% use microtype if available
  \usepackage[]{microtype}
  \UseMicrotypeSet[protrusion]{basicmath} % disable protrusion for tt fonts
}{}
\makeatletter
\@ifundefined{KOMAClassName}{% if non-KOMA class
  \IfFileExists{parskip.sty}{%
    \usepackage{parskip}
  }{% else
    \setlength{\parindent}{0pt}
    \setlength{\parskip}{6pt plus 2pt minus 1pt}}
}{% if KOMA class
  \KOMAoptions{parskip=half}}
\makeatother
% Make \paragraph and \subparagraph free-standing
\makeatletter
\ifx\paragraph\undefined\else
  \let\oldparagraph\paragraph
  \renewcommand{\paragraph}{
    \@ifstar
      \xxxParagraphStar
      \xxxParagraphNoStar
  }
  \newcommand{\xxxParagraphStar}[1]{\oldparagraph*{#1}\mbox{}}
  \newcommand{\xxxParagraphNoStar}[1]{\oldparagraph{#1}\mbox{}}
\fi
\ifx\subparagraph\undefined\else
  \let\oldsubparagraph\subparagraph
  \renewcommand{\subparagraph}{
    \@ifstar
      \xxxSubParagraphStar
      \xxxSubParagraphNoStar
  }
  \newcommand{\xxxSubParagraphStar}[1]{\oldsubparagraph*{#1}\mbox{}}
  \newcommand{\xxxSubParagraphNoStar}[1]{\oldsubparagraph{#1}\mbox{}}
\fi
\makeatother

\usepackage{color}
\usepackage{fancyvrb}
\newcommand{\VerbBar}{|}
\newcommand{\VERB}{\Verb[commandchars=\\\{\}]}
\DefineVerbatimEnvironment{Highlighting}{Verbatim}{commandchars=\\\{\}}
% Add ',fontsize=\small' for more characters per line
\usepackage{framed}
\definecolor{shadecolor}{RGB}{241,243,245}
\newenvironment{Shaded}{\begin{snugshade}}{\end{snugshade}}
\newcommand{\AlertTok}[1]{\textcolor[rgb]{0.68,0.00,0.00}{#1}}
\newcommand{\AnnotationTok}[1]{\textcolor[rgb]{0.37,0.37,0.37}{#1}}
\newcommand{\AttributeTok}[1]{\textcolor[rgb]{0.40,0.45,0.13}{#1}}
\newcommand{\BaseNTok}[1]{\textcolor[rgb]{0.68,0.00,0.00}{#1}}
\newcommand{\BuiltInTok}[1]{\textcolor[rgb]{0.00,0.23,0.31}{#1}}
\newcommand{\CharTok}[1]{\textcolor[rgb]{0.13,0.47,0.30}{#1}}
\newcommand{\CommentTok}[1]{\textcolor[rgb]{0.37,0.37,0.37}{#1}}
\newcommand{\CommentVarTok}[1]{\textcolor[rgb]{0.37,0.37,0.37}{\textit{#1}}}
\newcommand{\ConstantTok}[1]{\textcolor[rgb]{0.56,0.35,0.01}{#1}}
\newcommand{\ControlFlowTok}[1]{\textcolor[rgb]{0.00,0.23,0.31}{\textbf{#1}}}
\newcommand{\DataTypeTok}[1]{\textcolor[rgb]{0.68,0.00,0.00}{#1}}
\newcommand{\DecValTok}[1]{\textcolor[rgb]{0.68,0.00,0.00}{#1}}
\newcommand{\DocumentationTok}[1]{\textcolor[rgb]{0.37,0.37,0.37}{\textit{#1}}}
\newcommand{\ErrorTok}[1]{\textcolor[rgb]{0.68,0.00,0.00}{#1}}
\newcommand{\ExtensionTok}[1]{\textcolor[rgb]{0.00,0.23,0.31}{#1}}
\newcommand{\FloatTok}[1]{\textcolor[rgb]{0.68,0.00,0.00}{#1}}
\newcommand{\FunctionTok}[1]{\textcolor[rgb]{0.28,0.35,0.67}{#1}}
\newcommand{\ImportTok}[1]{\textcolor[rgb]{0.00,0.46,0.62}{#1}}
\newcommand{\InformationTok}[1]{\textcolor[rgb]{0.37,0.37,0.37}{#1}}
\newcommand{\KeywordTok}[1]{\textcolor[rgb]{0.00,0.23,0.31}{\textbf{#1}}}
\newcommand{\NormalTok}[1]{\textcolor[rgb]{0.00,0.23,0.31}{#1}}
\newcommand{\OperatorTok}[1]{\textcolor[rgb]{0.37,0.37,0.37}{#1}}
\newcommand{\OtherTok}[1]{\textcolor[rgb]{0.00,0.23,0.31}{#1}}
\newcommand{\PreprocessorTok}[1]{\textcolor[rgb]{0.68,0.00,0.00}{#1}}
\newcommand{\RegionMarkerTok}[1]{\textcolor[rgb]{0.00,0.23,0.31}{#1}}
\newcommand{\SpecialCharTok}[1]{\textcolor[rgb]{0.37,0.37,0.37}{#1}}
\newcommand{\SpecialStringTok}[1]{\textcolor[rgb]{0.13,0.47,0.30}{#1}}
\newcommand{\StringTok}[1]{\textcolor[rgb]{0.13,0.47,0.30}{#1}}
\newcommand{\VariableTok}[1]{\textcolor[rgb]{0.07,0.07,0.07}{#1}}
\newcommand{\VerbatimStringTok}[1]{\textcolor[rgb]{0.13,0.47,0.30}{#1}}
\newcommand{\WarningTok}[1]{\textcolor[rgb]{0.37,0.37,0.37}{\textit{#1}}}

\usepackage{longtable,booktabs,array}
\newcounter{none} % for unnumbered tables
\usepackage{calc} % for calculating minipage widths
% Correct order of tables after \paragraph or \subparagraph
\usepackage{etoolbox}
\makeatletter
\patchcmd\longtable{\par}{\if@noskipsec\mbox{}\fi\par}{}{}
\makeatother
% Allow footnotes in longtable head/foot
\IfFileExists{footnotehyper.sty}{\usepackage{footnotehyper}}{\usepackage{footnote}}
\makesavenoteenv{longtable}
\usepackage{graphicx}
\makeatletter
\newsavebox\pandoc@box
\newcommand*\pandocbounded[1]{% scales image to fit in text height/width
  \sbox\pandoc@box{#1}%
  \Gscale@div\@tempa{\textheight}{\dimexpr\ht\pandoc@box+\dp\pandoc@box\relax}%
  \Gscale@div\@tempb{\linewidth}{\wd\pandoc@box}%
  \ifdim\@tempb\p@<\@tempa\p@\let\@tempa\@tempb\fi% select the smaller of both
  \ifdim\@tempa\p@<\p@\scalebox{\@tempa}{\usebox\pandoc@box}%
  \else\usebox{\pandoc@box}%
  \fi%
}
% Set default figure placement to htbp
\def\fps@figure{htbp}
\makeatother





\setlength{\emergencystretch}{3em} % prevent overfull lines

\providecommand{\tightlist}{%
  \setlength{\itemsep}{0pt}\setlength{\parskip}{0pt}}



 


\KOMAoption{captions}{tableheading}
\makeatletter
\@ifpackageloaded{bookmark}{}{\usepackage{bookmark}}
\makeatother
\makeatletter
\@ifpackageloaded{caption}{}{\usepackage{caption}}
\AtBeginDocument{%
\ifdefined\contentsname
  \renewcommand*\contentsname{Table of contents}
\else
  \newcommand\contentsname{Table of contents}
\fi
\ifdefined\listfigurename
  \renewcommand*\listfigurename{List of Figures}
\else
  \newcommand\listfigurename{List of Figures}
\fi
\ifdefined\listtablename
  \renewcommand*\listtablename{List of Tables}
\else
  \newcommand\listtablename{List of Tables}
\fi
\ifdefined\figurename
  \renewcommand*\figurename{Figure}
\else
  \newcommand\figurename{Figure}
\fi
\ifdefined\tablename
  \renewcommand*\tablename{Table}
\else
  \newcommand\tablename{Table}
\fi
}
\@ifpackageloaded{float}{}{\usepackage{float}}
\floatstyle{ruled}
\@ifundefined{c@chapter}{\newfloat{codelisting}{h}{lop}}{\newfloat{codelisting}{h}{lop}[chapter]}
\floatname{codelisting}{Listing}
\newcommand*\listoflistings{\listof{codelisting}{List of Listings}}
\makeatother
\makeatletter
\makeatother
\makeatletter
\@ifpackageloaded{caption}{}{\usepackage{caption}}
\@ifpackageloaded{subcaption}{}{\usepackage{subcaption}}
\makeatother
\usepackage{bookmark}
\IfFileExists{xurl.sty}{\usepackage{xurl}}{} % add URL line breaks if available
\urlstyle{same}
\hypersetup{
  pdftitle={Random Variable untuk Pengambilan Keputusan},
  pdfauthor={Armein Z. R. Langi},
  colorlinks=true,
  linkcolor={blue},
  filecolor={Maroon},
  citecolor={Blue},
  urlcolor={Blue},
  pdfcreator={LaTeX via pandoc}}


\title{Random Variable untuk Pengambilan Keputusan}
\usepackage{etoolbox}
\makeatletter
\providecommand{\subtitle}[1]{% add subtitle to \maketitle
  \apptocmd{\@title}{\par {\large #1 \par}}{}{}
}
\makeatother
\subtitle{Probabilitas Statistika untuk Mahasiswa Generasi Z}
\author{Armein Z. R. Langi}
\date{2026-03-31}
\begin{document}
\maketitle

\renewcommand*\contentsname{Table of contents}
{
\setcounter{tocdepth}{2}
\tableofcontents
}

\bookmarksetup{startatroot}

\chapter{Random Variable untuk Pengambilan
Keputusan}\label{random-variable-untuk-pengambilan-keputusan}

Probabilitas Statistika untuk Mahasiswa Generasi Z

\hfill\break

\section{Selamat Datang}\label{selamat-datang}

Buku ini ditulis untuk mendampingi mahasiswa \textbf{II-2111
Probabilitas Statistika} memahami topik \textbf{Random Variable} dengan
cara yang lebih hidup, lebih aplikatif, dan lebih dekat dengan keputusan
nyata.

Alih-alih memulai dari definisi yang terasa jauh, buku ini mengajak
pembaca masuk dari pertanyaan yang lebih akrab:

\begin{itemize}
\tightlist
\item
  Bagaimana membuat keputusan ketika masa depan belum pasti?
\item
  Bagaimana membandingkan dua pilihan yang sama-sama mungkin
  menguntungkan?
\item
  Bagaimana memahami risiko, bukan hanya rata-rata?
\item
  Bagaimana menggunakan Python untuk melihat pola, mencoba simulasi, dan
  memeriksa logika?
\end{itemize}

Buku ini dibangun dengan tiga semangat utama:

\begin{enumerate}
\def\labelenumi{\arabic{enumi}.}
\item
  \textbf{Decision-making oriented}\\
  Probabilitas dipelajari sebagai alat untuk memilih dengan lebih
  bijaksana.
\item
  \textbf{Python first for intuition}\\
  Mahasiswa diajak mendapatkan \emph{quick wins} lebih dulu melalui
  simulasi, visualisasi, dan eksperimen.
\item
  \textbf{Logic with meaning}\\
  Definisi, rumus, dan teori tetap penting, tetapi selalu dihubungkan
  dengan makna dan konteks nyata.
\end{enumerate}

\section{Cara Menggunakan Buku Ini}\label{cara-menggunakan-buku-ini}

Setiap bab dalam buku ini dirancang agar bisa dipelajari dengan pola:

\begin{enumerate}
\def\labelenumi{\arabic{enumi}.}
\tightlist
\item
  \textbf{Mulai dari konteks masalah}
\item
  \textbf{Lihat quick win dengan Python}
\item
  \textbf{Bangun intuisi}
\item
  \textbf{Masuk ke model formal}
\item
  \textbf{Jawab pertanyaan berbasis model}
\item
  \textbf{Tarik implikasi keputusan}
\end{enumerate}

Pola ini disebut pendekatan \textbf{KMQA}:

\begin{itemize}
\tightlist
\item
  \textbf{K} --- Konteks\\
\item
  \textbf{M} --- Model\\
\item
  \textbf{Q} --- Questions\\
\item
  \textbf{A} --- Apply
\end{itemize}

\section{Struktur Buku}\label{struktur-buku}

Buku ini terdiri dari enam bab utama:

\begin{enumerate}
\def\labelenumi{\arabic{enumi}.}
\item
  \textbf{Pendahuluan Pengambilan Keputusan}\\
  Mengaitkan probabilitas dengan keputusan nyata, risiko, dan Python.
\item
  \textbf{Random Variable Umum}\\
  Fondasi formal: range, diskrit vs kontinu, PMF, CDF, PDF, ekspektasi,
  dan varians.
\item
  \textbf{Distribusi Random Variable Diskrit}\\
  Uniform diskrit, Bernoulli, Binomial, Geometric, Poisson, dan
  distribusi empiris dari histogram.
\item
  \textbf{Distribusi Random Variable Kontinu}\\
  Uniform kontinu, Normal, Gamma, Exponential, Erlang, Weibull, Pareto,
  Chi-Square, dan hubungan antardistribusi.
\item
  \textbf{Random Variable Multivariat dan Fungsi Random Variable}\\
  Joint, marginal, conditional, covariance, correlation, independensi,
  dan transformasi peubah acak.
\item
  \textbf{Penutup}\\
  Ringkasan ide besar buku, checklist kompetensi, dan peneguhan bahwa
  probabilitas adalah cara berpikir di bawah ketidakpastian.
\end{enumerate}

\section{Untuk Siapa Buku Ini}\label{untuk-siapa-buku-ini}

Buku ini terutama ditujukan untuk mahasiswa tingkat dua, khususnya yang:
- ingin memahami probabilitas secara konseptual, - ingin melihat contoh
nyata dan relevan, - ingin belajar Python sambil tetap berpikir, - dan
ingin merasakan bahwa statistika bukan sekadar soal ujian, tetapi alat
untuk hidup dan bekerja dengan lebih cerdas.

\section{Harapan}\label{harapan}

Semoga buku ini membantu Anda mengalami sesuatu yang penting: bahwa
probabilitas bukanlah kumpulan simbol yang dingin, melainkan bahasa yang
sangat berguna saat kita harus berpikir, memilih, dan bertindak di dunia
yang tidak sepenuhnya pasti.

Silakan mulai dari \textbf{Kata Pengantar}, lalu masuk ke \textbf{Bab
1}.

\bookmarksetup{startatroot}

\chapter{Kata Pengantar}\label{kata-pengantar}

Buku ini disusun untuk menemani para mahasiswa kuliah \textbf{II-2111
Probabilitas Statistika} belajar, berpikir, dan bertumbuh. Mungkin
muncul pertanyaan: di tengah begitu banyak buku, video, dan sumber
belajar yang tersedia, mengapa masih perlu ada buku ini?

Jawaban paling jujurnya sederhana: karena belajar bukan hanya soal
tersedianya materi, tetapi soal \textbf{bagaimana hati dan pikiran kita
digerakkan untuk mau memahami}. Buku ini lahir dari keinginan untuk
membantu mahasiswa, khususnya generasi Z, agar tidak hanya ``mengikuti
kuliah'' atau ``mengerjakan soal'', tetapi sungguh-sungguh melihat bahwa
probabilitas dan statistika adalah ilmu yang hidup, dekat, dan berguna
untuk mengambil keputusan dalam dunia nyata.

Buku ini dibangun di atas dua keyakinan. Pertama, bahwa materi
probabilitas dan statistika akan menjadi lebih bermakna bila dibingkai
sebagai bagian dari \textbf{pengambilan keputusan} atau inferensi.
Kedua, bahwa pada zaman ini, pembelajaran akan lebih kuat bila didukung
oleh \textbf{Python} sebagai alat untuk mencoba, bereksperimen,
mensimulasikan, dan memahami.

Pada dasarnya, hidup manusia penuh dengan keputusan. Setiap hari kita
memilih, menimbang, memperkirakan, berharap, dan menghadapi konsekuensi.
Dalam pengertian itu, pengambilan keputusan adalah salah satu wujud
paling nyata dari kecerdasan manusia. Kita tidak disebut cerdas hanya
karena tahu banyak hal, tetapi karena mampu menggunakan pengetahuan itu
untuk memilih jalan yang lebih baik.

Menariknya, bentuk pengambilan keputusan tidak jauh berbeda dari
menjawab soal ujian. Ada yang sederhana seperti \textbf{true-false}, ada
yang berupa \textbf{pilihan ganda}, ada yang berbentuk isian singkat,
jawaban satu kalimat, satu alinea, esai, bahkan disertasi. Semua itu
pada dasarnya adalah latihan untuk memilih jawaban terbaik berdasarkan
informasi yang tersedia. Bedanya, dalam kehidupan nyata, keputusan tidak
hanya dinilai benar atau salah, tetapi juga membawa akibat yang sungguh
kita rasakan.

Karena itu, untuk mengambil keputusan yang baik, kita memerlukan
\textbf{model}. Model membantu kita menyederhanakan kenyataan tanpa
kehilangan makna yang penting. Dan model dibangun oleh
\textbf{variabel}.

Kita mengenal variabel sebagai besaran yang nilainya dapat berubah.
Misalnya, panjang dan lebar sebidang tanah kita nyatakan dengan \(p\)
dan \(l\). Luas tanah kemudian dapat dimodelkan sebagai

\[
L = p \times l.
\]

Tentu luas tanah yang sesungguhnya bukanlah huruf \(L\), melainkan suatu
kenyataan fisik. Namun variabel itu menolong kita untuk memikirkan
kenyataan tersebut secara lebih teratur. Dengan cara itulah manusia
membangun ilmu: bukan dengan langsung menggenggam realitas sepenuhnya,
tetapi dengan membuat model yang cukup baik untuk memahami dan
bertindak.

Ada besaran yang ketika diukur hasilnya relatif tetap. Misalnya, luas
meja atau panjang sebuah kabel. Tetapi ada juga besaran yang ketika
diukur berulang kali hasilnya berubah-ubah. Misalnya, jumlah kendaraan
yang lewat di gerbang tol antara pukul 09.00--10.00, jumlah pasien yang
datang ke klinik dalam satu hari, atau skor pertandingan sepak bola.
Besaran seperti inilah yang kita modelkan sebagai \textbf{variabel acak}
(\emph{random variable}).

Jika \(X\) adalah variabel yang memodelkan jumlah kendaraan yang
melewati suatu gerbang tol dalam rentang waktu tertentu, dan nilainya
dapat berbeda dari satu hari ke hari berikutnya, maka \(X\) adalah
variabel acak. Variabel acak memiliki \textbf{range}, yaitu himpunan
nilai-nilai yang mungkin diambilnya. Jika nilai-nilainya berupa bilangan
diskrit, seperti 0, 1, 2, 3, dan seterusnya, maka kita menyebutnya
variabel acak diskrit.

Dalam praktik, banyak pengukuran dilakukan secara diskrit karena kita
menggunakan satuan tertentu. Namun secara teoritis, ada besaran yang
lebih mudah dipahami sebagai kontinu. Waktu adalah contoh yang baik.
Dalam praktik kita mengukurnya dalam tahun, bulan, hari, jam, menit, dan
detik. Namun bila ketelitian itu terus diperhalus, kita dapat
memandangnya sebagai kontinu. Pemilihan model diskrit atau kontinu bukan
sekadar soal teknik, tetapi juga soal cara berpikir yang paling tepat
untuk memahami suatu persoalan.

Hal penting yang perlu disadari adalah ini: sebelum suatu peristiwa
terjadi, kita biasanya tidak tahu pasti berapa nilai aktual dari \(X\).
Yang kita miliki hanyalah pengetahuan tentang
kemungkinan-kemungkinannya. Namun mengetahui range saja belum cukup
untuk membuat keputusan. Kita membutuhkan sesuatu yang lebih kuat, yaitu
ukuran yang mewakili kecenderungan nilai-nilai yang mungkin muncul. Di
sinilah \textbf{ekspektasi} menjadi sangat penting. Nilai \(E[X]\) dapat
dipandang sebagai tebakan terbaik kita terhadap \(X\) sebelum kenyataan
terjadi.

Dari titik inilah probabilitas menjadi sangat berarti. Probabilitas
menolong kita memberi bobot pada kemungkinan-kemungkinan. Dengan
probabilitas, kita tidak hanya tahu apa saja yang mungkin terjadi,
tetapi juga seberapa masuk akal masing-masing kemungkinan itu. Dari sana
kita dapat menghitung ekspektasi, memahami risiko, dan akhirnya
mengambil keputusan dengan dasar yang lebih bijaksana.

Bayangkan, misalnya, seseorang harus memilih apakah pesta pernikahannya
akan diadakan \textbf{indoor} atau \textbf{outdoor}. Pilihan itu tampak
sederhana. Namun begitu muncul pertanyaan, ``Bagaimana kalau hujan?'',
keputusan itu berubah. Tiba-tiba kita sadar bahwa memilih bukan hanya
soal selera, tetapi juga soal menghadapi ketidakpastian. Dalam banyak
hal, demikianlah hidup bekerja. Kita jarang mengambil keputusan dalam
keadaan serba pasti. Justru karena ada ketidakpastian, kita perlu
probabilitas dan statistika.

Maka, probabilitas dan statistika bukan sekadar kumpulan rumus. Keduanya
adalah cara berpikir. Keduanya mengajarkan kita untuk rendah hati di
hadapan ketidakpastian, tetapi juga berani bertindak dengan dasar yang
lebih baik. Keduanya menolong kita menyadari bahwa meskipun masa depan
tidak bisa diketahui sepenuhnya, masa depan tetap bisa dihadapi dengan
pikiran yang jernih.

Lalu, mengapa \textbf{Python}? Karena pada zaman ini, Python telah
menjadi salah satu bahasa utama untuk berbicara dengan komputer. Dengan
Python, kita bisa mencoba ide, membuat simulasi, memeriksa tebakan, dan
melihat konsep bekerja secara nyata. Memang, sekarang kita juga dapat
meminta AI memberikan jawaban. Namun bila kita hanya menerima jawaban
akhir, sering kali kita kehilangan proses belajar yang sesungguhnya.
Python memberi kita kesempatan untuk tetap berpikir, tetap
bereksperimen, tetap bertanya, dan tetap bertumbuh---bahkan saat kita
memanfaatkan AI.

Karena itu, buku ini berusaha memadukan \textbf{konsep},
\textbf{intuisi}, \textbf{pengambilan keputusan}, dan \textbf{eksperimen
komputasional}. Harapannya, mahasiswa tidak hanya mampu mengerjakan soal
ujian, tetapi juga mampu melihat makna di balik soal itu: bahwa setiap
konsep yang dipelajari sesungguhnya sedang melatih kita menjadi pribadi
yang lebih cermat, lebih rasional, dan lebih bertanggung jawab dalam
memilih.

Akhirnya, buku ini sendiri adalah hasil dari sebuah keputusan: keputusan
untuk menyusun bahan ajar yang lebih dekat, lebih hidup, dan lebih
bersahabat bagi para pembaca. Semoga keputusan ini adalah keputusan yang
baik. Dan lebih daripada itu, semoga buku ini tidak hanya menambah
pengetahuan, tetapi juga menumbuhkan keberanian untuk berpikir, mencoba,
dan terus belajar.

\textbf{Bandung, 31 Maret 2026}\\
\textbf{Armein Z. R. Langi}

\bookmarksetup{startatroot}

\chapter{Bab 1. Pendahuluan Pengambilan
Keputusan}\label{bab-1.-pendahuluan-pengambilan-keputusan}

\section{Tujuan Bab}\label{tujuan-bab}

Setelah mempelajari bab ini, mahasiswa diharapkan mampu:

\begin{enumerate}
\def\labelenumi{\arabic{enumi}.}
\tightlist
\item
  melihat probabilitas dan statistika sebagai \textbf{alat pengambilan
  keputusan},
\item
  memahami mengapa \textbf{random variable} diperlukan untuk memodelkan
  dunia yang tidak pasti,
\item
  mengenali peran \textbf{ekspektasi} dan \textbf{varians} dalam
  membandingkan pilihan,
\item
  menggunakan \textbf{Python} untuk mendapatkan \emph{quick wins}
  sebelum masuk ke logika formal,
\item
  menghubungkan konsep probabilitas dengan kasus nyata di bidang teknik,
  bisnis, dan layanan.
\end{enumerate}

\section{Pembuka}\label{pembuka}

Banyak mahasiswa pertama kali bertemu probabilitas dan statistika
sebagai kumpulan simbol, rumus, dan notasi yang terasa dingin. Padahal,
di balik semua itu, probabilitas dan statistika sesungguhnya adalah
bahasa untuk menjawab pertanyaan yang sangat manusiawi:

\begin{itemize}
\tightlist
\item
  Apa keputusan yang paling masuk akal bila masa depan belum pasti?
\item
  Mana pilihan yang lebih aman?
\item
  Mana yang lebih menguntungkan?
\item
  Seberapa besar risiko yang harus kita tanggung?
\item
  Kapan kita harus berhenti mengandalkan intuisi dan mulai menghitung?
\end{itemize}

Bab ini mengajak Anda masuk ke probabilitas bukan dari pintu definisi
formal terlebih dahulu, tetapi dari pintu yang lebih akrab:
\textbf{pengambilan keputusan}.

\section{1.1 Dari Pilihan ke Model}\label{dari-pilihan-ke-model}

Setiap hari manusia membuat keputusan. Sebagian kecil: memilih rute ke
kampus, memilih tempat makan siang, memilih waktu belajar. Sebagian lain
lebih besar: memilih investasi, menentukan kebijakan produksi,
menetapkan garansi, menambah jumlah dokter, atau menentukan batas aman
komponen listrik.

Masalahnya, keputusan hampir selalu dibuat \textbf{sebelum} kita tahu
apa yang akan terjadi. Masa depan belum datang, data belum lengkap,
keadaan tidak sepenuhnya pasti. Karena itu, kita membutuhkan
\textbf{model}.

Model membantu kita menyederhanakan kenyataan agar bisa dipikirkan,
dihitung, dan diputuskan. Salah satu model paling penting dalam
probabilitas adalah \textbf{random variable}. Dengan random variable,
kejadian yang semula acak dipetakan menjadi angka. Begitu dunia acak
berhasil dipetakan ke bilangan, kita bisa mulai menghitung: - peluang, -
rata-rata jangka panjang, - seberapa liar penyimpangannya, - dan
konsekuensi keputusan.

\section{1.2 Quick Win: Melihat Mengapa Rata-Rata Saja Tidak
Cukup}\label{quick-win-melihat-mengapa-rata-rata-saja-tidak-cukup}

Mari mulai dari contoh yang sangat sederhana. Misalkan Anda harus
memilih satu dari dua opsi investasi:

\begin{itemize}
\tightlist
\item
  \textbf{A}: stabil, keuntungan kecil tetapi konsisten
\item
  \textbf{B}: bisa sangat menguntungkan, tetapi bisa juga sangat buruk
\end{itemize}

Secara kasatmata, dua opsi bisa saja memiliki rata-rata hasil yang sama.
Tetapi apakah itu berarti keduanya sama baik?

\begin{Shaded}
\begin{Highlighting}[]
\ImportTok{import}\NormalTok{ numpy }\ImportTok{as}\NormalTok{ np}
\ImportTok{import}\NormalTok{ matplotlib.pyplot }\ImportTok{as}\NormalTok{ plt}

\NormalTok{rng }\OperatorTok{=}\NormalTok{ np.random.default\_rng(}\DecValTok{42}\NormalTok{)}

\CommentTok{\#\# Dua investasi dengan mean hampir sama, tetapi variasi berbeda}
\NormalTok{A }\OperatorTok{=}\NormalTok{ rng.normal(loc}\OperatorTok{=}\DecValTok{10}\NormalTok{, scale}\OperatorTok{=}\DecValTok{2}\NormalTok{, size}\OperatorTok{=}\DecValTok{10000}\NormalTok{)}
\NormalTok{B }\OperatorTok{=}\NormalTok{ rng.normal(loc}\OperatorTok{=}\DecValTok{10}\NormalTok{, scale}\OperatorTok{=}\DecValTok{12}\NormalTok{, size}\OperatorTok{=}\DecValTok{10000}\NormalTok{)}

\BuiltInTok{print}\NormalTok{(}\StringTok{"Mean A ="}\NormalTok{, A.mean())}
\BuiltInTok{print}\NormalTok{(}\StringTok{"Var A  ="}\NormalTok{, A.var())}

\BuiltInTok{print}\NormalTok{(}\StringTok{"Mean B ="}\NormalTok{, B.mean())}
\BuiltInTok{print}\NormalTok{(}\StringTok{"Var B  ="}\NormalTok{, B.var())}
\end{Highlighting}
\end{Shaded}

\begin{verbatim}
Mean A = 9.979500249171975
Var A  = 4.050444191840732
Mean B = 10.24422094568761
Var B  = 144.83317214885628
\end{verbatim}

\begin{Shaded}
\begin{Highlighting}[]
\NormalTok{plt.figure(figsize}\OperatorTok{=}\NormalTok{(}\DecValTok{9}\NormalTok{, }\DecValTok{4}\NormalTok{))}
\NormalTok{plt.hist(A, bins}\OperatorTok{=}\DecValTok{60}\NormalTok{, alpha}\OperatorTok{=}\FloatTok{0.6}\NormalTok{, density}\OperatorTok{=}\VariableTok{True}\NormalTok{, label}\OperatorTok{=}\StringTok{"Investasi A"}\NormalTok{)}
\NormalTok{plt.hist(B, bins}\OperatorTok{=}\DecValTok{60}\NormalTok{, alpha}\OperatorTok{=}\FloatTok{0.6}\NormalTok{, density}\OperatorTok{=}\VariableTok{True}\NormalTok{, label}\OperatorTok{=}\StringTok{"Investasi B"}\NormalTok{)}
\NormalTok{plt.xlabel(}\StringTok{"Return"}\NormalTok{)}
\NormalTok{plt.ylabel(}\StringTok{"Density"}\NormalTok{)}
\NormalTok{plt.title(}\StringTok{"Dua investasi dengan mean hampir sama, variance sangat berbeda"}\NormalTok{)}
\NormalTok{plt.legend()}
\NormalTok{plt.grid(alpha}\OperatorTok{=}\FloatTok{0.3}\NormalTok{)}
\NormalTok{plt.show()}
\end{Highlighting}
\end{Shaded}

\pandocbounded{\includegraphics[keepaspectratio]{01-pendahuluan-pengambilan-keputusan_files/figure-pdf/cell-3-output-1.pdf}}

Dari hasil ini, kita segera mendapat \emph{quick win}:

\begin{itemize}
\tightlist
\item
  \textbf{mean} memberi gambaran hasil rata-rata,
\item
  \textbf{variance} memberi gambaran seberapa liar hasil itu,
\item
  keputusan yang matang tidak cukup melihat rata-rata, tetapi juga perlu
  melihat \textbf{risiko}.
\end{itemize}

Di sinilah probabilitas dan random variable mulai terasa hidup. Kita
tidak sedang belajar rumus demi rumus. Kita sedang belajar cara
membandingkan pilihan secara bertanggung jawab.

\section{1.3 Pola Umum Pengambilan
Keputusan}\label{pola-umum-pengambilan-keputusan}

Dalam buku ini, kita akan sering memakai pola berikut:

\subsection{K --- Konteks}\label{k-konteks}

Apa situasi nyatanya? Siapa pengambil keputusan? Apa kebutuhan dan
risikonya?

\subsection{M --- Model}\label{m-model}

Apa random variable yang relevan? Apa asumsi masuk akalnya? Apa
parameter pentingnya?

\subsection{Q --- Questions}\label{q-questions}

Apa pertanyaan yang benar-benar ingin dijawab? Peluang? Rata-rata?
Risiko bangkrut? Peluang cacat? Jumlah pelanggan? Waktu tunggu?

\subsection{A --- Apply}\label{a-apply}

Gunakan Python, simulasi, visualisasi, dan teori probabilitas untuk
menjawab pertanyaan dan menarik implikasi keputusan.

Pola ini akan kita sebut sebagai pendekatan \textbf{KMQA}:
\textbf{Konteks → Model → Questions → Apply}.

\section{1.4 Random Variable sebagai Bahasa Angka untuk Dunia
Acak}\label{random-variable-sebagai-bahasa-angka-untuk-dunia-acak}

Misalkan sebuah gerbang tol mencatat jumlah kendaraan yang lewat antara
pukul 09.00--10.00. Besaran itu tidak pasti. Setiap hari nilainya bisa
berbeda. Kita dapat menyatakan besaran ini sebagai sebuah random
variable:

\[
X = \text{jumlah kendaraan yang melewati gerbang tol antara pukul 09.00–10.00}
\]

Dengan begitu, kejadian nyata yang berubah-ubah kita petakan menjadi
bilangan. Begitu sudah menjadi bilangan, kita bisa bertanya: - berapa
nilai yang mungkin? - berapa peluang setiap nilai? - berapa
rata-ratanya? - seberapa besar variasinya?

Itulah inti random variable.

\section{1.5 Kasus 1 --- Ekspektasi dan Varians dalam Pilihan
Investasi}\label{kasus-1-ekspektasi-dan-varians-dalam-pilihan-investasi}

\subsection{K --- Konteks}\label{k-konteks-1}

Seorang mahasiswa mendapat dua opsi investasi kecil untuk dana
tabungannya:

\begin{itemize}
\tightlist
\item
  Investasi A: hasil tidak terlalu tinggi, tetapi relatif stabil.
\item
  Investasi B: kadang sangat tinggi, kadang turun tajam.
\end{itemize}

Ia bertanya: \emph{kalau saya ingin mengambil keputusan yang lebih
dewasa, ukuran apa yang harus saya lihat?}

\subsection{M --- Model}\label{m-model-1}

Misalkan: - \(X_A\) = return investasi A - \(X_B\) = return investasi B

Masing-masing adalah random variable. Dua ukuran dasar yang penting
adalah:

\[
E[X] = \text{ekspektasi}
\]

dan

\[
\mathrm{Var}(X) = E[(X - E[X])^2]
\]

Ekspektasi memberi gambaran hasil rata-rata jangka panjang. Varians
memberi ukuran seberapa jauh hasil itu bisa menyimpang.

\subsection{Q --- Questions}\label{q-questions-1}

\begin{enumerate}
\def\labelenumi{\arabic{enumi}.}
\tightlist
\item
  Mana investasi dengan rata-rata hasil lebih baik?
\item
  Mana investasi dengan risiko lebih besar?
\item
  Bila dua investasi punya rata-rata hampir sama, apakah variance dapat
  mengubah keputusan?
\end{enumerate}

\subsection{A --- Apply}\label{a-apply-1}

\begin{Shaded}
\begin{Highlighting}[]
\ImportTok{import}\NormalTok{ numpy }\ImportTok{as}\NormalTok{ np}

\NormalTok{rng }\OperatorTok{=}\NormalTok{ np.random.default\_rng(}\DecValTok{7}\NormalTok{)}

\NormalTok{XA }\OperatorTok{=}\NormalTok{ rng.choice([}\DecValTok{8}\NormalTok{, }\DecValTok{10}\NormalTok{, }\DecValTok{12}\NormalTok{], size}\OperatorTok{=}\DecValTok{100000}\NormalTok{, p}\OperatorTok{=}\NormalTok{[}\FloatTok{0.25}\NormalTok{, }\FloatTok{0.50}\NormalTok{, }\FloatTok{0.25}\NormalTok{])}
\NormalTok{XB }\OperatorTok{=}\NormalTok{ rng.choice([}\OperatorTok{{-}}\DecValTok{15}\NormalTok{, }\DecValTok{5}\NormalTok{, }\DecValTok{10}\NormalTok{, }\DecValTok{20}\NormalTok{, }\DecValTok{30}\NormalTok{], size}\OperatorTok{=}\DecValTok{100000}\NormalTok{, p}\OperatorTok{=}\NormalTok{[}\FloatTok{0.10}\NormalTok{, }\FloatTok{0.20}\NormalTok{, }\FloatTok{0.30}\NormalTok{, }\FloatTok{0.25}\NormalTok{, }\FloatTok{0.15}\NormalTok{])}

\BuiltInTok{print}\NormalTok{(}\StringTok{"E[XA] \textasciitilde{}"}\NormalTok{, XA.mean(), }\StringTok{" Var[XA] \textasciitilde{}"}\NormalTok{, XA.var())}
\BuiltInTok{print}\NormalTok{(}\StringTok{"E[XB] \textasciitilde{}"}\NormalTok{, XB.mean(), }\StringTok{" Var[XB] \textasciitilde{}"}\NormalTok{, XB.var())}
\end{Highlighting}
\end{Shaded}

\begin{verbatim}
E[XA] ~ 10.0018  Var[XA] ~ 2.0022367599999997
E[XB] ~ 12.0198  Var[XB] ~ 147.42240796
\end{verbatim}

\begin{Shaded}
\begin{Highlighting}[]
\NormalTok{plt.figure(figsize}\OperatorTok{=}\NormalTok{(}\DecValTok{9}\NormalTok{,}\DecValTok{4}\NormalTok{))}
\NormalTok{plt.hist(XA, bins}\OperatorTok{=}\NormalTok{np.arange(XA.}\BuiltInTok{min}\NormalTok{()}\OperatorTok{{-}}\FloatTok{0.5}\NormalTok{, XA.}\BuiltInTok{max}\NormalTok{()}\OperatorTok{+}\FloatTok{1.5}\NormalTok{), alpha}\OperatorTok{=}\FloatTok{0.6}\NormalTok{, density}\OperatorTok{=}\VariableTok{True}\NormalTok{, label}\OperatorTok{=}\StringTok{"XA"}\NormalTok{)}
\NormalTok{plt.hist(XB, bins}\OperatorTok{=}\NormalTok{np.arange(XB.}\BuiltInTok{min}\NormalTok{()}\OperatorTok{{-}}\FloatTok{0.5}\NormalTok{, XB.}\BuiltInTok{max}\NormalTok{()}\OperatorTok{+}\FloatTok{1.5}\NormalTok{), alpha}\OperatorTok{=}\FloatTok{0.6}\NormalTok{, density}\OperatorTok{=}\VariableTok{True}\NormalTok{, label}\OperatorTok{=}\StringTok{"XB"}\NormalTok{)}
\NormalTok{plt.xlabel(}\StringTok{"Return"}\NormalTok{)}
\NormalTok{plt.ylabel(}\StringTok{"Probability / Density"}\NormalTok{)}
\NormalTok{plt.title(}\StringTok{"Distribusi return dua investasi"}\NormalTok{)}
\NormalTok{plt.legend()}
\NormalTok{plt.grid(alpha}\OperatorTok{=}\FloatTok{0.3}\NormalTok{)}
\NormalTok{plt.show()}
\end{Highlighting}
\end{Shaded}

\pandocbounded{\includegraphics[keepaspectratio]{01-pendahuluan-pengambilan-keputusan_files/figure-pdf/cell-5-output-1.pdf}}

\subsubsection{Interpretasi}\label{interpretasi}

Investasi dengan ekspektasi lebih tinggi memang tampak menarik. Tetapi
bila variance-nya juga jauh lebih besar, maka keputusan tidak
sesederhana ``pilih yang rata-ratanya terbesar''. Dalam dunia teknik,
sains data, dan bisnis, keputusan hampir selalu menimbang \textbf{nilai
harapan} dan \textbf{ketidakpastian} sekaligus.

\subsubsection{Poin inti}\label{poin-inti}

\begin{itemize}
\tightlist
\item
  Ekspektasi = pusat distribusi
\item
  Varians = ukuran penyebaran
\item
  Dua pilihan dengan mean sama belum tentu sama baik
\item
  Probabilitas membantu kita memilih secara lebih sadar, bukan sekadar
  ``feeling''
\end{itemize}

\subsection{Latihan singkat}\label{latihan-singkat}

\begin{enumerate}
\def\labelenumi{\arabic{enumi}.}
\tightlist
\item
  Buat dua investasi lain yang punya mean hampir sama tetapi variance
  berbeda.
\item
  Simulasikan 10.000 kali dan bandingkan histogramnya.
\item
  Investasi mana yang akan Anda pilih jika Anda tidak tahan melihat
  hasil negatif besar?
\end{enumerate}

\section{1.6 Kasus 2 --- Produk Cacat dan Quality
Control}\label{kasus-2-produk-cacat-dan-quality-control}

\subsection{K --- Konteks}\label{k-konteks-2}

Sebuah pabrik mengklaim bahwa peluang sebuah produk cacat hanya 1\%. Tim
quality control mengambil sampel 100 unit. Mereka ingin tahu peluang
menemukan tepat 2 produk cacat. Keputusan yang akan diambil: apakah
produksi lanjut seperti biasa, perlu inspeksi tambahan, atau bahkan
perlu dihentikan sementara?

\subsection{M --- Model}\label{m-model-2}

Jika: - setiap unit independen, - tiap unit punya peluang cacat
\(p = 0.01\), - jumlah sampel \(n = 100\),

maka banyaknya produk cacat \(X\) dapat dimodelkan sebagai:

\[
X \sim \mathrm{Binomial}(n=100, p=0.01)
\]

Untuk kejadian langka, distribusi Poisson juga sering menjadi
pendekatan:

\[
X \approx \mathrm{Poisson}(\lambda = np = 1)
\]

\subsection{Q --- Questions}\label{q-questions-2}

\begin{enumerate}
\def\labelenumi{\arabic{enumi}.}
\tightlist
\item
  Berapa peluang tepat 2 cacat?
\item
  Seberapa baik Poisson mendekati Binomial pada kasus ini?
\item
  Jika ditemukan terlalu banyak cacat, kapan QC harus curiga?
\end{enumerate}

\subsection{A --- Apply}\label{a-apply-2}

\begin{Shaded}
\begin{Highlighting}[]
\ImportTok{from}\NormalTok{ math }\ImportTok{import}\NormalTok{ comb, exp, factorial}

\NormalTok{n }\OperatorTok{=} \DecValTok{100}
\NormalTok{p }\OperatorTok{=} \FloatTok{0.01}
\NormalTok{k }\OperatorTok{=} \DecValTok{2}
\NormalTok{lam }\OperatorTok{=}\NormalTok{ n }\OperatorTok{*}\NormalTok{ p}

\NormalTok{binom\_p }\OperatorTok{=}\NormalTok{ comb(n, k) }\OperatorTok{*}\NormalTok{ (p}\OperatorTok{**}\NormalTok{k) }\OperatorTok{*}\NormalTok{ ((}\DecValTok{1}\OperatorTok{{-}}\NormalTok{p)}\OperatorTok{**}\NormalTok{(n}\OperatorTok{{-}}\NormalTok{k))}
\NormalTok{pois\_p }\OperatorTok{=}\NormalTok{ exp(}\OperatorTok{{-}}\NormalTok{lam) }\OperatorTok{*}\NormalTok{ (lam}\OperatorTok{**}\NormalTok{k) }\OperatorTok{/}\NormalTok{ factorial(k)}

\BuiltInTok{print}\NormalTok{(}\StringTok{"P(X=2) Binomial ="}\NormalTok{, binom\_p)}
\BuiltInTok{print}\NormalTok{(}\StringTok{"P(X=2) Poisson  ="}\NormalTok{, pois\_p)}
\end{Highlighting}
\end{Shaded}

\begin{verbatim}
P(X=2) Binomial = 0.18486481882486325
P(X=2) Poisson  = 0.18393972058572117
\end{verbatim}

\begin{Shaded}
\begin{Highlighting}[]
\ImportTok{from}\NormalTok{ scipy.stats }\ImportTok{import}\NormalTok{ binom, poisson}

\NormalTok{x }\OperatorTok{=}\NormalTok{ np.arange(}\DecValTok{0}\NormalTok{, }\DecValTok{8}\NormalTok{)}

\NormalTok{plt.figure(figsize}\OperatorTok{=}\NormalTok{(}\DecValTok{9}\NormalTok{,}\DecValTok{4}\NormalTok{))}
\NormalTok{plt.stem(x, binom.pmf(x, n, p), linefmt}\OperatorTok{=}\StringTok{\textquotesingle{}C0{-}\textquotesingle{}}\NormalTok{, markerfmt}\OperatorTok{=}\StringTok{\textquotesingle{}C0o\textquotesingle{}}\NormalTok{, basefmt}\OperatorTok{=}\StringTok{" "}\NormalTok{, label}\OperatorTok{=}\StringTok{"Binomial"}\NormalTok{)}
\NormalTok{plt.stem(x}\OperatorTok{+}\FloatTok{0.1}\NormalTok{, poisson.pmf(x, lam), linefmt}\OperatorTok{=}\StringTok{\textquotesingle{}C1{-}\textquotesingle{}}\NormalTok{, markerfmt}\OperatorTok{=}\StringTok{\textquotesingle{}C1s\textquotesingle{}}\NormalTok{, basefmt}\OperatorTok{=}\StringTok{" "}\NormalTok{, label}\OperatorTok{=}\StringTok{"Poisson approx"}\NormalTok{)}
\NormalTok{plt.xlabel(}\StringTok{"Jumlah produk cacat"}\NormalTok{)}
\NormalTok{plt.ylabel(}\StringTok{"Peluang"}\NormalTok{)}
\NormalTok{plt.title(}\StringTok{"Binomial vs Poisson untuk produk cacat"}\NormalTok{)}
\NormalTok{plt.legend()}
\NormalTok{plt.grid(alpha}\OperatorTok{=}\FloatTok{0.3}\NormalTok{)}
\NormalTok{plt.show()}
\end{Highlighting}
\end{Shaded}

\pandocbounded{\includegraphics[keepaspectratio]{01-pendahuluan-pengambilan-keputusan_files/figure-pdf/cell-7-output-1.pdf}}

\subsubsection{Interpretasi}\label{interpretasi-1}

Kasus ini penting karena memperkenalkan dua hal sekaligus: 1. random
variable diskrit untuk menghitung banyak kejadian, 2. pentingnya memilih
model yang tepat.

Engineer yang baik bukan yang hafal rumus paling banyak, tetapi yang
tahu model mana yang masuk akal untuk suatu konteks.

\subsection{Latihan singkat}\label{latihan-singkat-1}

\begin{enumerate}
\def\labelenumi{\arabic{enumi}.}
\tightlist
\item
  Hitung peluang tepat 0, 1, 2, dan 3 cacat.
\item
  Bandingkan hasil Binomial dan Poisson.
\item
  Jika pabrik menemukan 6 produk cacat dari 100, apakah itu masih terasa
  wajar?
\end{enumerate}

\section{1.7 Kasus 3 --- Garansi Produk dan Risiko
Klaim}\label{kasus-3-garansi-produk-dan-risiko-klaim}

\subsection{K --- Konteks}\label{k-konteks-3}

Sebuah perusahaan lampu ingin menetapkan masa garansi. Umur lampu
bersifat acak. Jika masa garansi terlalu pendek, pelanggan kecewa. Jika
terlalu panjang, biaya klaim bisa membengkak. Keputusan harus
menyeimbangkan daya saing produk dan risiko kerugian.

\subsection{M --- Model}\label{m-model-3}

Misalkan umur lampu \(X\) dinyatakan dalam jam. Untuk contoh awal,
anggap:

\[
X \sim \mathcal{N}(900, 50^2)
\]

Artinya: - rata-rata umur = 900 jam - simpangan baku = 50 jam

Yang ingin dihitung misalnya: \[
P(X < T)
\] yaitu peluang lampu rusak sebelum waktu garansi \(T\).

\subsection{Q --- Questions}\label{q-questions-3}

\begin{enumerate}
\def\labelenumi{\arabic{enumi}.}
\tightlist
\item
  Berapa peluang rusak sebelum 800 jam?
\item
  Bagaimana jika garansi 850 jam? 900 jam?
\item
  Mana kebijakan yang lebih aman bagi perusahaan?
\end{enumerate}

\subsection{A --- Apply}\label{a-apply-3}

\begin{Shaded}
\begin{Highlighting}[]
\ImportTok{from}\NormalTok{ scipy.stats }\ImportTok{import}\NormalTok{ norm}

\NormalTok{mu }\OperatorTok{=} \DecValTok{900}
\NormalTok{sigma }\OperatorTok{=} \DecValTok{50}

\ControlFlowTok{for}\NormalTok{ T }\KeywordTok{in}\NormalTok{ [}\DecValTok{800}\NormalTok{, }\DecValTok{850}\NormalTok{, }\DecValTok{900}\NormalTok{]:}
    \BuiltInTok{print}\NormalTok{(}\SpecialStringTok{f"P(X \textless{} }\SpecialCharTok{\{}\NormalTok{T}\SpecialCharTok{\}}\SpecialStringTok{) ="}\NormalTok{, norm.cdf(T, loc}\OperatorTok{=}\NormalTok{mu, scale}\OperatorTok{=}\NormalTok{sigma))}
\end{Highlighting}
\end{Shaded}

\begin{verbatim}
P(X < 800) = 0.0227501319481792
P(X < 850) = 0.15865525393145707
P(X < 900) = 0.5
\end{verbatim}

\begin{Shaded}
\begin{Highlighting}[]
\NormalTok{x }\OperatorTok{=}\NormalTok{ np.linspace(}\DecValTok{700}\NormalTok{, }\DecValTok{1100}\NormalTok{, }\DecValTok{400}\NormalTok{)}
\NormalTok{pdf }\OperatorTok{=}\NormalTok{ norm.pdf(x, loc}\OperatorTok{=}\NormalTok{mu, scale}\OperatorTok{=}\NormalTok{sigma)}

\NormalTok{plt.figure(figsize}\OperatorTok{=}\NormalTok{(}\DecValTok{9}\NormalTok{,}\DecValTok{4}\NormalTok{))}
\NormalTok{plt.plot(x, pdf, label}\OperatorTok{=}\StringTok{"PDF umur lampu"}\NormalTok{)}
\ControlFlowTok{for}\NormalTok{ T }\KeywordTok{in}\NormalTok{ [}\DecValTok{800}\NormalTok{, }\DecValTok{850}\NormalTok{, }\DecValTok{900}\NormalTok{]:}
\NormalTok{    plt.axvline(T, linestyle}\OperatorTok{=}\StringTok{"{-}{-}"}\NormalTok{, label}\OperatorTok{=}\SpecialStringTok{f"T=}\SpecialCharTok{\{}\NormalTok{T}\SpecialCharTok{\}}\SpecialStringTok{"}\NormalTok{)}
\NormalTok{plt.xlabel(}\StringTok{"Jam hidup"}\NormalTok{)}
\NormalTok{plt.ylabel(}\StringTok{"Density"}\NormalTok{)}
\NormalTok{plt.title(}\StringTok{"Distribusi umur lampu dan pilihan masa garansi"}\NormalTok{)}
\NormalTok{plt.legend()}
\NormalTok{plt.grid(alpha}\OperatorTok{=}\FloatTok{0.3}\NormalTok{)}
\NormalTok{plt.show()}
\end{Highlighting}
\end{Shaded}

\pandocbounded{\includegraphics[keepaspectratio]{01-pendahuluan-pengambilan-keputusan_files/figure-pdf/cell-9-output-1.pdf}}

\subsubsection{Interpretasi}\label{interpretasi-2}

Statistika di sini bukan sekadar menjelaskan data. Ia menjadi alat
desain keputusan. Dengan probabilitas, perusahaan dapat memperkirakan
biaya klaim sebelum kebijakan garansi benar-benar diterapkan.

\subsection{Latihan singkat}\label{latihan-singkat-2}

\begin{enumerate}
\def\labelenumi{\arabic{enumi}.}
\tightlist
\item
  Hitung peluang rusak sebelum 875 jam.
\item
  Jika target klaim maksimum 10\%, berapa garansi maksimum yang bisa
  diberikan?
\item
  Bagaimana jika simpangan baku membesar menjadi 80 jam?
\end{enumerate}

\section{1.8 Kasus 4 --- Fungsi Random Variable dan Keamanan
Sistem}\label{kasus-4-fungsi-random-variable-dan-keamanan-sistem}

\subsection{K --- Konteks}\label{k-konteks-4}

Dalam sistem kelistrikan sederhana, daya dapat dinyatakan sebagai:

\[
P = I^2 R
\]

Jika arus \(I\) memiliki variasi kecil, daya \(P\) juga akan ikut acak.
Tetapi karena hubungan ini kuadratik, perubahan kecil pada arus bisa
membesar pada daya. Untuk sistem teknik, pertanyaan pentingnya adalah:
\emph{seberapa sering daya melewati ambang aman?}

\subsection{M --- Model}\label{m-model-4}

Misalkan: - \(I\) random, - \(R\) konstan, - \(P = g(I) = I^2 R\)

Berarti \(P\) adalah fungsi dari random variable lain. Karakter
distribusi \(P\) bisa sangat berbeda dari distribusi \(I\).

\subsection{Q --- Questions}\label{q-questions-4}

\begin{enumerate}
\def\labelenumi{\arabic{enumi}.}
\tightlist
\item
  Jika \(I\) berdistribusi normal, seperti apa distribusi \(P\)?
\item
  Berapa peluang \(P\) melampaui batas aman tertentu?
\item
  Mengapa transformasi nonlinear penting dalam desain?
\end{enumerate}

\subsection{A --- Apply}\label{a-apply-4}

\begin{Shaded}
\begin{Highlighting}[]
\NormalTok{rng }\OperatorTok{=}\NormalTok{ np.random.default\_rng(}\DecValTok{123}\NormalTok{)}
\NormalTok{I }\OperatorTok{=}\NormalTok{ rng.normal(loc}\OperatorTok{=}\DecValTok{10}\NormalTok{, scale}\OperatorTok{=}\FloatTok{0.8}\NormalTok{, size}\OperatorTok{=}\DecValTok{100000}\NormalTok{)}
\NormalTok{R }\OperatorTok{=} \DecValTok{5}
\NormalTok{P }\OperatorTok{=}\NormalTok{ (I}\OperatorTok{**}\DecValTok{2}\NormalTok{) }\OperatorTok{*}\NormalTok{ R}

\BuiltInTok{print}\NormalTok{(}\StringTok{"Mean(I) ="}\NormalTok{, I.mean(), }\StringTok{"Var(I) ="}\NormalTok{, I.var())}
\BuiltInTok{print}\NormalTok{(}\StringTok{"Mean(P) ="}\NormalTok{, P.mean(), }\StringTok{"Var(P) ="}\NormalTok{, P.var())}
\end{Highlighting}
\end{Shaded}

\begin{verbatim}
Mean(I) = 10.001059538027341 Var(I) = 0.6397573277010377
Mean(P) = 503.3047460543435 Var(P) = 6420.356765767546
\end{verbatim}

\begin{Shaded}
\begin{Highlighting}[]
\NormalTok{plt.figure(figsize}\OperatorTok{=}\NormalTok{(}\DecValTok{9}\NormalTok{,}\DecValTok{4}\NormalTok{))}
\NormalTok{plt.hist(I, bins}\OperatorTok{=}\DecValTok{60}\NormalTok{, alpha}\OperatorTok{=}\FloatTok{0.6}\NormalTok{, density}\OperatorTok{=}\VariableTok{True}\NormalTok{, label}\OperatorTok{=}\StringTok{"Arus I"}\NormalTok{)}
\NormalTok{plt.hist(P, bins}\OperatorTok{=}\DecValTok{60}\NormalTok{, alpha}\OperatorTok{=}\FloatTok{0.6}\NormalTok{, density}\OperatorTok{=}\VariableTok{True}\NormalTok{, label}\OperatorTok{=}\StringTok{"Daya P = I\^{}2 R"}\NormalTok{)}
\NormalTok{plt.xlabel(}\StringTok{"Nilai"}\NormalTok{)}
\NormalTok{plt.ylabel(}\StringTok{"Density"}\NormalTok{)}
\NormalTok{plt.title(}\StringTok{"Distribusi input vs distribusi output setelah transformasi nonlinear"}\NormalTok{)}
\NormalTok{plt.legend()}
\NormalTok{plt.grid(alpha}\OperatorTok{=}\FloatTok{0.3}\NormalTok{)}
\NormalTok{plt.show()}
\end{Highlighting}
\end{Shaded}

\pandocbounded{\includegraphics[keepaspectratio]{01-pendahuluan-pengambilan-keputusan_files/figure-pdf/cell-11-output-1.pdf}}

\subsubsection{Interpretasi}\label{interpretasi-3}

Kasus ini menunjukkan gagasan penting: - random variable tidak selalu
berdiri sendiri, - fungsi dari random variable dapat mengubah bentuk
distribusi, - keputusan teknik sering bergantung pada distribusi output,
bukan sekadar distribusi input.

\subsection{Latihan singkat}\label{latihan-singkat-3}

\begin{enumerate}
\def\labelenumi{\arabic{enumi}.}
\tightlist
\item
  Hitung peluang \(P > 600\).
\item
  Ubah simpangan baku arus menjadi 1.5 dan lihat perubahan distribusi
  daya.
\item
  Mengapa fungsi kuadrat lebih ``berbahaya'' daripada fungsi linear
  dalam konteks batas aman?
\end{enumerate}

\section{1.9 Kasus 5 --- Varians, Profit, dan Risiko
Bangkrut}\label{kasus-5-varians-profit-dan-risiko-bangkrut}

\subsection{K --- Konteks}\label{k-konteks-5}

Sebuah perusahaan memproduksi 1000 unit per hari. Biaya produksi tiap
produk adalah \$100. Produk dijual dengan margin 10\% per unit yang
laku. Penjualan per hari bersifat acak dengan rata-rata 1000 unit.
Profit harian ditambahkan ke saldo modal. Bila saldo modal negatif,
perusahaan bangkrut.

Pertanyaannya: - kalau rata-rata profit positif, apakah perusahaan pasti
aman? - apakah variance besar pada penjualan bisa tetap membunuh
perusahaan?

\subsection{M --- Model}\label{m-model-5}

\begin{itemize}
\tightlist
\item
  biaya per produk = \$100
\item
  produksi per hari = 1000
\item
  harga jual per produk = \$110
\item
  biaya produksi harian = \$100\{,\}000
\end{itemize}

Jika penjualan harian \(X_t\), maka profit harian:

\[
\Pi_t = 110X_t - 100000
\]

Saldo modal diperbarui setiap hari:

\[
M_t = M_{t-1} + \Pi_t
\]

Jika \(M_t < 0\), perusahaan bangkrut.

\subsection{Q --- Questions}\label{q-questions-5}

\begin{enumerate}
\def\labelenumi{\arabic{enumi}.}
\tightlist
\item
  Jika rata-rata penjualan 1000, apakah expected profit harian positif?
\item
  Bagaimana pengaruh variance penjualan terhadap risiko bangkrut?
\item
  Bagaimana peran modal awal 1x, 2x, 3x biaya produksi harian?
\end{enumerate}

\subsection{A --- Apply}\label{a-apply-5}

\begin{Shaded}
\begin{Highlighting}[]
\ImportTok{import}\NormalTok{ numpy }\ImportTok{as}\NormalTok{ np}
\ImportTok{import}\NormalTok{ matplotlib.pyplot }\ImportTok{as}\NormalTok{ plt}

\NormalTok{rng }\OperatorTok{=}\NormalTok{ np.random.default\_rng(}\DecValTok{123}\NormalTok{)}

\NormalTok{cost\_per\_product }\OperatorTok{=} \DecValTok{100}
\NormalTok{production }\OperatorTok{=} \DecValTok{1000}
\NormalTok{selling\_price }\OperatorTok{=} \DecValTok{110}
\NormalTok{daily\_cost }\OperatorTok{=}\NormalTok{ cost\_per\_product }\OperatorTok{*}\NormalTok{ production}
\NormalTok{mean\_sales }\OperatorTok{=} \DecValTok{1000}

\KeywordTok{def}\NormalTok{ simulate\_company(days}\OperatorTok{=}\DecValTok{365}\NormalTok{, initial\_capital}\OperatorTok{=}\DecValTok{100000}\NormalTok{, sales\_std}\OperatorTok{=}\DecValTok{250}\NormalTok{, n\_sims}\OperatorTok{=}\DecValTok{2000}\NormalTok{):}
\NormalTok{    bankrupt }\OperatorTok{=} \DecValTok{0}
\NormalTok{    final\_capitals }\OperatorTok{=}\NormalTok{ []}

    \ControlFlowTok{for}\NormalTok{ \_ }\KeywordTok{in} \BuiltInTok{range}\NormalTok{(n\_sims):}
\NormalTok{        sales }\OperatorTok{=}\NormalTok{ rng.normal(loc}\OperatorTok{=}\NormalTok{mean\_sales, scale}\OperatorTok{=}\NormalTok{sales\_std, size}\OperatorTok{=}\NormalTok{days)}
\NormalTok{        sales }\OperatorTok{=}\NormalTok{ np.maximum(}\DecValTok{0}\NormalTok{, np.}\BuiltInTok{round}\NormalTok{(sales))}
\NormalTok{        profit }\OperatorTok{=}\NormalTok{ selling\_price }\OperatorTok{*}\NormalTok{ sales }\OperatorTok{{-}}\NormalTok{ daily\_cost}
\NormalTok{        capital }\OperatorTok{=}\NormalTok{ initial\_capital }\OperatorTok{+}\NormalTok{ np.cumsum(profit)}

        \ControlFlowTok{if}\NormalTok{ np.}\BuiltInTok{any}\NormalTok{(capital }\OperatorTok{\textless{}} \DecValTok{0}\NormalTok{):}
\NormalTok{            bankrupt }\OperatorTok{+=} \DecValTok{1}
\NormalTok{        final\_capitals.append(capital[}\OperatorTok{{-}}\DecValTok{1}\NormalTok{])}

    \ControlFlowTok{return}\NormalTok{ np.array(final\_capitals), bankrupt }\OperatorTok{/}\NormalTok{ n\_sims}
\end{Highlighting}
\end{Shaded}

\begin{Shaded}
\begin{Highlighting}[]
\NormalTok{capital\_levels }\OperatorTok{=}\NormalTok{ [}\DecValTok{1}\NormalTok{, }\DecValTok{2}\NormalTok{, }\DecValTok{3}\NormalTok{]}
\NormalTok{sales\_stds }\OperatorTok{=}\NormalTok{ [}\DecValTok{50}\NormalTok{, }\DecValTok{150}\NormalTok{, }\DecValTok{300}\NormalTok{, }\DecValTok{500}\NormalTok{]}

\NormalTok{results }\OperatorTok{=}\NormalTok{ \{\}}

\ControlFlowTok{for}\NormalTok{ c }\KeywordTok{in}\NormalTok{ capital\_levels:}
\NormalTok{    init\_cap }\OperatorTok{=}\NormalTok{ c }\OperatorTok{*}\NormalTok{ daily\_cost}
\NormalTok{    results[c] }\OperatorTok{=}\NormalTok{ []}
    \ControlFlowTok{for}\NormalTok{ s }\KeywordTok{in}\NormalTok{ sales\_stds:}
\NormalTok{        finals, p\_bankrupt }\OperatorTok{=}\NormalTok{ simulate\_company(initial\_capital}\OperatorTok{=}\NormalTok{init\_cap, sales\_std}\OperatorTok{=}\NormalTok{s)}
\NormalTok{        results[c].append((s, finals.mean(), p\_bankrupt))}

\ControlFlowTok{for}\NormalTok{ c }\KeywordTok{in}\NormalTok{ capital\_levels:}
    \BuiltInTok{print}\NormalTok{(}\SpecialStringTok{f"}\CharTok{\textbackslash{}n}\SpecialStringTok{Modal awal = }\SpecialCharTok{\{}\NormalTok{c}\SpecialCharTok{\}}\SpecialStringTok{ x biaya harian"}\NormalTok{)}
    \ControlFlowTok{for}\NormalTok{ s, mean\_final, p\_bankrupt }\KeywordTok{in}\NormalTok{ results[c]:}
        \BuiltInTok{print}\NormalTok{(}\SpecialStringTok{f"std sales=}\SpecialCharTok{\{}\NormalTok{s}\SpecialCharTok{:\textgreater{}3\}}\SpecialStringTok{ | mean final capital=}\SpecialCharTok{\{}\NormalTok{mean\_final}\SpecialCharTok{:,.0f\}}\SpecialStringTok{ | P(bankrupt)=}\SpecialCharTok{\{}\NormalTok{p\_bankrupt}\SpecialCharTok{:.3f\}}\SpecialStringTok{"}\NormalTok{)}
\end{Highlighting}
\end{Shaded}

\begin{verbatim}

Modal awal = 1 x biaya harian
std sales= 50 | mean final capital=3,751,041 | P(bankrupt)=0.000
std sales=150 | mean final capital=3,755,095 | P(bankrupt)=0.001
std sales=300 | mean final capital=3,780,991 | P(bankrupt)=0.105
std sales=500 | mean final capital=3,903,849 | P(bankrupt)=0.377

Modal awal = 2 x biaya harian
std sales= 50 | mean final capital=3,847,908 | P(bankrupt)=0.000
std sales=150 | mean final capital=3,857,446 | P(bankrupt)=0.000
std sales=300 | mean final capital=3,849,650 | P(bankrupt)=0.022
std sales=500 | mean final capital=4,004,548 | P(bankrupt)=0.192

Modal awal = 3 x biaya harian
std sales= 50 | mean final capital=3,950,421 | P(bankrupt)=0.000
std sales=150 | mean final capital=3,955,070 | P(bankrupt)=0.000
std sales=300 | mean final capital=3,950,022 | P(bankrupt)=0.002
std sales=500 | mean final capital=4,118,620 | P(bankrupt)=0.087
\end{verbatim}

\begin{Shaded}
\begin{Highlighting}[]
\NormalTok{plt.figure(figsize}\OperatorTok{=}\NormalTok{(}\DecValTok{9}\NormalTok{,}\DecValTok{4}\NormalTok{))}
\ControlFlowTok{for}\NormalTok{ c }\KeywordTok{in}\NormalTok{ capital\_levels:}
\NormalTok{    xs }\OperatorTok{=}\NormalTok{ [r[}\DecValTok{0}\NormalTok{] }\ControlFlowTok{for}\NormalTok{ r }\KeywordTok{in}\NormalTok{ results[c]]}
\NormalTok{    ys }\OperatorTok{=}\NormalTok{ [r[}\DecValTok{2}\NormalTok{] }\ControlFlowTok{for}\NormalTok{ r }\KeywordTok{in}\NormalTok{ results[c]]}
\NormalTok{    plt.plot(xs, ys, marker}\OperatorTok{=}\StringTok{\textquotesingle{}o\textquotesingle{}}\NormalTok{, label}\OperatorTok{=}\SpecialStringTok{f"Modal awal }\SpecialCharTok{\{}\NormalTok{c}\SpecialCharTok{\}}\SpecialStringTok{x biaya harian"}\NormalTok{)}
\NormalTok{plt.xlabel(}\StringTok{"Std penjualan harian"}\NormalTok{)}
\NormalTok{plt.ylabel(}\StringTok{"Peluang bangkrut"}\NormalTok{)}
\NormalTok{plt.title(}\StringTok{"Variance besar dapat mematikan perusahaan walau expected profit positif"}\NormalTok{)}
\NormalTok{plt.legend()}
\NormalTok{plt.grid(alpha}\OperatorTok{=}\FloatTok{0.3}\NormalTok{)}
\NormalTok{plt.show()}
\end{Highlighting}
\end{Shaded}

\pandocbounded{\includegraphics[keepaspectratio]{01-pendahuluan-pengambilan-keputusan_files/figure-pdf/cell-14-output-1.pdf}}

\subsubsection{Interpretasi}\label{interpretasi-4}

Ini adalah salah satu pelajaran paling penting dalam probabilitas untuk
keputusan nyata:

\[
\text{Expected profit positif} \not\Rightarrow \text{pasti selamat}
\]

Kalau variance terlalu besar, perusahaan bisa kehabisan napas sebelum
sempat menikmati rata-rata jangka panjangnya.

\subsubsection{Poin inti}\label{poin-inti-1}

\begin{itemize}
\tightlist
\item
  Mean mengukur ``arah''
\item
  Variance mengukur ``ombak''
\item
  Modal awal adalah ``bantalan''
\item
  Survival bergantung pada kombinasi mean, variance, dan cadangan modal
\end{itemize}

\subsection{Latihan singkat}\label{latihan-singkat-4}

\begin{enumerate}
\def\labelenumi{\arabic{enumi}.}
\tightlist
\item
  Ubah penjualan harian dari Normal menjadi Poisson atau Negative
  Binomial.
\item
  Bandingkan peluang bangkrutnya.
\item
  Mana yang lebih penting ditambah: margin, modal awal, atau pengurangan
  variance?
\end{enumerate}

\section{1.10 Kasus 6 --- Klinik Gigi: Layanan, Throughput, dan
Profit}\label{kasus-6-klinik-gigi-layanan-throughput-dan-profit}

\subsection{K --- Konteks}\label{k-konteks-6}

Sebuah klinik gigi buka 8 jam per hari. Satu dokter rata-rata memerlukan
30 menit untuk melayani seorang pasien. Karena jenis keluhan pasien
berbeda-beda, waktu layanan bersifat acak dan dianggap memoryless.
Rata-rata pasien datang 16 orang per hari.

Pertanyaan awal: - berapa rata-rata pasien yang dapat dilayani satu
dokter per hari? - berapa distribusi profit harian jika dokter dibayar
tetap dan pasien membayar proporsional dengan durasi layanan?

\subsection{M --- Model}\label{m-model-6}

Klinik buka 8 jam = 480 menit per hari.\\
Jika satu layanan rata-rata 30 menit, maka laju layanan dokter:

\[
\mu = 2 \text{ pasien/jam}
\]

Jika rata-rata pasien datang 16 per hari dalam 8 jam, maka laju
kedatangan:

\[
\lambda = 2 \text{ pasien/jam}
\]

Ini adalah model \textbf{M/M/1}: - kedatangan Poisson, - layanan
Eksponensial, - satu pelayan.

Untuk profit: - gaji dokter = \$350/hari - biaya operasional =
\$200/hari - pasien membayar \$1/menit layanan

Jika total waktu layanan aktual dalam sehari adalah \(B\) menit, maka:

\[
\text{Revenue} = B
\] \[
\text{Profit} = B - 350 - 200 = B - 550
\]

Karena \(B \le 480\), maka profit maksimum:

\[
480 - 550 = -70
\]

Artinya klinik pasti rugi pada tarif \$1/menit.

\subsection{Q --- Questions}\label{q-questions-6}

\begin{enumerate}
\def\labelenumi{\arabic{enumi}.}
\tightlist
\item
  Berapa rata-rata pasien yang selesai dilayani per hari?
\item
  Bagaimana histogram pasien terlayani per hari?
\item
  Bagaimana distribusi profit per hari?
\item
  Jika ada dua dokter, apa yang berubah pada throughput dan profit?
\end{enumerate}

\subsection{A --- Apply}\label{a-apply-6}

\begin{Shaded}
\begin{Highlighting}[]
\ImportTok{import}\NormalTok{ numpy }\ImportTok{as}\NormalTok{ np}
\ImportTok{import}\NormalTok{ matplotlib.pyplot }\ImportTok{as}\NormalTok{ plt}

\NormalTok{rng }\OperatorTok{=}\NormalTok{ np.random.default\_rng(}\DecValTok{42}\NormalTok{)}

\KeywordTok{def}\NormalTok{ simulate\_mm1\_day(arrival\_rate\_per\_hour}\OperatorTok{=}\FloatTok{2.0}\NormalTok{, service\_rate\_per\_hour}\OperatorTok{=}\FloatTok{2.0}\NormalTok{, hours}\OperatorTok{=}\FloatTok{8.0}\NormalTok{):}
    \CommentTok{\#\# Generate arrivals}
\NormalTok{    arrivals }\OperatorTok{=}\NormalTok{ []}
\NormalTok{    t }\OperatorTok{=} \FloatTok{0.0}
    \ControlFlowTok{while} \VariableTok{True}\NormalTok{:}
\NormalTok{        t }\OperatorTok{+=}\NormalTok{ rng.exponential(scale}\OperatorTok{=}\DecValTok{1}\OperatorTok{/}\NormalTok{arrival\_rate\_per\_hour)}
        \ControlFlowTok{if}\NormalTok{ t }\OperatorTok{\textgreater{}}\NormalTok{ hours:}
            \ControlFlowTok{break}
\NormalTok{        arrivals.append(t)}

\NormalTok{    service\_end\_prev }\OperatorTok{=} \FloatTok{0.0}
\NormalTok{    served\_count }\OperatorTok{=} \DecValTok{0}
\NormalTok{    busy\_hours }\OperatorTok{=} \FloatTok{0.0}

    \ControlFlowTok{for}\NormalTok{ a }\KeywordTok{in}\NormalTok{ arrivals:}
\NormalTok{        service\_time }\OperatorTok{=}\NormalTok{ rng.exponential(scale}\OperatorTok{=}\DecValTok{1}\OperatorTok{/}\NormalTok{service\_rate\_per\_hour)}
\NormalTok{        start }\OperatorTok{=} \BuiltInTok{max}\NormalTok{(a, service\_end\_prev)}
\NormalTok{        end }\OperatorTok{=}\NormalTok{ start }\OperatorTok{+}\NormalTok{ service\_time}

\NormalTok{        overlap\_start }\OperatorTok{=} \BuiltInTok{max}\NormalTok{(start, }\FloatTok{0.0}\NormalTok{)}
\NormalTok{        overlap\_end }\OperatorTok{=} \BuiltInTok{min}\NormalTok{(end, hours)}
        \ControlFlowTok{if}\NormalTok{ overlap\_end }\OperatorTok{\textgreater{}}\NormalTok{ overlap\_start:}
\NormalTok{            busy\_hours }\OperatorTok{+=}\NormalTok{ overlap\_end }\OperatorTok{{-}}\NormalTok{ overlap\_start}

        \ControlFlowTok{if}\NormalTok{ end }\OperatorTok{\textless{}=}\NormalTok{ hours:}
\NormalTok{            served\_count }\OperatorTok{+=} \DecValTok{1}
\NormalTok{            service\_end\_prev }\OperatorTok{=}\NormalTok{ end}
        \ControlFlowTok{else}\NormalTok{:}
            \ControlFlowTok{break}

\NormalTok{    busy\_minutes }\OperatorTok{=} \DecValTok{60} \OperatorTok{*}\NormalTok{ busy\_hours}
    \ControlFlowTok{return}\NormalTok{ served\_count, busy\_minutes}
\end{Highlighting}
\end{Shaded}

\begin{Shaded}
\begin{Highlighting}[]
\NormalTok{served }\OperatorTok{=}\NormalTok{ []}
\NormalTok{profits }\OperatorTok{=}\NormalTok{ []}

\ControlFlowTok{for}\NormalTok{ \_ }\KeywordTok{in} \BuiltInTok{range}\NormalTok{(}\DecValTok{10000}\NormalTok{):}
\NormalTok{    s, b }\OperatorTok{=}\NormalTok{ simulate\_mm1\_day()}
\NormalTok{    served.append(s)}
\NormalTok{    profits.append(b }\OperatorTok{{-}} \DecValTok{350} \OperatorTok{{-}} \DecValTok{200}\NormalTok{)}

\NormalTok{served }\OperatorTok{=}\NormalTok{ np.array(served)}
\NormalTok{profits }\OperatorTok{=}\NormalTok{ np.array(profits)}

\BuiltInTok{print}\NormalTok{(}\StringTok{"Rata{-}rata pasien terlayani per hari ="}\NormalTok{, served.mean())}
\BuiltInTok{print}\NormalTok{(}\StringTok{"Rata{-}rata profit per hari ="}\NormalTok{, profits.mean())}
\BuiltInTok{print}\NormalTok{(}\StringTok{"Profit maksimum simulasi ="}\NormalTok{, profits.}\BuiltInTok{max}\NormalTok{())}
\end{Highlighting}
\end{Shaded}

\begin{verbatim}
Rata-rata pasien terlayani per hari = 11.995
Rata-rata profit per hari = -190.51550749020004
Profit maksimum simulasi = -70.07101058383995
\end{verbatim}

\begin{Shaded}
\begin{Highlighting}[]
\NormalTok{plt.figure(figsize}\OperatorTok{=}\NormalTok{(}\DecValTok{9}\NormalTok{,}\DecValTok{4}\NormalTok{))}
\NormalTok{bins }\OperatorTok{=}\NormalTok{ np.arange(served.}\BuiltInTok{min}\NormalTok{()}\OperatorTok{{-}}\FloatTok{0.5}\NormalTok{, served.}\BuiltInTok{max}\NormalTok{()}\OperatorTok{+}\FloatTok{1.5}\NormalTok{, }\DecValTok{1}\NormalTok{)}
\NormalTok{plt.hist(served, bins}\OperatorTok{=}\NormalTok{bins, density}\OperatorTok{=}\VariableTok{True}\NormalTok{, alpha}\OperatorTok{=}\FloatTok{0.75}\NormalTok{, edgecolor}\OperatorTok{=}\StringTok{\textquotesingle{}black\textquotesingle{}}\NormalTok{)}
\NormalTok{plt.xlabel(}\StringTok{"Jumlah pasien selesai dilayani"}\NormalTok{)}
\NormalTok{plt.ylabel(}\StringTok{"Probabilitas relatif"}\NormalTok{)}
\NormalTok{plt.title(}\StringTok{"Histogram jumlah pasien terlayani per hari"}\NormalTok{)}
\NormalTok{plt.grid(alpha}\OperatorTok{=}\FloatTok{0.3}\NormalTok{)}
\NormalTok{plt.show()}
\end{Highlighting}
\end{Shaded}

\pandocbounded{\includegraphics[keepaspectratio]{01-pendahuluan-pengambilan-keputusan_files/figure-pdf/cell-17-output-1.pdf}}

\begin{Shaded}
\begin{Highlighting}[]
\NormalTok{plt.figure(figsize}\OperatorTok{=}\NormalTok{(}\DecValTok{9}\NormalTok{,}\DecValTok{4}\NormalTok{))}
\NormalTok{plt.hist(profits, bins}\OperatorTok{=}\DecValTok{60}\NormalTok{, density}\OperatorTok{=}\VariableTok{True}\NormalTok{, alpha}\OperatorTok{=}\FloatTok{0.75}\NormalTok{, edgecolor}\OperatorTok{=}\StringTok{\textquotesingle{}black\textquotesingle{}}\NormalTok{)}
\NormalTok{plt.xlabel(}\StringTok{"Profit per hari ($)"}\NormalTok{)}
\NormalTok{plt.ylabel(}\StringTok{"Density"}\NormalTok{)}
\NormalTok{plt.title(}\StringTok{"Distribusi profit per hari klinik gigi"}\NormalTok{)}
\NormalTok{plt.grid(alpha}\OperatorTok{=}\FloatTok{0.3}\NormalTok{)}
\NormalTok{plt.show()}
\end{Highlighting}
\end{Shaded}

\pandocbounded{\includegraphics[keepaspectratio]{01-pendahuluan-pengambilan-keputusan_files/figure-pdf/cell-18-output-1.pdf}}

\subsubsection{Interpretasi}\label{interpretasi-5}

Kasus ini mengajarkan beberapa hal sekaligus: - kapasitas rata-rata
dokter, - pengaruh ketidakpastian layanan dan kedatangan, - perbedaan
antara ``rata-rata teoretis'' dan ``jumlah selesai terlayani sebelum
tutup'', - kaitan antara model antrian dan model profit.

Untuk dua dokter, throughput meningkat, tetapi biaya tetap juga naik.
Keputusan menambah dokter harus menimbang: - peningkatan kapasitas, -
perubahan waiting time, - perubahan busy time, - dan perubahan profit.

\subsection{Latihan singkat}\label{latihan-singkat-5}

\begin{enumerate}
\def\labelenumi{\arabic{enumi}.}
\tightlist
\item
  Simulasikan versi dua dokter.
\item
  Bandingkan histogram pasien terlayani.
\item
  Cari tarif minimum per menit agar secara teori klinik bisa impas.
\end{enumerate}

\section{1.11 Menyimpulkan Bab Ini}\label{menyimpulkan-bab-ini}

Bab ini memperkenalkan satu ide besar: \textbf{probabilitas dan
statistika adalah alat untuk membuat keputusan ketika masa depan belum
pasti}.

Kita telah melihat bahwa: - investasi menuntut kita membandingkan
ekspektasi dan risiko, - quality control membutuhkan model jumlah cacat,
- garansi produk membutuhkan peluang kegagalan sebelum waktu tertentu, -
fungsi random variable penting untuk keselamatan sistem, - variance
dapat membuat perusahaan bangkrut walaupun mean profit positif, - model
antrian dapat menjelaskan throughput dan profit layanan.

Di balik semua contoh itu, ada benang merah yang sama: kita memerlukan
cara untuk memetakan ketidakpastian menjadi objek yang bisa dihitung.
Itulah peran random variable.

\section{1.12 Ringkasan Poin Inti}\label{ringkasan-poin-inti}

\begin{enumerate}
\def\labelenumi{\arabic{enumi}.}
\tightlist
\item
  \textbf{Pengambilan keputusan} adalah konteks alami bagi probabilitas
  dan statistika.
\item
  \textbf{Random variable} memetakan dunia acak ke bilangan.
\item
  \textbf{Ekspektasi} membantu kita melihat kecenderungan rata-rata.
\item
  \textbf{Varians} membantu kita mengukur seberapa liar hasil yang
  mungkin terjadi.
\item
  \textbf{Python} memberi quick wins: melihat, mensimulasikan, dan
  memahami sebelum formalisasi penuh.
\item
  \textbf{Model yang tepat} penting: Binomial, Poisson, Normal,
  Exponential, dan lainnya tidak dapat dipertukarkan sembarangan.
\item
  Dalam banyak masalah nyata, keputusan yang baik tidak cukup melihat
  mean; ia juga harus melihat \textbf{risiko, bentuk distribusi, dan
  konsekuensi}.
\end{enumerate}

\section{1.13 Latihan Bab 1}\label{latihan-bab-1}

\subsection{A. Konseptual}\label{a.-konseptual}

\begin{enumerate}
\def\labelenumi{\arabic{enumi}.}
\tightlist
\item
  Mengapa mean saja tidak cukup untuk membandingkan dua pilihan?
\item
  Apa perbedaan antara hasil acak dan random variable?
\item
  Mengapa keputusan garansi produk adalah masalah probabilistik?
\end{enumerate}

\subsection{B. Komputasional}\label{b.-komputasional}

\begin{enumerate}
\def\labelenumi{\arabic{enumi}.}
\tightlist
\item
  Simulasikan dua investasi dengan mean sama tetapi variance berbeda.
\item
  Hitung peluang tepat 3 produk cacat dari 100 jika \(p=0.02\).
\item
  Hitung peluang sebuah lampu rusak sebelum 850 jam jika umur lampu
  normal dengan mean 900 dan simpangan baku 60.
\end{enumerate}

\subsection{C. Aplikatif}\label{c.-aplikatif}

\begin{enumerate}
\def\labelenumi{\arabic{enumi}.}
\tightlist
\item
  Dalam kasus perusahaan, apakah Anda lebih memilih menaikkan margin,
  menambah modal awal, atau menurunkan variance penjualan? Jelaskan.
\item
  Dalam kasus klinik gigi, apakah Anda akan menambah dokter kedua jika
  tarif masih \$1/menit? Mengapa?
\item
  Berikan satu contoh lain dari kehidupan sehari-hari yang menurut Anda
  cocok dimodelkan dengan random variable.
\end{enumerate}

\section{1.14 Penutup Kecil}\label{penutup-kecil}

Kalau setelah membaca bab ini Anda merasa, ``Ternyata probabilitas bukan
sekadar rumus, tetapi cara berpikir,'' maka bab ini sudah mencapai
tujuannya.

Di bab berikutnya, kita akan memperdalam apa sebenarnya random variable
itu, bagaimana range-nya didefinisikan, dan bagaimana PMF, CDF, PDF,
ekspektasi, serta varians menjadi fondasi seluruh analisis berikutnya.

\bookmarksetup{startatroot}

\chapter{Bab 2. Random Variable Umum}\label{bab-2.-random-variable-umum}

\section{Tujuan Bab}\label{tujuan-bab-1}

Setelah mempelajari bab ini, mahasiswa diharapkan mampu:

\begin{enumerate}
\def\labelenumi{\arabic{enumi}.}
\tightlist
\item
  memahami \textbf{random variable} sebagai pemetaan dari hasil acak ke
  bilangan,
\item
  membedakan \textbf{event} dengan \textbf{nilai random variable},
\item
  memahami \textbf{range} sebagai himpunan nilai yang mungkin,
\item
  membedakan random variable \textbf{diskrit} dan \textbf{kontinu},
\item
  memahami dan menggunakan \textbf{PMF}, \textbf{CDF}, dan \textbf{PDF},
\item
  menghitung dan menafsirkan \textbf{ekspektasi}, \textbf{varians}, dan
  \textbf{simpangan baku},
\item
  menghubungkan ukuran-ukuran tersebut dengan \textbf{pengambilan
  keputusan}.
\end{enumerate}

\section{Pembuka}\label{pembuka-1}

Pada bab sebelumnya, kita sudah melihat bahwa probabilitas menjadi
penting ketika kita harus mengambil keputusan di bawah ketidakpastian.
Kita juga sudah melihat bahwa salah satu cara paling kuat untuk
``menjinakkan'' ketidakpastian adalah dengan membuat \textbf{model}.

Di bab ini, kita masuk ke fondasi model tersebut: \textbf{random
variable}.

Random variable adalah salah satu ide yang tampak sederhana, tetapi
pengaruhnya sangat besar. Begitu kita berhasil memetakan hasil-hasil
acak menjadi angka, kita memperoleh akses ke dunia matematika: kita bisa
menggambar, menghitung, mensimulasikan, membandingkan, dan akhirnya
memutuskan.

\section{2.1 Quick Win: Dari Dunia Acak ke Garis
Bilangan}\label{quick-win-dari-dunia-acak-ke-garis-bilangan}

Mari mulai dari contoh sederhana. Misalkan kita melempar dua koin. Ruang
sampelnya adalah:

\[
S = \{HH, HT, TH, TT\}
\]

Sekarang definisikan random variable:

\[
X = \text{jumlah sisi Head yang muncul}
\]

Maka: - \(X(HH)=2\) - \(X(HT)=1\) - \(X(TH)=1\) - \(X(TT)=0\)

Artinya, hasil-hasil acak yang tadinya berupa simbol \(HH, HT, TH, TT\)
dipetakan menjadi angka \(0,1,2\).

Mari lihat dengan Python.

\begin{Shaded}
\begin{Highlighting}[]
\ImportTok{import}\NormalTok{ itertools}
\ImportTok{import}\NormalTok{ numpy }\ImportTok{as}\NormalTok{ np}
\ImportTok{import}\NormalTok{ matplotlib.pyplot }\ImportTok{as}\NormalTok{ plt}

\NormalTok{outcomes }\OperatorTok{=} \BuiltInTok{list}\NormalTok{(itertools.product([}\StringTok{\textquotesingle{}H\textquotesingle{}}\NormalTok{, }\StringTok{\textquotesingle{}T\textquotesingle{}}\NormalTok{], repeat}\OperatorTok{=}\DecValTok{2}\NormalTok{))}

\KeywordTok{def}\NormalTok{ X(outcome):}
    \ControlFlowTok{return}\NormalTok{ outcome.count(}\StringTok{\textquotesingle{}H\textquotesingle{}}\NormalTok{)}

\NormalTok{mapped }\OperatorTok{=}\NormalTok{ [(o, X(o)) }\ControlFlowTok{for}\NormalTok{ o }\KeywordTok{in}\NormalTok{ outcomes]}
\NormalTok{mapped}
\end{Highlighting}
\end{Shaded}

\begin{verbatim}
[(('H', 'H'), 2), (('H', 'T'), 1), (('T', 'H'), 1), (('T', 'T'), 0)]
\end{verbatim}

\begin{Shaded}
\begin{Highlighting}[]
\NormalTok{values }\OperatorTok{=}\NormalTok{ [X(o) }\ControlFlowTok{for}\NormalTok{ o }\KeywordTok{in}\NormalTok{ outcomes]}
\NormalTok{unique, counts }\OperatorTok{=}\NormalTok{ np.unique(values, return\_counts}\OperatorTok{=}\VariableTok{True}\NormalTok{)}

\NormalTok{plt.figure(figsize}\OperatorTok{=}\NormalTok{(}\DecValTok{7}\NormalTok{,}\DecValTok{4}\NormalTok{))}
\NormalTok{plt.bar(unique, counts }\OperatorTok{/}\NormalTok{ counts.}\BuiltInTok{sum}\NormalTok{())}
\NormalTok{plt.xlabel(}\StringTok{"Nilai X = jumlah Head"}\NormalTok{)}
\NormalTok{plt.ylabel(}\StringTok{"Probabilitas"}\NormalTok{)}
\NormalTok{plt.title(}\StringTok{"Random variable sebagai pemetaan ke bilangan"}\NormalTok{)}
\NormalTok{plt.xticks(unique)}
\NormalTok{plt.grid(alpha}\OperatorTok{=}\FloatTok{0.3}\NormalTok{)}
\NormalTok{plt.show()}
\end{Highlighting}
\end{Shaded}

\pandocbounded{\includegraphics[keepaspectratio]{02-random-variable-umum_files/figure-pdf/cell-3-output-1.pdf}}

Dari sini kita melihat \emph{quick win} yang penting: - event hidup di
ruang sampel, - random variable memetakan event ke angka, - distribusi
random variable memberi tahu kita bagaimana peluang tersebar pada
angka-angka itu.

\section{2.2 K --- Konteks: Mengapa Kita Perlu Random
Variable?}\label{k-konteks-mengapa-kita-perlu-random-variable}

Bayangkan beberapa pertanyaan berikut: - Berapa jumlah kendaraan yang
lewat gerbang tol dalam satu jam? - Berapa waktu tunggu pasien sampai
dilayani? - Berapa profit harian sebuah klinik? - Berapa banyak produk
cacat dalam 100 unit? - Berapa lama sebuah lampu bertahan sebelum rusak?

Semua pertanyaan ini punya satu pola: hasilnya \textbf{tidak pasti},
tetapi hasil itu \textbf{berupa bilangan}.

Itulah alasan random variable menjadi penting. Ia menjadi jembatan
antara: - kejadian acak di dunia nyata, dan - analisis matematis
berbasis angka.

\section{2.3 M --- Model: Definisi Random
Variable}\label{m-model-definisi-random-variable}

Secara formal, random variable \(X\) adalah sebuah fungsi yang memetakan
setiap hasil \(\omega\) dalam ruang sampel \(S\) ke sebuah bilangan
real:

\[
X: S \to \mathbb{R}
\]

dengan

\[
\omega \mapsto X(\omega)
\]

Di sini: - \(S\) = ruang sampel, - \(\omega\) = satu hasil acak, -
\(X(\omega)\) = angka yang dihasilkan oleh pemetaan itu.

\subsection{Catatan penting}\label{catatan-penting}

Random variable bukan ``angka yang acak'' dalam arti longgar. Secara
matematis, random variable adalah \textbf{fungsi}. Keacakan datang dari
kenyataan bahwa hasil \(\omega\) belum diketahui sebelum eksperimen
terjadi.

\section{2.4 Event vs Nilai Random
Variable}\label{event-vs-nilai-random-variable}

Mahasiswa sering bingung membedakan dua hal ini:

\begin{itemize}
\tightlist
\item
  \textbf{event} = himpunan hasil dalam ruang sampel,
\item
  \textbf{nilai random variable} = angka hasil pemetaan.
\end{itemize}

Contoh: lempar dua dadu, lalu definisikan

\[
X = \text{jumlah mata dadu}
\]

Maka event: \[
E = \{(1,6), (2,5), (3,4), (4,3), (5,2), (6,1)\}
\]

adalah event ``jumlah mata dadu = 7''.

Tetapi nilai random variable yang terkait adalah: \[
X = 7
\]

Jadi: - event adalah kumpulan outcome, - random variable adalah angka
yang dihasilkan dari outcome itu.

Mari visualisasikan cepat.

\begin{Shaded}
\begin{Highlighting}[]
\NormalTok{outcomes }\OperatorTok{=} \BuiltInTok{list}\NormalTok{(itertools.product(}\BuiltInTok{range}\NormalTok{(}\DecValTok{1}\NormalTok{, }\DecValTok{7}\NormalTok{), repeat}\OperatorTok{=}\DecValTok{2}\NormalTok{))}
\NormalTok{sums }\OperatorTok{=}\NormalTok{ [}\BuiltInTok{sum}\NormalTok{(o) }\ControlFlowTok{for}\NormalTok{ o }\KeywordTok{in}\NormalTok{ outcomes]}

\NormalTok{unique, counts }\OperatorTok{=}\NormalTok{ np.unique(sums, return\_counts}\OperatorTok{=}\VariableTok{True}\NormalTok{)}
\NormalTok{probs }\OperatorTok{=}\NormalTok{ counts }\OperatorTok{/}\NormalTok{ counts.}\BuiltInTok{sum}\NormalTok{()}

\NormalTok{plt.figure(figsize}\OperatorTok{=}\NormalTok{(}\DecValTok{8}\NormalTok{,}\DecValTok{4}\NormalTok{))}
\NormalTok{plt.bar(unique, probs)}
\NormalTok{plt.xlabel(}\StringTok{"Nilai X = jumlah dua dadu"}\NormalTok{)}
\NormalTok{plt.ylabel(}\StringTok{"Probabilitas"}\NormalTok{)}
\NormalTok{plt.title(}\StringTok{"Distribusi random variable jumlah dua dadu"}\NormalTok{)}
\NormalTok{plt.xticks(unique)}
\NormalTok{plt.grid(alpha}\OperatorTok{=}\FloatTok{0.3}\NormalTok{)}
\NormalTok{plt.show()}
\end{Highlighting}
\end{Shaded}

\pandocbounded{\includegraphics[keepaspectratio]{02-random-variable-umum_files/figure-pdf/cell-4-output-1.pdf}}

\section{2.5 Range: Nilai-Nilai yang
Mungkin}\label{range-nilai-nilai-yang-mungkin}

Range dari random variable adalah himpunan semua nilai yang mungkin
diambil oleh random variable tersebut.

Untuk contoh dua dadu dengan: \[
X = \text{jumlah mata dadu}
\]

range-nya adalah: \[
\{2,3,4,\dots,12\}
\]

Range ini penting karena memberi tahu: - apa saja kemungkinan nilai, -
apakah nilai-nilai itu diskrit atau kontinu, - bagaimana bentuk
distribusinya.

\subsection{Diskrit vs Kontinu}\label{diskrit-vs-kontinu}

\subsubsection{Random Variable Diskrit}\label{random-variable-diskrit}

Jika range terdiri dari titik-titik terpisah, misalnya: \[
\{0,1,2,3,\dots\}
\] maka random variable itu diskrit.

Contoh: - jumlah pelanggan datang, - jumlah cacat, - jumlah pesan masuk,
- banyaknya Head.

\subsubsection{Random Variable Kontinu}\label{random-variable-kontinu}

Jika range berupa interval pada garis bilangan, misalnya: \[
[0,\infty)
\] atau \[
(-\infty,\infty)
\] maka random variable itu kontinu.

Contoh: - waktu tunggu, - lama hidup komponen, - tinggi badan, - suhu.

\section{2.6 PMF untuk Random Variable
Diskrit}\label{pmf-untuk-random-variable-diskrit}

Jika \(X\) diskrit, maka distribusinya dapat dinyatakan dengan
\textbf{Probability Mass Function (PMF)}:

\[
p_X(x) = P(X=x)
\]

PMF memberi peluang bahwa random variable mengambil nilai tepat \(x\).

\subsection{Sifat PMF}\label{sifat-pmf}

\begin{enumerate}
\def\labelenumi{\arabic{enumi}.}
\item
  \[
  p_X(x) \ge 0
  \]
\item
  \[
  \sum_x p_X(x) = 1
  \]
\end{enumerate}

\subsubsection{Contoh: jumlah Head pada dua lemparan
koin}\label{contoh-jumlah-head-pada-dua-lemparan-koin}

\[
P(X=0)=\frac{1}{4}, \quad P(X=1)=\frac{1}{2}, \quad P(X=2)=\frac{1}{4}
\]

\begin{Shaded}
\begin{Highlighting}[]
\NormalTok{x }\OperatorTok{=}\NormalTok{ np.array([}\DecValTok{0}\NormalTok{, }\DecValTok{1}\NormalTok{, }\DecValTok{2}\NormalTok{])}
\NormalTok{pmf }\OperatorTok{=}\NormalTok{ np.array([}\DecValTok{1}\OperatorTok{/}\DecValTok{4}\NormalTok{, }\DecValTok{1}\OperatorTok{/}\DecValTok{2}\NormalTok{, }\DecValTok{1}\OperatorTok{/}\DecValTok{4}\NormalTok{])}

\NormalTok{plt.figure(figsize}\OperatorTok{=}\NormalTok{(}\DecValTok{7}\NormalTok{,}\DecValTok{4}\NormalTok{))}
\NormalTok{plt.stem(x, pmf, basefmt}\OperatorTok{=}\StringTok{" "}\NormalTok{)}
\NormalTok{plt.xlabel(}\StringTok{"x"}\NormalTok{)}
\NormalTok{plt.ylabel(}\StringTok{"P(X=x)"}\NormalTok{)}
\NormalTok{plt.title(}\StringTok{"PMF jumlah Head pada dua lemparan koin"}\NormalTok{)}
\NormalTok{plt.grid(alpha}\OperatorTok{=}\FloatTok{0.3}\NormalTok{)}
\NormalTok{plt.show()}
\end{Highlighting}
\end{Shaded}

\pandocbounded{\includegraphics[keepaspectratio]{02-random-variable-umum_files/figure-pdf/cell-5-output-1.pdf}}

\subsubsection{Interpretasi keputusan}\label{interpretasi-keputusan}

PMF membantu menjawab pertanyaan seperti: - seberapa mungkin tepat 2
cacat? - seberapa mungkin tepat 5 pelanggan datang? - seberapa mungkin
tepat 1 dari 3 server gagal?

\section{2.7 CDF: Bahasa Probabilitas
Kumulatif}\label{cdf-bahasa-probabilitas-kumulatif}

Baik untuk random variable diskrit maupun kontinu, salah satu fungsi
terpenting adalah \textbf{Cumulative Distribution Function (CDF)}:

\[
F_X(x) = P(X \le x)
\]

CDF memberi peluang bahwa nilai random variable \textbf{tidak melebihi}
\(x\).

Ini sangat penting karena banyak keputusan berbentuk: - peluang rusak
sebelum waktu tertentu, - peluang demand tidak melebihi stok, - peluang
profit lebih kecil dari nol, - peluang nilai ujian di bawah batas lulus.

\subsubsection{Contoh CDF diskrit}\label{contoh-cdf-diskrit}

Untuk \(X\) = jumlah Head pada dua lemparan koin: -
\(F(0)=P(X\le 0)=1/4\) - \(F(1)=P(X\le 1)=3/4\) - \(F(2)=1\)

\begin{Shaded}
\begin{Highlighting}[]
\NormalTok{cdf }\OperatorTok{=}\NormalTok{ np.cumsum(pmf)}

\NormalTok{plt.figure(figsize}\OperatorTok{=}\NormalTok{(}\DecValTok{7}\NormalTok{,}\DecValTok{4}\NormalTok{))}
\NormalTok{plt.step(x, cdf, where}\OperatorTok{=}\StringTok{\textquotesingle{}post\textquotesingle{}}\NormalTok{)}
\NormalTok{plt.xlabel(}\StringTok{"x"}\NormalTok{)}
\NormalTok{plt.ylabel(}\StringTok{"F(x)=P(X≤x)"}\NormalTok{)}
\NormalTok{plt.title(}\StringTok{"CDF random variable diskrit"}\NormalTok{)}
\NormalTok{plt.ylim(}\OperatorTok{{-}}\FloatTok{0.05}\NormalTok{, }\FloatTok{1.05}\NormalTok{)}
\NormalTok{plt.grid(alpha}\OperatorTok{=}\FloatTok{0.3}\NormalTok{)}
\NormalTok{plt.show()}
\end{Highlighting}
\end{Shaded}

\pandocbounded{\includegraphics[keepaspectratio]{02-random-variable-umum_files/figure-pdf/cell-6-output-1.pdf}}

\section{2.8 PDF untuk Random Variable
Kontinu}\label{pdf-untuk-random-variable-kontinu}

Untuk random variable kontinu, kita biasanya tidak berbicara tentang
peluang tepat pada satu titik, karena:

\[
P(X=x)=0
\]

Sebagai gantinya, distribusi dijelaskan dengan \textbf{Probability
Density Function (PDF)}:

\[
f_X(x)
\]

Peluang pada interval diperoleh melalui integral:

\[
P(a \le X \le b) = \int_a^b f_X(x)\,dx
\]

\subsection{Sifat PDF}\label{sifat-pdf}

\begin{enumerate}
\def\labelenumi{\arabic{enumi}.}
\item
  \[
  f_X(x)\ge 0
  \]
\item
  \[
  \int_{-\infty}^{\infty} f_X(x)\,dx = 1
  \]
\end{enumerate}

\subsubsection{Contoh cepat: Uniform kontinu pada
{[}0,1{]}}\label{contoh-cepat-uniform-kontinu-pada-01}

\[
f(x)=
\begin{cases}
1, & 0 \le x \le 1\\
0, & \text{lainnya}
\end{cases}
\]

Maka: \[
P(0.2 \le X \le 0.7)=0.5
\]

\begin{Shaded}
\begin{Highlighting}[]
\NormalTok{x }\OperatorTok{=}\NormalTok{ np.linspace(}\OperatorTok{{-}}\FloatTok{0.2}\NormalTok{, }\FloatTok{1.2}\NormalTok{, }\DecValTok{400}\NormalTok{)}
\NormalTok{pdf }\OperatorTok{=}\NormalTok{ np.where((x }\OperatorTok{\textgreater{}=} \DecValTok{0}\NormalTok{) }\OperatorTok{\&}\NormalTok{ (x }\OperatorTok{\textless{}=} \DecValTok{1}\NormalTok{), }\FloatTok{1.0}\NormalTok{, }\FloatTok{0.0}\NormalTok{)}

\NormalTok{plt.figure(figsize}\OperatorTok{=}\NormalTok{(}\DecValTok{8}\NormalTok{,}\DecValTok{4}\NormalTok{))}
\NormalTok{plt.plot(x, pdf)}
\NormalTok{plt.fill\_between(x, }\DecValTok{0}\NormalTok{, pdf, where}\OperatorTok{=}\NormalTok{((x }\OperatorTok{\textgreater{}=} \FloatTok{0.2}\NormalTok{) }\OperatorTok{\&}\NormalTok{ (x }\OperatorTok{\textless{}=} \FloatTok{0.7}\NormalTok{)), alpha}\OperatorTok{=}\FloatTok{0.3}\NormalTok{)}
\NormalTok{plt.xlabel(}\StringTok{"x"}\NormalTok{)}
\NormalTok{plt.ylabel(}\StringTok{"f(x)"}\NormalTok{)}
\NormalTok{plt.title(}\StringTok{"PDF Uniform(0,1)"}\NormalTok{)}
\NormalTok{plt.grid(alpha}\OperatorTok{=}\FloatTok{0.3}\NormalTok{)}
\NormalTok{plt.show()}
\end{Highlighting}
\end{Shaded}

\pandocbounded{\includegraphics[keepaspectratio]{02-random-variable-umum_files/figure-pdf/cell-7-output-1.pdf}}

\subsubsection{Interpretasi penting}\label{interpretasi-penting}

Tinggi PDF bukan peluang langsung.\\
Peluang adalah \textbf{luas area} di bawah kurva PDF.

\section{2.9 Menghubungkan PMF, CDF, dan
PDF}\label{menghubungkan-pmf-cdf-dan-pdf}

Hubungan penting:

\subsection{Untuk diskrit}\label{untuk-diskrit}

\[
F_X(x)=\sum_{t \le x} p_X(t)
\]

\subsection{Untuk kontinu}\label{untuk-kontinu}

\[
F_X(x)=\int_{-\infty}^{x} f_X(t)\,dt
\]

Jika PDF cukup halus, maka: \[
f_X(x)=\frac{d}{dx}F_X(x)
\]

Jadi: - PMF/PDF memberi struktur lokal distribusi, - CDF memberi
akumulasi peluang sampai titik tertentu.

\section{2.10 Ekspektasi: Tebakan Terbaik Sebelum Fakta
Terjadi}\label{ekspektasi-tebakan-terbaik-sebelum-fakta-terjadi}

Ekspektasi adalah ukuran pusat distribusi. Ia sering dipandang sebagai
\textbf{rata-rata jangka panjang}.

\subsection{Untuk random variable
diskrit}\label{untuk-random-variable-diskrit}

\[
E[X] = \sum_x x\,p_X(x)
\]

\subsection{Untuk random variable
kontinu}\label{untuk-random-variable-kontinu}

\[
E[X] = \int_{-\infty}^{\infty} x f_X(x)\,dx
\]

\subsubsection{Contoh: jumlah Head pada dua lemparan
koin}\label{contoh-jumlah-head-pada-dua-lemparan-koin-1}

\[
E[X] = 0 \cdot \frac14 + 1 \cdot \frac12 + 2 \cdot \frac14 = 1
\]

Mari cek dengan simulasi.

\begin{Shaded}
\begin{Highlighting}[]
\NormalTok{rng }\OperatorTok{=}\NormalTok{ np.random.default\_rng(}\DecValTok{123}\NormalTok{)}

\NormalTok{n }\OperatorTok{=} \DecValTok{100000}
\NormalTok{coin1 }\OperatorTok{=}\NormalTok{ rng.integers(}\DecValTok{0}\NormalTok{, }\DecValTok{2}\NormalTok{, size}\OperatorTok{=}\NormalTok{n)}
\NormalTok{coin2 }\OperatorTok{=}\NormalTok{ rng.integers(}\DecValTok{0}\NormalTok{, }\DecValTok{2}\NormalTok{, size}\OperatorTok{=}\NormalTok{n)}
\NormalTok{Xsim }\OperatorTok{=}\NormalTok{ coin1 }\OperatorTok{+}\NormalTok{ coin2}

\BuiltInTok{print}\NormalTok{(}\StringTok{"Rata{-}rata simulasi ="}\NormalTok{, Xsim.mean())}
\end{Highlighting}
\end{Shaded}

\begin{verbatim}
Rata-rata simulasi = 1.00235
\end{verbatim}

\subsubsection{Interpretasi keputusan}\label{interpretasi-keputusan-1}

Ekspektasi menjawab pertanyaan seperti: - rata-rata keuntungan harian, -
rata-rata jumlah pelanggan, - rata-rata cacat per batch, - rata-rata
waktu tunggu.

Tetapi ekspektasi \textbf{tidak mengatakan seluruh cerita}. Dua
distribusi dengan mean sama bisa memiliki risiko yang sangat berbeda.

\section{2.11 Varians dan Simpangan Baku: Seberapa Liar
Penyimpangannya?}\label{varians-dan-simpangan-baku-seberapa-liar-penyimpangannya}

Varians mengukur seberapa jauh random variable menyebar dari nilai
ekspektasinya.

\subsection{Definisi}\label{definisi}

\[
\mathrm{Var}(X)=E[(X-E[X])^2]
\]

Simpangan baku adalah akar kuadrat varians: \[
\sigma_X = \sqrt{\mathrm{Var}(X)}
\]

\subsection{Bentuk yang sering lebih
praktis}\label{bentuk-yang-sering-lebih-praktis}

\[
\mathrm{Var}(X)=E[X^2] - (E[X])^2
\]

\subsubsection{Contoh: dua investasi dengan mean sama, variance
beda}\label{contoh-dua-investasi-dengan-mean-sama-variance-beda}

\begin{Shaded}
\begin{Highlighting}[]
\NormalTok{rng }\OperatorTok{=}\NormalTok{ np.random.default\_rng(}\DecValTok{2026}\NormalTok{)}

\NormalTok{A }\OperatorTok{=}\NormalTok{ rng.choice([}\DecValTok{8}\NormalTok{, }\DecValTok{10}\NormalTok{, }\DecValTok{12}\NormalTok{], size}\OperatorTok{=}\DecValTok{100000}\NormalTok{, p}\OperatorTok{=}\NormalTok{[}\FloatTok{0.25}\NormalTok{, }\FloatTok{0.50}\NormalTok{, }\FloatTok{0.25}\NormalTok{])}
\NormalTok{B }\OperatorTok{=}\NormalTok{ rng.choice([}\DecValTok{0}\NormalTok{, }\DecValTok{10}\NormalTok{, }\DecValTok{20}\NormalTok{], size}\OperatorTok{=}\DecValTok{100000}\NormalTok{, p}\OperatorTok{=}\NormalTok{[}\FloatTok{0.25}\NormalTok{, }\FloatTok{0.50}\NormalTok{, }\FloatTok{0.25}\NormalTok{])}

\BuiltInTok{print}\NormalTok{(}\StringTok{"Mean A ="}\NormalTok{, A.mean(), }\StringTok{"Var A ="}\NormalTok{, A.var(), }\StringTok{"Std A ="}\NormalTok{, A.std())}
\BuiltInTok{print}\NormalTok{(}\StringTok{"Mean B ="}\NormalTok{, B.mean(), }\StringTok{"Var B ="}\NormalTok{, B.var(), }\StringTok{"Std B ="}\NormalTok{, B.std())}
\end{Highlighting}
\end{Shaded}

\begin{verbatim}
Mean A = 9.98924 Var A = 2.0013242224 Std A = 1.414681668220805
Mean B = 9.9584 Var B = 49.78426944 Std B = 7.055796867824356
\end{verbatim}

\begin{Shaded}
\begin{Highlighting}[]
\NormalTok{plt.figure(figsize}\OperatorTok{=}\NormalTok{(}\DecValTok{8}\NormalTok{,}\DecValTok{4}\NormalTok{))}
\NormalTok{plt.hist(A, bins}\OperatorTok{=}\NormalTok{np.arange(}\OperatorTok{{-}}\FloatTok{0.5}\NormalTok{, }\FloatTok{21.5}\NormalTok{, }\DecValTok{1}\NormalTok{), alpha}\OperatorTok{=}\FloatTok{0.6}\NormalTok{, density}\OperatorTok{=}\VariableTok{True}\NormalTok{, label}\OperatorTok{=}\StringTok{"A"}\NormalTok{)}
\NormalTok{plt.hist(B, bins}\OperatorTok{=}\NormalTok{np.arange(}\OperatorTok{{-}}\FloatTok{0.5}\NormalTok{, }\FloatTok{21.5}\NormalTok{, }\DecValTok{1}\NormalTok{), alpha}\OperatorTok{=}\FloatTok{0.6}\NormalTok{, density}\OperatorTok{=}\VariableTok{True}\NormalTok{, label}\OperatorTok{=}\StringTok{"B"}\NormalTok{)}
\NormalTok{plt.xlabel(}\StringTok{"Nilai"}\NormalTok{)}
\NormalTok{plt.ylabel(}\StringTok{"Density / Probability"}\NormalTok{)}
\NormalTok{plt.title(}\StringTok{"Mean sama, variance berbeda"}\NormalTok{)}
\NormalTok{plt.legend()}
\NormalTok{plt.grid(alpha}\OperatorTok{=}\FloatTok{0.3}\NormalTok{)}
\NormalTok{plt.show()}
\end{Highlighting}
\end{Shaded}

\pandocbounded{\includegraphics[keepaspectratio]{02-random-variable-umum_files/figure-pdf/cell-10-output-1.pdf}}

\subsubsection{Interpretasi keputusan}\label{interpretasi-keputusan-2}

Jika kita hanya melihat mean, dua pilihan ini tampak setara. Tetapi jika
kita juga melihat variance, kita sadar bahwa salah satu pilihan jauh
lebih berisiko.

\section{2.12 Sifat-Sifat Penting Ekspektasi dan
Varians}\label{sifat-sifat-penting-ekspektasi-dan-varians}

\subsection{Linearitas ekspektasi}\label{linearitas-ekspektasi}

Untuk konstanta \(a,b\): \[
E[aX+b] = aE[X]+b
\]

Ini sangat penting karena memungkinkan kita memindahkan model ke unit
baru atau mengubah skala.

\subsection{Varians terhadap transformasi
linear}\label{varians-terhadap-transformasi-linear}

\[
\mathrm{Var}(aX+b)=a^2 \mathrm{Var}(X)
\]

Artinya: - menambah konstanta tidak mengubah variance, - mengalikan
skala dengan \(a\) memperbesar variance sebesar \(a^2\).

Mari cek.

\begin{Shaded}
\begin{Highlighting}[]
\NormalTok{rng }\OperatorTok{=}\NormalTok{ np.random.default\_rng(}\DecValTok{0}\NormalTok{)}
\NormalTok{X }\OperatorTok{=}\NormalTok{ rng.normal(loc}\OperatorTok{=}\DecValTok{5}\NormalTok{, scale}\OperatorTok{=}\DecValTok{2}\NormalTok{, size}\OperatorTok{=}\DecValTok{100000}\NormalTok{)}

\NormalTok{Y }\OperatorTok{=} \DecValTok{3}\OperatorTok{*}\NormalTok{X }\OperatorTok{+} \DecValTok{10}

\BuiltInTok{print}\NormalTok{(}\StringTok{"E[X] \textasciitilde{}"}\NormalTok{, X.mean())}
\BuiltInTok{print}\NormalTok{(}\StringTok{"E[Y] \textasciitilde{}"}\NormalTok{, Y.mean(), }\StringTok{"   teori ="}\NormalTok{, }\DecValTok{3}\OperatorTok{*}\NormalTok{X.mean() }\OperatorTok{+} \DecValTok{10}\NormalTok{)}

\BuiltInTok{print}\NormalTok{(}\StringTok{"Var[X] \textasciitilde{}"}\NormalTok{, X.var())}
\BuiltInTok{print}\NormalTok{(}\StringTok{"Var[Y] \textasciitilde{}"}\NormalTok{, Y.var(), }\StringTok{"   teori ="}\NormalTok{, }\DecValTok{9}\OperatorTok{*}\NormalTok{X.var())}
\end{Highlighting}
\end{Shaded}

\begin{verbatim}
E[X] ~ 4.998183498453758
E[Y] ~ 24.994550495361278    teori = 24.994550495361274
Var[X] ~ 4.001028098478324
Var[Y] ~ 36.00925288630492    teori = 36.009252886304914
\end{verbatim}

\section{2.13 Diskrit dan Kontinu dalam
Praktik}\label{diskrit-dan-kontinu-dalam-praktik}

Secara praktik, banyak data diukur secara diskrit: - berat badan
dibulatkan ke 0.1 kg, - waktu dicatat dalam detik, - panjang dicatat
dalam milimeter.

Tetapi secara model, kadang jauh lebih sederhana menganggap besaran itu
kontinu. Ini bukan kesalahan, tetapi bentuk idealisasi yang membantu
analisis.

\subsection{Pedoman intuitif}\label{pedoman-intuitif}

\begin{itemize}
\tightlist
\item
  Jika nilai alami berupa hitungan → cenderung diskrit
\item
  Jika nilai alami berupa pengukuran halus → sering dimodelkan kontinu
\end{itemize}

Contoh: - jumlah mobil → diskrit - waktu hidup lampu → kontinu - jumlah
pasien → diskrit - suhu → kontinu

\section{2.14 Dari Histogram ke
Distribusi}\label{dari-histogram-ke-distribusi}

Dalam kenyataan, kita sering tidak mulai dari rumus distribusi, tetapi
dari data. Dari data, kita bisa membuat histogram untuk melihat bentuk
penyebaran.

Mari contohkan data simulasi.

\begin{Shaded}
\begin{Highlighting}[]
\NormalTok{rng }\OperatorTok{=}\NormalTok{ np.random.default\_rng(}\DecValTok{99}\NormalTok{)}
\NormalTok{samples }\OperatorTok{=}\NormalTok{ rng.normal(loc}\OperatorTok{=}\DecValTok{100}\NormalTok{, scale}\OperatorTok{=}\DecValTok{15}\NormalTok{, size}\OperatorTok{=}\DecValTok{5000}\NormalTok{)}

\NormalTok{plt.figure(figsize}\OperatorTok{=}\NormalTok{(}\DecValTok{8}\NormalTok{,}\DecValTok{4}\NormalTok{))}
\NormalTok{plt.hist(samples, bins}\OperatorTok{=}\DecValTok{40}\NormalTok{, density}\OperatorTok{=}\VariableTok{True}\NormalTok{, alpha}\OperatorTok{=}\FloatTok{0.75}\NormalTok{, edgecolor}\OperatorTok{=}\StringTok{\textquotesingle{}black\textquotesingle{}}\NormalTok{)}
\NormalTok{plt.xlabel(}\StringTok{"Nilai"}\NormalTok{)}
\NormalTok{plt.ylabel(}\StringTok{"Density"}\NormalTok{)}
\NormalTok{plt.title(}\StringTok{"Histogram sebagai pendekatan bentuk distribusi"}\NormalTok{)}
\NormalTok{plt.grid(alpha}\OperatorTok{=}\FloatTok{0.3}\NormalTok{)}
\NormalTok{plt.show()}
\end{Highlighting}
\end{Shaded}

\pandocbounded{\includegraphics[keepaspectratio]{02-random-variable-umum_files/figure-pdf/cell-12-output-1.pdf}}

Histogram tidak otomatis menjadi PMF atau PDF yang eksak, tetapi
histogram sangat berguna untuk: - menduga bentuk distribusi, - melihat
pusat, - melihat penyebaran, - mendeteksi skewness atau outlier, -
memilih model awal yang masuk akal.

\section{2.15 KMQA Mini-Case 1 --- Timbangan Berat
Badan}\label{kmqa-mini-case-1-timbangan-berat-badan}

\subsection{K --- Konteks}\label{k-konteks-7}

Sebuah timbangan digital menampilkan berat badan dalam satu angka
desimal.

\subsection{M --- Model}\label{m-model-7}

Jika \(X\) adalah berat badan orang yang naik ke timbangan, maka \(X\)
dapat dimodelkan sebagai random variable kontinu, walaupun tampilan
akhirnya dibulatkan.

\subsection{Q --- Questions}\label{q-questions-7}

\begin{enumerate}
\def\labelenumi{\arabic{enumi}.}
\tightlist
\item
  Apa range masuk akalnya?
\item
  Apakah lebih masuk akal dipandang diskrit atau kontinu?
\item
  Ukuran apa yang penting: mean, variance, atau keduanya?
\end{enumerate}

\subsection{A --- Apply}\label{a-apply-7}

Jawabannya: - range teknis mungkin 0--999.9, tetapi range realistis jauh
lebih sempit, - secara teori lebih nyaman dipandang kontinu, - mean
memberi pusat berat populasi, - variance memberi seberapa beragam berat
badan orang yang diukur.

\section{2.16 KMQA Mini-Case 2 --- Jumlah Kendaraan di Gerbang
Tol}\label{kmqa-mini-case-2-jumlah-kendaraan-di-gerbang-tol}

\subsection{K --- Konteks}\label{k-konteks-8}

Kita ingin memodelkan jumlah kendaraan yang lewat dalam satu jam.

\subsection{M --- Model}\label{m-model-8}

Jika \(X\) = jumlah kendaraan, maka \(X\) adalah random variable
diskrit.

\subsection{Q --- Questions}\label{q-questions-8}

\begin{enumerate}
\def\labelenumi{\arabic{enumi}.}
\tightlist
\item
  Mengapa diskrit?
\item
  Apa range yang masuk akal?
\item
  Apa arti \(P(X=120)\)? Apa arti \(P(X\le 120)\)?
\end{enumerate}

\subsection{A --- Apply}\label{a-apply-8}

\begin{itemize}
\tightlist
\item
  diskrit karena berupa hitungan,
\item
  \(P(X=120)\) adalah peluang tepat 120 kendaraan,
\item
  \(P(X\le 120)\) adalah peluang jumlah kendaraan tidak melebihi 120.
\end{itemize}

\section{2.17 KMQA Mini-Case 3 --- Waktu Tunggu
Pasien}\label{kmqa-mini-case-3-waktu-tunggu-pasien}

\subsection{K --- Konteks}\label{k-konteks-9}

Sebuah klinik ingin memodelkan waktu tunggu pasien hingga dilayani.

\subsection{M --- Model}\label{m-model-9}

Jika \(X\) = waktu tunggu (menit), maka \(X\) biasanya dimodelkan
kontinu.

\subsection{Q --- Questions}\label{q-questions-9}

\begin{enumerate}
\def\labelenumi{\arabic{enumi}.}
\tightlist
\item
  Mengapa kontinu?
\item
  Mengapa \(P(X=10)=0\) tetapi \(P(9.5 \le X \le 10.5)\) bisa positif?
\item
  Mengapa CDF penting dalam konteks ini?
\end{enumerate}

\subsection{A --- Apply}\label{a-apply-9}

\begin{itemize}
\tightlist
\item
  kontinu karena waktu dapat berubah halus,
\item
  pada model kontinu, peluang titik tunggal nol,
\item
  yang bermakna adalah peluang interval,
\item
  CDF membantu menjawab ``peluang menunggu tidak lebih dari x menit''.
\end{itemize}

\section{2.18 Python Toolbox Dasar untuk Bab
Ini}\label{python-toolbox-dasar-untuk-bab-ini}

Berikut beberapa alat Python yang sangat sering dipakai untuk random
variable.

\begin{Shaded}
\begin{Highlighting}[]
\ImportTok{import}\NormalTok{ numpy }\ImportTok{as}\NormalTok{ np}
\ImportTok{from}\NormalTok{ scipy }\ImportTok{import}\NormalTok{ stats}
\end{Highlighting}
\end{Shaded}

\subsection{Simulasi data}\label{simulasi-data}

\begin{Shaded}
\begin{Highlighting}[]
\NormalTok{sample\_binom }\OperatorTok{=}\NormalTok{ np.random.binomial(n}\OperatorTok{=}\DecValTok{10}\NormalTok{, p}\OperatorTok{=}\FloatTok{0.4}\NormalTok{, size}\OperatorTok{=}\DecValTok{10}\NormalTok{)}
\NormalTok{sample\_norm }\OperatorTok{=}\NormalTok{ np.random.normal(loc}\OperatorTok{=}\DecValTok{0}\NormalTok{, scale}\OperatorTok{=}\DecValTok{1}\NormalTok{, size}\OperatorTok{=}\DecValTok{10}\NormalTok{)}

\NormalTok{sample\_binom, sample\_norm[:}\DecValTok{5}\NormalTok{]}
\end{Highlighting}
\end{Shaded}

\begin{verbatim}
(array([3, 6, 4, 5, 4, 7, 3, 3, 5, 6], dtype=int32),
 array([-0.44419803,  0.11322245, -0.06029038,  0.32626076, -0.06141967]))
\end{verbatim}

\subsection{PMF / PDF / CDF dari SciPy}\label{pmf-pdf-cdf-dari-scipy}

\begin{Shaded}
\begin{Highlighting}[]
\NormalTok{x\_discrete }\OperatorTok{=}\NormalTok{ np.arange(}\DecValTok{0}\NormalTok{, }\DecValTok{11}\NormalTok{)}
\NormalTok{pmf\_binom }\OperatorTok{=}\NormalTok{ stats.binom.pmf(x\_discrete, n}\OperatorTok{=}\DecValTok{10}\NormalTok{, p}\OperatorTok{=}\FloatTok{0.4}\NormalTok{)}

\NormalTok{x\_cont }\OperatorTok{=}\NormalTok{ np.linspace(}\OperatorTok{{-}}\DecValTok{3}\NormalTok{, }\DecValTok{3}\NormalTok{, }\DecValTok{200}\NormalTok{)}
\NormalTok{pdf\_norm }\OperatorTok{=}\NormalTok{ stats.norm.pdf(x\_cont, loc}\OperatorTok{=}\DecValTok{0}\NormalTok{, scale}\OperatorTok{=}\DecValTok{1}\NormalTok{)}
\NormalTok{cdf\_norm }\OperatorTok{=}\NormalTok{ stats.norm.cdf(x\_cont, loc}\OperatorTok{=}\DecValTok{0}\NormalTok{, scale}\OperatorTok{=}\DecValTok{1}\NormalTok{)}

\NormalTok{pmf\_binom[:}\DecValTok{5}\NormalTok{], pdf\_norm[:}\DecValTok{5}\NormalTok{], cdf\_norm[:}\DecValTok{5}\NormalTok{]}
\end{Highlighting}
\end{Shaded}

\begin{verbatim}
(array([0.00604662, 0.04031078, 0.12093235, 0.21499085, 0.25082266]),
 array([0.00443185, 0.0048492 , 0.00530104, 0.00578971, 0.00631769]),
 array([0.0013499 , 0.00148973, 0.00164266, 0.00180976, 0.00199218]))
\end{verbatim}

\subsection{Plot cepat}\label{plot-cepat}

\begin{Shaded}
\begin{Highlighting}[]
\NormalTok{plt.figure(figsize}\OperatorTok{=}\NormalTok{(}\DecValTok{8}\NormalTok{,}\DecValTok{4}\NormalTok{))}
\NormalTok{plt.stem(x\_discrete, pmf\_binom, basefmt}\OperatorTok{=}\StringTok{" "}\NormalTok{)}
\NormalTok{plt.title(}\StringTok{"Contoh PMF Binomial dari scipy.stats"}\NormalTok{)}
\NormalTok{plt.xlabel(}\StringTok{"x"}\NormalTok{)}
\NormalTok{plt.ylabel(}\StringTok{"P(X=x)"}\NormalTok{)}
\NormalTok{plt.grid(alpha}\OperatorTok{=}\FloatTok{0.3}\NormalTok{)}
\NormalTok{plt.show()}
\end{Highlighting}
\end{Shaded}

\pandocbounded{\includegraphics[keepaspectratio]{02-random-variable-umum_files/figure-pdf/cell-16-output-1.pdf}}

\section{2.19 Kesalahan-Kesalahan Umum yang Harus
Dihindari}\label{kesalahan-kesalahan-umum-yang-harus-dihindari}

\begin{enumerate}
\def\labelenumi{\arabic{enumi}.}
\item
  \textbf{Menganggap random variable sama dengan event}\\
  Tidak. Event adalah himpunan outcome. Random variable adalah fungsi ke
  bilangan.
\item
  \textbf{Menganggap PDF adalah peluang langsung}\\
  Tidak. Untuk kontinu, peluang adalah luas area di bawah PDF.
\item
  \textbf{Menganggap mean cukup untuk semua keputusan}\\
  Tidak. Variance dan bentuk distribusi juga penting.
\item
  \textbf{Menganggap diskrit dan kontinu bisa dipertukarkan
  sembarangan}\\
  Tidak. Pemilihan model harus sesuai konteks.
\item
  \textbf{Menganggap \(P(X=x)\) untuk kontinu bisa positif}\\
  Tidak. Dalam model kontinu, peluang tepat satu titik adalah nol.
\end{enumerate}

\section{2.20 Menyimpulkan Bab Ini}\label{menyimpulkan-bab-ini-1}

Bab ini memberi fondasi yang akan dipakai terus-menerus pada bab-bab
berikutnya.

Kita telah melihat bahwa: - random variable adalah fungsi dari ruang
sampel ke bilangan real, - range membantu kita mengenali jenis peubah
acak, - PMF dipakai untuk diskrit, - PDF dipakai untuk kontinu, - CDF
berlaku untuk keduanya, - ekspektasi memberi ukuran pusat, - varians dan
simpangan baku memberi ukuran penyebaran.

Dengan fondasi ini, kita siap masuk ke distribusi-distribusi khusus yang
sangat sering dipakai dalam pemodelan nyata.

\section{2.21 Ringkasan Poin Inti}\label{ringkasan-poin-inti-1}

\begin{enumerate}
\def\labelenumi{\arabic{enumi}.}
\tightlist
\item
  \textbf{Random variable} memetakan hasil acak menjadi angka.
\item
  \textbf{Event} berbeda dari \textbf{nilai random variable}.
\item
  \textbf{Range} adalah himpunan nilai yang mungkin diambil.
\item
  \textbf{PMF} untuk diskrit, \textbf{PDF} untuk kontinu, dan
  \textbf{CDF} untuk keduanya.
\item
  \textbf{Ekspektasi} mengukur pusat atau rata-rata jangka panjang.
\item
  \textbf{Varians} dan \textbf{simpangan baku} mengukur penyebaran.
\item
  Dalam keputusan nyata, kita hampir selalu perlu melihat lebih dari
  sekadar mean.
\end{enumerate}

\section{2.22 Latihan Bab 2}\label{latihan-bab-2}

\subsection{A. Konseptual}\label{a.-konseptual-1}

\begin{enumerate}
\def\labelenumi{\arabic{enumi}.}
\tightlist
\item
  Jelaskan dengan kata-kata sendiri apa itu random variable.
\item
  Apa perbedaan antara event ``jumlah mata dua dadu = 7'' dan nilai
  random variable \(X=7\)?
\item
  Mengapa PDF bukan peluang langsung?
\end{enumerate}

\subsection{B. Hitungan}\label{b.-hitungan}

\begin{enumerate}
\def\labelenumi{\arabic{enumi}.}
\tightlist
\item
  Untuk dua lemparan koin, hitung PMF dan CDF dari jumlah Head.
\item
  Untuk dua dadu fair, hitung:

  \begin{itemize}
  \tightlist
  \item
    \(P(X=8)\)
  \item
    \(P(X\le 8)\) dengan \(X\) = jumlah dua dadu.
  \end{itemize}
\item
  Jika \(X\) uniform pada {[}0,1{]}, hitung:

  \begin{itemize}
  \tightlist
  \item
    \(P(X<0.3)\)
  \item
    \(P(0.3 \le X \le 0.8)\)
  \end{itemize}
\end{enumerate}

\subsection{C. Python}\label{c.-python}

\begin{enumerate}
\def\labelenumi{\arabic{enumi}.}
\tightlist
\item
  Simulasikan 50.000 lemparan dua koin dan bandingkan histogram hasil
  dengan PMF teoritis.
\item
  Simulasikan 50.000 lemparan dua dadu dan gambar CDF empiris jumlah
  mata dadu.
\item
  Simulasikan random variable kontinu sederhana dan tunjukkan bahwa
  peluang titik tunggal mendekati nol.
\end{enumerate}

\subsection{D. Aplikatif}\label{d.-aplikatif}

\begin{enumerate}
\def\labelenumi{\arabic{enumi}.}
\tightlist
\item
  Sebutkan satu random variable diskrit dan satu kontinu dari kehidupan
  sehari-hari.
\item
  Dalam konteks layanan klinik, ukuran mana yang lebih penting: mean
  waktu tunggu atau peluang menunggu lebih dari 30 menit? Jelaskan.
\item
  Dalam konteks quality control, mengapa \(P(X\le k)\) sering lebih
  berguna daripada hanya \(P(X=k)\)?
\end{enumerate}

\section{2.23 Penutup Kecil}\label{penutup-kecil-1}

Bab ini mungkin terlihat ``dasar'', tetapi fondasi yang kuat justru
dibuat dari konsep-konsep dasar yang dipahami dengan jernih. Begitu Anda
benar-benar memahami random variable, PMF, CDF, PDF, ekspektasi, dan
varians, banyak topik berikutnya akan terasa jauh lebih masuk akal.

Di bab berikutnya, kita akan mulai memasuki keluarga-keluarga
\textbf{distribusi diskrit khusus} yang sangat sering dipakai dalam
pemodelan: Bernoulli, Binomial, Geometric, Poisson, dan lainnya.

\bookmarksetup{startatroot}

\chapter{Bab 3. Distribusi Random Variable
Diskrit}\label{bab-3.-distribusi-random-variable-diskrit}

\section{Tujuan Bab}\label{tujuan-bab-2}

Setelah mempelajari bab ini, mahasiswa diharapkan mampu:

\begin{enumerate}
\def\labelenumi{\arabic{enumi}.}
\tightlist
\item
  memahami apa yang dimaksud dengan \textbf{distribusi random variable
  diskrit},
\item
  membangun distribusi diskrit dari \textbf{data histogram} atau dari
  \textbf{model probabilistik},
\item
  mengenali dan menggunakan distribusi:

  \begin{itemize}
  \tightlist
  \item
    discrete uniform,
  \item
    Bernoulli,
  \item
    Binomial,
  \item
    Geometric,
  \item
    Poisson,
  \end{itemize}
\item
  memahami makna parameter setiap distribusi,
\item
  menggunakan Python untuk:

  \begin{itemize}
  \tightlist
  \item
    simulasi,
  \item
    perhitungan PMF/CDF,
  \item
    visualisasi,
  \item
    dan pengambilan keputusan,
  \end{itemize}
\item
  memilih distribusi diskrit yang tepat untuk konteks masalah tertentu.
\end{enumerate}

\section{Pembuka}\label{pembuka-2}

Di bab sebelumnya, kita sudah belajar bahwa random variable adalah cara
memetakan dunia acak ke angka. Sekarang kita melangkah lebih jauh:
setelah punya random variable, kita perlu tahu \textbf{bagaimana peluang
tersebar} pada nilai-nilai yang mungkin.

Itulah yang disebut \textbf{distribusi}.

Untuk random variable diskrit, distribusi memberi jawaban atas
pertanyaan seperti: - seberapa mungkin tepat 2 produk cacat? - seberapa
mungkin tepat 10 pelanggan datang? - seberapa mungkin perlu 4 kali
percobaan sampai sukses pertama? - seberapa mungkin jumlah panggilan
masuk dalam satu jam sama dengan 7?

Bab ini akan memperkenalkan keluarga distribusi diskrit yang paling
penting dan paling sering dipakai.

\section{3.1 Quick Win: Dari Histogram ke
Model}\label{quick-win-dari-histogram-ke-model}

Bayangkan kita mengamati jumlah pelanggan yang datang ke sebuah kios
kopi per jam selama beberapa hari, lalu mendapatkan data seperti
berikut.

\begin{Shaded}
\begin{Highlighting}[]
\ImportTok{import}\NormalTok{ numpy }\ImportTok{as}\NormalTok{ np}
\ImportTok{import}\NormalTok{ matplotlib.pyplot }\ImportTok{as}\NormalTok{ plt}
\ImportTok{from}\NormalTok{ scipy }\ImportTok{import}\NormalTok{ stats}

\NormalTok{rng }\OperatorTok{=}\NormalTok{ np.random.default\_rng(}\DecValTok{42}\NormalTok{)}
\NormalTok{data }\OperatorTok{=}\NormalTok{ rng.poisson(lam}\OperatorTok{=}\DecValTok{5}\NormalTok{, size}\OperatorTok{=}\DecValTok{500}\NormalTok{)}

\NormalTok{values, counts }\OperatorTok{=}\NormalTok{ np.unique(data, return\_counts}\OperatorTok{=}\VariableTok{True}\NormalTok{)}
\NormalTok{probs }\OperatorTok{=}\NormalTok{ counts }\OperatorTok{/}\NormalTok{ counts.}\BuiltInTok{sum}\NormalTok{()}

\NormalTok{plt.figure(figsize}\OperatorTok{=}\NormalTok{(}\DecValTok{8}\NormalTok{,}\DecValTok{4}\NormalTok{))}
\NormalTok{plt.bar(values, probs)}
\NormalTok{plt.xlabel(}\StringTok{"Jumlah pelanggan per jam"}\NormalTok{)}
\NormalTok{plt.ylabel(}\StringTok{"Frekuensi relatif"}\NormalTok{)}
\NormalTok{plt.title(}\StringTok{"Histogram empiris jumlah pelanggan"}\NormalTok{)}
\NormalTok{plt.grid(alpha}\OperatorTok{=}\FloatTok{0.3}\NormalTok{)}
\NormalTok{plt.show()}
\end{Highlighting}
\end{Shaded}

\pandocbounded{\includegraphics[keepaspectratio]{03-distribusi-diskrit_files/figure-pdf/cell-2-output-1.pdf}}

Histogram ini belum tentu langsung memberi kita distribusi teoretis,
tetapi ia memberi tiga hal penting: 1. range nilai yang mungkin, 2.
lokasi pusat distribusi, 3. bentuk kasar penyebaran.

Dari sini kita bisa mulai bertanya: model mana yang paling masuk akal?

\section{3.2 Distribusi Diskrit Kustom dari
Histogram}\label{distribusi-diskrit-kustom-dari-histogram}

\subsection{K --- Konteks}\label{k-konteks-10}

Kadang-kadang kita tidak mulai dari rumus distribusi yang terkenal. Kita
mulai dari data empiris. Misalnya: - distribusi nilai kuis, - banyak
pelanggan datang per jam, - banyak tiket gangguan per hari, - jumlah
pembelian aplikasi per sesi.

\subsection{M --- Model}\label{m-model-10}

Jika kita sudah punya histogram atau frekuensi relatif, kita bisa
memperlakukannya sebagai \textbf{PMF empiris}:

\[
p_X(x_i) \approx \frac{\text{frekuensi pada } x_i}{\text{total data}}
\]

Ini adalah distribusi diskrit kustom.

\subsection{Q --- Questions}\label{q-questions-10}

\begin{enumerate}
\def\labelenumi{\arabic{enumi}.}
\tightlist
\item
  Bagaimana membangun PMF dari data?
\item
  Bagaimana menghitung CDF empiris?
\item
  Bagaimana menghitung mean dan variance dari distribusi diskrit yang
  dibangun dari histogram?
\end{enumerate}

\subsection{A --- Apply}\label{a-apply-10}

\begin{Shaded}
\begin{Highlighting}[]
\NormalTok{x }\OperatorTok{=}\NormalTok{ values}
\NormalTok{pmf\_emp }\OperatorTok{=}\NormalTok{ probs}

\NormalTok{mean\_emp }\OperatorTok{=}\NormalTok{ np.}\BuiltInTok{sum}\NormalTok{(x }\OperatorTok{*}\NormalTok{ pmf\_emp)}
\NormalTok{var\_emp }\OperatorTok{=}\NormalTok{ np.}\BuiltInTok{sum}\NormalTok{((x }\OperatorTok{{-}}\NormalTok{ mean\_emp)}\OperatorTok{**}\DecValTok{2} \OperatorTok{*}\NormalTok{ pmf\_emp)}
\NormalTok{cdf\_emp }\OperatorTok{=}\NormalTok{ np.cumsum(pmf\_emp)}

\BuiltInTok{print}\NormalTok{(}\StringTok{"Mean empiris ="}\NormalTok{, mean\_emp)}
\BuiltInTok{print}\NormalTok{(}\StringTok{"Variance empiris ="}\NormalTok{, var\_emp)}
\end{Highlighting}
\end{Shaded}

\begin{verbatim}
Mean empiris = 5.036
Variance empiris = 4.718704
\end{verbatim}

\begin{Shaded}
\begin{Highlighting}[]
\NormalTok{plt.figure(figsize}\OperatorTok{=}\NormalTok{(}\DecValTok{8}\NormalTok{,}\DecValTok{4}\NormalTok{))}
\NormalTok{plt.stem(x, pmf\_emp, basefmt}\OperatorTok{=}\StringTok{" "}\NormalTok{)}
\NormalTok{plt.xlabel(}\StringTok{"x"}\NormalTok{)}
\NormalTok{plt.ylabel(}\StringTok{"p(x)"}\NormalTok{)}
\NormalTok{plt.title(}\StringTok{"PMF empiris dari histogram"}\NormalTok{)}
\NormalTok{plt.grid(alpha}\OperatorTok{=}\FloatTok{0.3}\NormalTok{)}
\NormalTok{plt.show()}
\end{Highlighting}
\end{Shaded}

\pandocbounded{\includegraphics[keepaspectratio]{03-distribusi-diskrit_files/figure-pdf/cell-4-output-1.pdf}}

\begin{Shaded}
\begin{Highlighting}[]
\NormalTok{plt.figure(figsize}\OperatorTok{=}\NormalTok{(}\DecValTok{8}\NormalTok{,}\DecValTok{4}\NormalTok{))}
\NormalTok{plt.step(x, cdf\_emp, where}\OperatorTok{=}\StringTok{\textquotesingle{}post\textquotesingle{}}\NormalTok{)}
\NormalTok{plt.xlabel(}\StringTok{"x"}\NormalTok{)}
\NormalTok{plt.ylabel(}\StringTok{"F(x)"}\NormalTok{)}
\NormalTok{plt.title(}\StringTok{"CDF empiris"}\NormalTok{)}
\NormalTok{plt.ylim(}\OperatorTok{{-}}\FloatTok{0.05}\NormalTok{, }\FloatTok{1.05}\NormalTok{)}
\NormalTok{plt.grid(alpha}\OperatorTok{=}\FloatTok{0.3}\NormalTok{)}
\NormalTok{plt.show()}
\end{Highlighting}
\end{Shaded}

\pandocbounded{\includegraphics[keepaspectratio]{03-distribusi-diskrit_files/figure-pdf/cell-5-output-1.pdf}}

\subsubsection{Interpretasi}\label{interpretasi-6}

Distribusi diskrit kustom sangat berguna ketika: - kita punya data
nyata, - tetapi belum ingin atau belum bisa memaksakan model teoritis
tertentu, - dan kita tetap ingin menghitung ekspektasi, varians, serta
peluang kumulatif.

\subsection{Ringkasan mini}\label{ringkasan-mini}

\begin{itemize}
\tightlist
\item
  Histogram dapat menjadi PMF empiris
\item
  Dari PMF empiris kita bisa hitung mean, variance, dan CDF
\item
  Ini adalah langkah awal yang penting sebelum memilih model yang lebih
  teoritis
\end{itemize}

\section{3.3 Discrete Uniform
Distribution}\label{discrete-uniform-distribution}

\subsection{K --- Konteks}\label{k-konteks-11}

Ada situasi di mana semua nilai mungkin dianggap sama-sama mungkin.
Misalnya: - hasil lemparan dadu fair, - memilih satu nomor secara acak
dari 1 sampai 10, - memilih satu hari dari 7 hari secara uniform.

\subsection{M --- Model}\label{m-model-11}

Jika \(X\) uniform diskrit pada \(\{a, a+1, \dots, b\}\), maka:

\[
P(X=x)=\frac{1}{b-a+1}, \qquad x=a,a+1,\dots,b
\]

\subsection{Q --- Questions}\label{q-questions-11}

\begin{enumerate}
\def\labelenumi{\arabic{enumi}.}
\tightlist
\item
  Berapa peluang satu nilai tertentu?
\item
  Berapa ekspektasi dan variansnya?
\item
  Kapan asumsi ``semua sama mungkin'' masuk akal?
\end{enumerate}

\subsection{A --- Apply}\label{a-apply-11}

\begin{Shaded}
\begin{Highlighting}[]
\NormalTok{x }\OperatorTok{=}\NormalTok{ np.arange(}\DecValTok{1}\NormalTok{, }\DecValTok{7}\NormalTok{)  }\CommentTok{\#\# dadu fair}
\NormalTok{pmf }\OperatorTok{=}\NormalTok{ np.ones\_like(x) }\OperatorTok{/} \BuiltInTok{len}\NormalTok{(x)}

\BuiltInTok{print}\NormalTok{(}\StringTok{"Mean ="}\NormalTok{, np.}\BuiltInTok{sum}\NormalTok{(x }\OperatorTok{*}\NormalTok{ pmf))}
\BuiltInTok{print}\NormalTok{(}\StringTok{"Variance ="}\NormalTok{, np.}\BuiltInTok{sum}\NormalTok{((x }\OperatorTok{{-}}\NormalTok{ np.}\BuiltInTok{sum}\NormalTok{(x}\OperatorTok{*}\NormalTok{pmf))}\OperatorTok{**}\DecValTok{2} \OperatorTok{*}\NormalTok{ pmf))}
\end{Highlighting}
\end{Shaded}

\begin{verbatim}
Mean = 3.5
Variance = 2.9166666666666665
\end{verbatim}

\begin{Shaded}
\begin{Highlighting}[]
\NormalTok{plt.figure(figsize}\OperatorTok{=}\NormalTok{(}\DecValTok{7}\NormalTok{,}\DecValTok{4}\NormalTok{))}
\NormalTok{plt.bar(x, pmf)}
\NormalTok{plt.xlabel(}\StringTok{"Mata dadu"}\NormalTok{)}
\NormalTok{plt.ylabel(}\StringTok{"Probabilitas"}\NormalTok{)}
\NormalTok{plt.title(}\StringTok{"Distribusi uniform diskrit pada dadu fair"}\NormalTok{)}
\NormalTok{plt.grid(alpha}\OperatorTok{=}\FloatTok{0.3}\NormalTok{)}
\NormalTok{plt.show()}
\end{Highlighting}
\end{Shaded}

\pandocbounded{\includegraphics[keepaspectratio]{03-distribusi-diskrit_files/figure-pdf/cell-7-output-1.pdf}}

\subsubsection{Hasil teoritis}\label{hasil-teoritis}

Untuk uniform diskrit pada \(\{1,2,\dots,n\}\):

\[
E[X] = \frac{n+1}{2}
\]

\[
\mathrm{Var}(X) = \frac{n^2-1}{12}
\]

\subsubsection{Interpretasi keputusan}\label{interpretasi-keputusan-3}

Distribusi uniform diskrit dipakai ketika tidak ada alasan untuk memberi
bobot lebih pada satu nilai dibanding yang lain.

\subsection{Latihan singkat}\label{latihan-singkat-6}

\begin{enumerate}
\def\labelenumi{\arabic{enumi}.}
\tightlist
\item
  Hitung mean dan variance untuk uniform diskrit pada
  \(\{3,4,\dots,10\}\).
\item
  Simulasikan 10000 sampel dari distribusi itu.
\item
  Beri contoh keputusan nyata yang cocok dimodelkan uniform diskrit.
\end{enumerate}

\section{3.4 Bernoulli Distribution}\label{bernoulli-distribution}

\subsection{K --- Konteks}\label{k-konteks-12}

Banyak keputusan nyata berujung pada dua kemungkinan: - sukses / gagal,
- cacat / tidak cacat, - hujan / tidak hujan, - lolos / tidak lolos, -
klik / tidak klik.

\subsection{M --- Model}\label{m-model-12}

Random variable Bernoulli memiliki dua nilai: \[
X \in \{0,1\}
\]

dengan: \[
P(X=1)=p, \qquad P(X=0)=1-p
\]

\subsection{Q --- Questions}\label{q-questions-12}

\begin{enumerate}
\def\labelenumi{\arabic{enumi}.}
\tightlist
\item
  Apa arti parameter \(p\)?
\item
  Berapa ekspektasi dan variansnya?
\item
  Masalah apa yang paling cocok dimodelkan Bernoulli?
\end{enumerate}

\subsection{A --- Apply}\label{a-apply-12}

\begin{Shaded}
\begin{Highlighting}[]
\NormalTok{p }\OperatorTok{=} \FloatTok{0.7}
\NormalTok{x }\OperatorTok{=}\NormalTok{ np.array([}\DecValTok{0}\NormalTok{, }\DecValTok{1}\NormalTok{])}
\NormalTok{pmf }\OperatorTok{=}\NormalTok{ np.array([}\DecValTok{1}\OperatorTok{{-}}\NormalTok{p, p])}

\NormalTok{mean }\OperatorTok{=}\NormalTok{ np.}\BuiltInTok{sum}\NormalTok{(x }\OperatorTok{*}\NormalTok{ pmf)}
\NormalTok{var }\OperatorTok{=}\NormalTok{ np.}\BuiltInTok{sum}\NormalTok{((x }\OperatorTok{{-}}\NormalTok{ mean)}\OperatorTok{**}\DecValTok{2} \OperatorTok{*}\NormalTok{ pmf)}

\BuiltInTok{print}\NormalTok{(}\StringTok{"Mean ="}\NormalTok{, mean)}
\BuiltInTok{print}\NormalTok{(}\StringTok{"Variance ="}\NormalTok{, var)}
\end{Highlighting}
\end{Shaded}

\begin{verbatim}
Mean = 0.7
Variance = 0.21000000000000002
\end{verbatim}

\begin{Shaded}
\begin{Highlighting}[]
\NormalTok{plt.figure(figsize}\OperatorTok{=}\NormalTok{(}\DecValTok{6}\NormalTok{,}\DecValTok{4}\NormalTok{))}
\NormalTok{plt.bar(x, pmf)}
\NormalTok{plt.xticks([}\DecValTok{0}\NormalTok{,}\DecValTok{1}\NormalTok{], [}\StringTok{"0 (gagal)"}\NormalTok{, }\StringTok{"1 (sukses)"}\NormalTok{])}
\NormalTok{plt.ylabel(}\StringTok{"Probabilitas"}\NormalTok{)}
\NormalTok{plt.title(}\StringTok{"Distribusi Bernoulli"}\NormalTok{)}
\NormalTok{plt.grid(alpha}\OperatorTok{=}\FloatTok{0.3}\NormalTok{)}
\NormalTok{plt.show()}
\end{Highlighting}
\end{Shaded}

\pandocbounded{\includegraphics[keepaspectratio]{03-distribusi-diskrit_files/figure-pdf/cell-9-output-1.pdf}}

\subsubsection{Hasil teoritis}\label{hasil-teoritis-1}

\[
E[X]=p
\]

\[
\mathrm{Var}(X)=p(1-p)
\]

\subsubsection{Interpretasi}\label{interpretasi-7}

Bernoulli adalah batu bata paling dasar bagi banyak distribusi lain.
Binomial, misalnya, dibangun dari pengulangan Bernoulli.

\subsection{Latihan singkat}\label{latihan-singkat-7}

\begin{enumerate}
\def\labelenumi{\arabic{enumi}.}
\tightlist
\item
  Jika peluang klik iklan 0.03, tulis model Bernoullinya.
\item
  Hitung mean dan variance-nya.
\item
  Mengapa variance Bernoulli maksimum saat \(p=0.5\)?
\end{enumerate}

\section{3.5 Binomial Distribution}\label{binomial-distribution}

\subsection{K --- Konteks}\label{k-konteks-13}

Sekarang kita beralih dari satu percobaan Bernoulli ke \(n\) percobaan
Bernoulli independen.\\
Contoh: - dari 100 produk, berapa yang cacat? - dari 20 email, berapa
yang dibuka? - dari 10 lemparan koin, berapa Head?

\subsection{M --- Model}\label{m-model-13}

Jika: - ada \(n\) percobaan independen, - setiap percobaan punya peluang
sukses \(p\), - \(X\) = jumlah sukses,

maka:

\[
X \sim \mathrm{Binomial}(n,p)
\]

dengan PMF:

\[
P(X=k)=\binom{n}{k} p^k (1-p)^{n-k}, \qquad k=0,1,\dots,n
\]

\subsection{Q --- Questions}\label{q-questions-13}

\begin{enumerate}
\def\labelenumi{\arabic{enumi}.}
\tightlist
\item
  Bagaimana menghitung peluang tepat \(k\) sukses?
\item
  Berapa ekspektasi dan varians?
\item
  Kapan Binomial tepat dipakai?
\end{enumerate}

\subsection{A --- Apply}\label{a-apply-13}

\begin{Shaded}
\begin{Highlighting}[]
\NormalTok{n }\OperatorTok{=} \DecValTok{10}
\NormalTok{p }\OperatorTok{=} \FloatTok{0.3}
\NormalTok{k }\OperatorTok{=}\NormalTok{ np.arange(}\DecValTok{0}\NormalTok{, n}\OperatorTok{+}\DecValTok{1}\NormalTok{)}

\NormalTok{pmf }\OperatorTok{=}\NormalTok{ stats.binom.pmf(k, n}\OperatorTok{=}\NormalTok{n, p}\OperatorTok{=}\NormalTok{p)}

\NormalTok{plt.figure(figsize}\OperatorTok{=}\NormalTok{(}\DecValTok{8}\NormalTok{,}\DecValTok{4}\NormalTok{))}
\NormalTok{plt.stem(k, pmf, basefmt}\OperatorTok{=}\StringTok{" "}\NormalTok{)}
\NormalTok{plt.xlabel(}\StringTok{"Jumlah sukses"}\NormalTok{)}
\NormalTok{plt.ylabel(}\StringTok{"P(X=k)"}\NormalTok{)}
\NormalTok{plt.title(}\StringTok{"PMF Binomial(n=10, p=0.3)"}\NormalTok{)}
\NormalTok{plt.grid(alpha}\OperatorTok{=}\FloatTok{0.3}\NormalTok{)}
\NormalTok{plt.show()}
\end{Highlighting}
\end{Shaded}

\pandocbounded{\includegraphics[keepaspectratio]{03-distribusi-diskrit_files/figure-pdf/cell-10-output-1.pdf}}

\begin{Shaded}
\begin{Highlighting}[]
\BuiltInTok{print}\NormalTok{(}\StringTok{"P(X=3) ="}\NormalTok{, stats.binom.pmf(}\DecValTok{3}\NormalTok{, n}\OperatorTok{=}\DecValTok{10}\NormalTok{, p}\OperatorTok{=}\FloatTok{0.3}\NormalTok{))}
\BuiltInTok{print}\NormalTok{(}\StringTok{"P(X\textless{}=3) ="}\NormalTok{, stats.binom.cdf(}\DecValTok{3}\NormalTok{, n}\OperatorTok{=}\DecValTok{10}\NormalTok{, p}\OperatorTok{=}\FloatTok{0.3}\NormalTok{))}
\BuiltInTok{print}\NormalTok{(}\StringTok{"E[X] ="}\NormalTok{, n}\OperatorTok{*}\NormalTok{p)}
\BuiltInTok{print}\NormalTok{(}\StringTok{"Var(X) ="}\NormalTok{, n}\OperatorTok{*}\NormalTok{p}\OperatorTok{*}\NormalTok{(}\DecValTok{1}\OperatorTok{{-}}\NormalTok{p))}
\end{Highlighting}
\end{Shaded}

\begin{verbatim}
P(X=3) = 0.2668279319999998
P(X<=3) = 0.6496107184000001
E[X] = 3.0
Var(X) = 2.0999999999999996
\end{verbatim}

\subsubsection{Hasil teoritis}\label{hasil-teoritis-2}

\[
E[X]=np
\]

\[
\mathrm{Var}(X)=np(1-p)
\]

\subsubsection{Interpretasi keputusan}\label{interpretasi-keputusan-4}

Gunakan Binomial ketika: - jumlah percobaan tetap, - tiap percobaan
independen, - peluang sukses tetap, - yang kita hitung adalah
\textbf{jumlah sukses}.

\subsubsection{Contoh keputusan QC}\label{contoh-keputusan-qc}

Jika sebuah pabrik mengklaim tingkat cacat 1\%, maka banyak cacat pada
sampel 100 unit cocok dimodelkan Binomial(100, 0.01). Dari sana kita
bisa menghitung apakah hasil inspeksi masih masuk akal.

\subsection{Latihan singkat}\label{latihan-singkat-8}

\begin{enumerate}
\def\labelenumi{\arabic{enumi}.}
\tightlist
\item
  Hitung peluang tepat 2 cacat dari 100 unit jika \(p=0.01\).
\item
  Hitung peluang paling banyak 2 cacat.
\item
  Simulasikan dan bandingkan hasil empiris dengan PMF teoritis.
\end{enumerate}

\section{3.6 Geometric Distribution}\label{geometric-distribution}

\subsection{K --- Konteks}\label{k-konteks-14}

Kadang bukan jumlah sukses yang ingin kita hitung, tetapi \textbf{berapa
lama menunggu sampai sukses pertama}.

Contoh: - berapa kali mencoba login sampai berhasil? - berapa banyak
panggilan sampai ada yang menjawab? - berapa kali inspeksi sampai
menemukan cacat pertama?

\subsection{M --- Model}\label{m-model-14}

Jika setiap percobaan independen dan peluang sukses tiap percobaan
adalah \(p\), maka random variable Geometric dapat didefinisikan
sebagai:

\[
X = \text{jumlah percobaan sampai sukses pertama}
\]

dengan PMF:

\[
P(X=k)=(1-p)^{k-1}p, \qquad k=1,2,3,\dots
\]

\subsection{Q --- Questions}\label{q-questions-14}

\begin{enumerate}
\def\labelenumi{\arabic{enumi}.}
\tightlist
\item
  Berapa peluang sukses pertama terjadi pada percobaan ke-\(k\)?
\item
  Berapa rata-rata jumlah percobaan sampai sukses pertama?
\item
  Mengapa distribusi ini punya sifat memoryless?
\end{enumerate}

\subsection{A --- Apply}\label{a-apply-14}

\begin{Shaded}
\begin{Highlighting}[]
\NormalTok{p }\OperatorTok{=} \FloatTok{0.25}
\NormalTok{k }\OperatorTok{=}\NormalTok{ np.arange(}\DecValTok{1}\NormalTok{, }\DecValTok{15}\NormalTok{)}

\NormalTok{pmf }\OperatorTok{=}\NormalTok{ stats.geom.pmf(k, p)}

\NormalTok{plt.figure(figsize}\OperatorTok{=}\NormalTok{(}\DecValTok{8}\NormalTok{,}\DecValTok{4}\NormalTok{))}
\NormalTok{plt.stem(k, pmf, basefmt}\OperatorTok{=}\StringTok{" "}\NormalTok{)}
\NormalTok{plt.xlabel(}\StringTok{"k = percobaan sampai sukses pertama"}\NormalTok{)}
\NormalTok{plt.ylabel(}\StringTok{"P(X=k)"}\NormalTok{)}
\NormalTok{plt.title(}\StringTok{"PMF Geometric(p=0.25)"}\NormalTok{)}
\NormalTok{plt.grid(alpha}\OperatorTok{=}\FloatTok{0.3}\NormalTok{)}
\NormalTok{plt.show()}
\end{Highlighting}
\end{Shaded}

\pandocbounded{\includegraphics[keepaspectratio]{03-distribusi-diskrit_files/figure-pdf/cell-12-output-1.pdf}}

\begin{Shaded}
\begin{Highlighting}[]
\BuiltInTok{print}\NormalTok{(}\StringTok{"P(X=4) ="}\NormalTok{, stats.geom.pmf(}\DecValTok{4}\NormalTok{, p))}
\BuiltInTok{print}\NormalTok{(}\StringTok{"P(X\textless{}=4) ="}\NormalTok{, stats.geom.cdf(}\DecValTok{4}\NormalTok{, p))}
\BuiltInTok{print}\NormalTok{(}\StringTok{"Mean ="}\NormalTok{, }\DecValTok{1}\OperatorTok{/}\NormalTok{p)}
\BuiltInTok{print}\NormalTok{(}\StringTok{"Variance ="}\NormalTok{, (}\DecValTok{1}\OperatorTok{{-}}\NormalTok{p)}\OperatorTok{/}\NormalTok{(p}\OperatorTok{**}\DecValTok{2}\NormalTok{))}
\end{Highlighting}
\end{Shaded}

\begin{verbatim}
P(X=4) = 0.10546875
P(X<=4) = 0.68359375
Mean = 4.0
Variance = 12.0
\end{verbatim}

\subsubsection{Hasil teoritis}\label{hasil-teoritis-3}

\[
E[X]=\frac{1}{p}
\]

\[
\mathrm{Var}(X)=\frac{1-p}{p^2}
\]

\subsubsection{Sifat memoryless}\label{sifat-memoryless}

Distribusi Geometric memenuhi:

\[
P(X>m+n \mid X>m)=P(X>n)
\]

Artinya, jika kita sudah gagal sampai titik tertentu, peluang menunggu
lebih lama lagi tidak ``mengingat'' masa lalu. Sifat ini akan muncul
lagi pada distribusi Eksponensial di bab kontinu.

\subsection{Latihan singkat}\label{latihan-singkat-9}

\begin{enumerate}
\def\labelenumi{\arabic{enumi}.}
\tightlist
\item
  Jika peluang sukses 0.2, berapa rata-rata jumlah percobaan sampai
  sukses pertama?
\item
  Hitung peluang perlu lebih dari 5 percobaan.
\item
  Verifikasi sifat memoryless dengan simulasi.
\end{enumerate}

\section{3.7 Poisson Distribution}\label{poisson-distribution}

\subsection{K --- Konteks}\label{k-konteks-15}

Poisson sangat penting untuk memodelkan \textbf{jumlah kejadian} dalam
suatu interval tetap, ketika kejadian itu: - relatif jarang, -
independen, - dan rata-rata kejadian per interval diketahui.

Contoh: - jumlah panggilan masuk per jam, - jumlah pelanggan datang per
10 menit, - jumlah cacat pada lembar besar material, - jumlah kendaraan
lewat per menit.

\subsection{M --- Model}\label{m-model-15}

Jika \(X\) = jumlah kejadian dalam interval, dengan rata-rata
\(\lambda\), maka:

\[
X \sim \mathrm{Poisson}(\lambda)
\]

dan PMF-nya:

\[
P(X=k)=e^{-\lambda}\frac{\lambda^k}{k!}, \qquad k=0,1,2,\dots
\]

\subsection{Q --- Questions}\label{q-questions-15}

\begin{enumerate}
\def\labelenumi{\arabic{enumi}.}
\tightlist
\item
  Berapa peluang tepat \(k\) kejadian?
\item
  Mengapa mean dan variance Poisson sama?
\item
  Kapan Poisson menjadi pendekatan yang baik untuk Binomial?
\end{enumerate}

\subsection{A --- Apply}\label{a-apply-15}

\begin{Shaded}
\begin{Highlighting}[]
\NormalTok{lam }\OperatorTok{=} \DecValTok{4}
\NormalTok{k }\OperatorTok{=}\NormalTok{ np.arange(}\DecValTok{0}\NormalTok{, }\DecValTok{15}\NormalTok{)}

\NormalTok{pmf }\OperatorTok{=}\NormalTok{ stats.poisson.pmf(k, mu}\OperatorTok{=}\NormalTok{lam)}

\NormalTok{plt.figure(figsize}\OperatorTok{=}\NormalTok{(}\DecValTok{8}\NormalTok{,}\DecValTok{4}\NormalTok{))}
\NormalTok{plt.stem(k, pmf, basefmt}\OperatorTok{=}\StringTok{" "}\NormalTok{)}
\NormalTok{plt.xlabel(}\StringTok{"Jumlah kejadian"}\NormalTok{)}
\NormalTok{plt.ylabel(}\StringTok{"P(X=k)"}\NormalTok{)}
\NormalTok{plt.title(}\StringTok{"PMF Poisson(λ=4)"}\NormalTok{)}
\NormalTok{plt.grid(alpha}\OperatorTok{=}\FloatTok{0.3}\NormalTok{)}
\NormalTok{plt.show()}
\end{Highlighting}
\end{Shaded}

\pandocbounded{\includegraphics[keepaspectratio]{03-distribusi-diskrit_files/figure-pdf/cell-14-output-1.pdf}}

\begin{Shaded}
\begin{Highlighting}[]
\BuiltInTok{print}\NormalTok{(}\StringTok{"P(X=3) ="}\NormalTok{, stats.poisson.pmf(}\DecValTok{3}\NormalTok{, mu}\OperatorTok{=}\NormalTok{lam))}
\BuiltInTok{print}\NormalTok{(}\StringTok{"P(X\textless{}=3) ="}\NormalTok{, stats.poisson.cdf(}\DecValTok{3}\NormalTok{, mu}\OperatorTok{=}\NormalTok{lam))}
\BuiltInTok{print}\NormalTok{(}\StringTok{"Mean ="}\NormalTok{, lam)}
\BuiltInTok{print}\NormalTok{(}\StringTok{"Variance ="}\NormalTok{, lam)}
\end{Highlighting}
\end{Shaded}

\begin{verbatim}
P(X=3) = 0.19536681481316454
P(X<=3) = 0.43347012036670896
Mean = 4
Variance = 4
\end{verbatim}

\subsubsection{Hasil teoritis}\label{hasil-teoritis-4}

\[
E[X]=\lambda
\]

\[
\mathrm{Var}(X)=\lambda
\]

\subsubsection{Poisson sebagai pendekatan
Binomial}\label{poisson-sebagai-pendekatan-binomial}

Jika: - \(n\) besar, - \(p\) kecil, - \(\lambda = np\) moderat,

maka: \[
\mathrm{Binomial}(n,p) \approx \mathrm{Poisson}(\lambda)
\]

Mari bandingkan.

\begin{Shaded}
\begin{Highlighting}[]
\NormalTok{n }\OperatorTok{=} \DecValTok{100}
\NormalTok{p }\OperatorTok{=} \FloatTok{0.02}
\NormalTok{lam }\OperatorTok{=}\NormalTok{ n }\OperatorTok{*}\NormalTok{ p}
\NormalTok{x }\OperatorTok{=}\NormalTok{ np.arange(}\DecValTok{0}\NormalTok{, }\DecValTok{10}\NormalTok{)}

\NormalTok{pmf\_binom }\OperatorTok{=}\NormalTok{ stats.binom.pmf(x, n}\OperatorTok{=}\NormalTok{n, p}\OperatorTok{=}\NormalTok{p)}
\NormalTok{pmf\_pois }\OperatorTok{=}\NormalTok{ stats.poisson.pmf(x, mu}\OperatorTok{=}\NormalTok{lam)}

\NormalTok{plt.figure(figsize}\OperatorTok{=}\NormalTok{(}\DecValTok{8}\NormalTok{,}\DecValTok{4}\NormalTok{))}
\NormalTok{plt.stem(x, pmf\_binom, linefmt}\OperatorTok{=}\StringTok{\textquotesingle{}C0{-}\textquotesingle{}}\NormalTok{, markerfmt}\OperatorTok{=}\StringTok{\textquotesingle{}C0o\textquotesingle{}}\NormalTok{, basefmt}\OperatorTok{=}\StringTok{" "}\NormalTok{, label}\OperatorTok{=}\StringTok{\textquotesingle{}Binomial\textquotesingle{}}\NormalTok{)}
\NormalTok{plt.stem(x}\OperatorTok{+}\FloatTok{0.1}\NormalTok{, pmf\_pois, linefmt}\OperatorTok{=}\StringTok{\textquotesingle{}C1{-}\textquotesingle{}}\NormalTok{, markerfmt}\OperatorTok{=}\StringTok{\textquotesingle{}C1s\textquotesingle{}}\NormalTok{, basefmt}\OperatorTok{=}\StringTok{" "}\NormalTok{, label}\OperatorTok{=}\StringTok{\textquotesingle{}Poisson approx\textquotesingle{}}\NormalTok{)}
\NormalTok{plt.xlabel(}\StringTok{"x"}\NormalTok{)}
\NormalTok{plt.ylabel(}\StringTok{"P(X=x)"}\NormalTok{)}
\NormalTok{plt.title(}\StringTok{"Pendekatan Poisson terhadap Binomial"}\NormalTok{)}
\NormalTok{plt.legend()}
\NormalTok{plt.grid(alpha}\OperatorTok{=}\FloatTok{0.3}\NormalTok{)}
\NormalTok{plt.show()}
\end{Highlighting}
\end{Shaded}

\pandocbounded{\includegraphics[keepaspectratio]{03-distribusi-diskrit_files/figure-pdf/cell-16-output-1.pdf}}

\subsubsection{Interpretasi keputusan}\label{interpretasi-keputusan-5}

Poisson cocok untuk menghitung volume kedatangan, insiden, atau kejadian
jarang dalam interval. Dalam keputusan operasional, ini sangat berguna
untuk: - kapasitas layanan, - staffing, - quality control, -
reliability.

\subsection{Latihan singkat}\label{latihan-singkat-10}

\begin{enumerate}
\def\labelenumi{\arabic{enumi}.}
\tightlist
\item
  Jika rata-rata pelanggan datang 5 per jam, hitung peluang tepat 7
  pelanggan.
\item
  Hitung peluang paling banyak 3 pelanggan.
\item
  Bandingkan hasil Poisson dan Binomial pada kasus kejadian langka.
\end{enumerate}

\section{3.8 Memilih Distribusi yang
Tepat}\label{memilih-distribusi-yang-tepat}

Salah satu kemampuan terpenting bukan hanya menghitung, tetapi memilih
model yang masuk akal. Berikut panduan ringkas.

\subsection{Bernoulli}\label{bernoulli}

Gunakan jika: - hanya ada dua hasil, - fokus pada satu percobaan.

Contoh: - klik / tidak klik, - cacat / tidak cacat.

\subsection{Binomial}\label{binomial}

Gunakan jika: - ada \(n\) percobaan independen, - peluang sukses tetap,
- fokus pada jumlah sukses.

Contoh: - jumlah cacat dari 100 unit.

\subsection{Geometric}\label{geometric}

Gunakan jika: - fokus pada berapa lama menunggu sampai sukses pertama, -
percobaan independen, - peluang sukses tetap.

Contoh: - jumlah percobaan sampai login berhasil.

\subsection{Poisson}\label{poisson}

Gunakan jika: - fokus pada jumlah kejadian dalam interval, - rata-rata
kejadian diketahui, - kejadian independen dan relatif jarang.

Contoh: - jumlah pasien datang per jam.

\subsection{Uniform diskrit}\label{uniform-diskrit}

Gunakan jika: - semua nilai dianggap sama mungkin.

Contoh: - hasil dadu fair.

\section{3.9 Mini-Case KMQA 1 --- Produk
Cacat}\label{mini-case-kmqa-1-produk-cacat}

\subsection{K --- Konteks}\label{k-konteks-16}

Pabrik mengklaim tingkat cacat 1\%. Tim QC mengambil sampel 100 unit.

\subsection{M --- Model}\label{m-model-16}

Jika \(X\) = jumlah cacat, maka: \[
X \sim \text{Binomial}(100, 0.01)
\]

Pendekatan Poisson: \[
X \approx \text{Poisson}(1)
\]

\subsection{Q --- Questions}\label{q-questions-16}

\begin{enumerate}
\def\labelenumi{\arabic{enumi}.}
\tightlist
\item
  Berapa peluang tepat 2 cacat?
\item
  Berapa peluang paling banyak 2 cacat?
\item
  Jika ditemukan 6 cacat, apakah itu masih terasa wajar?
\end{enumerate}

\subsection{A --- Apply}\label{a-apply-16}

\begin{Shaded}
\begin{Highlighting}[]
\NormalTok{n }\OperatorTok{=} \DecValTok{100}
\NormalTok{p }\OperatorTok{=} \FloatTok{0.01}
\NormalTok{lam }\OperatorTok{=}\NormalTok{ n}\OperatorTok{*}\NormalTok{p}

\BuiltInTok{print}\NormalTok{(}\StringTok{"P(X=2) Binomial ="}\NormalTok{, stats.binom.pmf(}\DecValTok{2}\NormalTok{, n, p))}
\BuiltInTok{print}\NormalTok{(}\StringTok{"P(X=2) Poisson  ="}\NormalTok{, stats.poisson.pmf(}\DecValTok{2}\NormalTok{, lam))}

\BuiltInTok{print}\NormalTok{(}\StringTok{"P(X\textless{}=2) Binomial ="}\NormalTok{, stats.binom.cdf(}\DecValTok{2}\NormalTok{, n, p))}
\BuiltInTok{print}\NormalTok{(}\StringTok{"P(X\textless{}=2) Poisson  ="}\NormalTok{, stats.poisson.cdf(}\DecValTok{2}\NormalTok{, lam))}

\BuiltInTok{print}\NormalTok{(}\StringTok{"P(X\textgreater{}=6) Binomial ="}\NormalTok{, }\DecValTok{1} \OperatorTok{{-}}\NormalTok{ stats.binom.cdf(}\DecValTok{5}\NormalTok{, n, p))}
\end{Highlighting}
\end{Shaded}

\begin{verbatim}
P(X=2) Binomial = 0.18486481882486347
P(X=2) Poisson  = 0.18393972058572114
P(X<=2) Binomial = 0.9206267977478195
P(X<=2) Poisson  = 0.9196986029286058
P(X>=6) Binomial = 0.0005345344639929861
\end{verbatim}

\subsubsection{Interpretasi}\label{interpretasi-8}

Kalau peluang mendapatkan 6 atau lebih cacat ternyata sangat kecil, maka
tim QC pantas curiga bahwa klaim tingkat cacat 1\% mungkin tidak lagi
valid.

\section{3.10 Mini-Case KMQA 2 --- Kedatangan
Pelanggan}\label{mini-case-kmqa-2-kedatangan-pelanggan}

\subsection{K --- Konteks}\label{k-konteks-17}

Sebuah kios kopi ingin tahu distribusi jumlah pelanggan per jam agar
bisa menyiapkan stok dan staf.

\subsection{M --- Model}\label{m-model-17}

Jika \(X\) = jumlah pelanggan datang per jam, gunakan: \[
X \sim \text{Poisson}(\lambda)
\]

\subsection{Q --- Questions}\label{q-questions-17}

\begin{enumerate}
\def\labelenumi{\arabic{enumi}.}
\tightlist
\item
  Berapa peluang tepat 10 pelanggan datang?
\item
  Berapa peluang lebih dari 12 pelanggan datang?
\item
  Apa arti mean dan variance dalam konteks operasional?
\end{enumerate}

\subsection{A --- Apply}\label{a-apply-17}

\begin{Shaded}
\begin{Highlighting}[]
\NormalTok{lam }\OperatorTok{=} \DecValTok{8}

\BuiltInTok{print}\NormalTok{(}\StringTok{"P(X=10) ="}\NormalTok{, stats.poisson.pmf(}\DecValTok{10}\NormalTok{, mu}\OperatorTok{=}\NormalTok{lam))}
\BuiltInTok{print}\NormalTok{(}\StringTok{"P(X\textgreater{}12) ="}\NormalTok{, }\DecValTok{1} \OperatorTok{{-}}\NormalTok{ stats.poisson.cdf(}\DecValTok{12}\NormalTok{, mu}\OperatorTok{=}\NormalTok{lam))}
\end{Highlighting}
\end{Shaded}

\begin{verbatim}
P(X=10) = 0.09926153383153544
P(X>12) = 0.06379719673656181
\end{verbatim}

\subsubsection{Interpretasi}\label{interpretasi-9}

Distribusi Poisson membantu kios kopi menilai: - seberapa sering hari
ramai, - seberapa besar stok minimum masuk akal, - dan berapa kapasitas
pelayanan yang aman.

\section{3.11 Mini-Case KMQA 3 --- Login Sampai
Berhasil}\label{mini-case-kmqa-3-login-sampai-berhasil}

\subsection{K --- Konteks}\label{k-konteks-18}

Sebuah sistem login punya probabilitas sukses 0.8 untuk setiap percobaan
independen.

\subsection{M --- Model}\label{m-model-18}

Jika \(X\) = jumlah percobaan sampai sukses pertama: \[
X \sim \text{Geometric}(p=0.8)
\]

\subsection{Q --- Questions}\label{q-questions-18}

\begin{enumerate}
\def\labelenumi{\arabic{enumi}.}
\tightlist
\item
  Berapa peluang sukses pada percobaan pertama?
\item
  Berapa peluang butuh lebih dari 3 percobaan?
\item
  Berapa rata-rata jumlah percobaan sampai berhasil?
\end{enumerate}

\subsection{A --- Apply}\label{a-apply-18}

\begin{Shaded}
\begin{Highlighting}[]
\NormalTok{p }\OperatorTok{=} \FloatTok{0.8}
\BuiltInTok{print}\NormalTok{(}\StringTok{"P(X=1) ="}\NormalTok{, stats.geom.pmf(}\DecValTok{1}\NormalTok{, p))}
\BuiltInTok{print}\NormalTok{(}\StringTok{"P(X\textgreater{}3) ="}\NormalTok{, }\DecValTok{1} \OperatorTok{{-}}\NormalTok{ stats.geom.cdf(}\DecValTok{3}\NormalTok{, p))}
\BuiltInTok{print}\NormalTok{(}\StringTok{"Mean ="}\NormalTok{, }\DecValTok{1}\OperatorTok{/}\NormalTok{p)}
\end{Highlighting}
\end{Shaded}

\begin{verbatim}
P(X=1) = 0.8
P(X>3) = 0.008000000000000007
Mean = 1.25
\end{verbatim}

\section{3.12 Python Toolbox untuk Bab
Ini}\label{python-toolbox-untuk-bab-ini}

Berikut beberapa fungsi Python yang sangat sering dipakai.

\begin{Shaded}
\begin{Highlighting}[]
\ImportTok{from}\NormalTok{ scipy }\ImportTok{import}\NormalTok{ stats}
\ImportTok{import}\NormalTok{ numpy }\ImportTok{as}\NormalTok{ np}
\ImportTok{import}\NormalTok{ matplotlib.pyplot }\ImportTok{as}\NormalTok{ plt}
\end{Highlighting}
\end{Shaded}

\subsection{Bernoulli}\label{bernoulli-1}

\begin{Shaded}
\begin{Highlighting}[]
\NormalTok{stats.bernoulli.pmf([}\DecValTok{0}\NormalTok{,}\DecValTok{1}\NormalTok{], p}\OperatorTok{=}\FloatTok{0.3}\NormalTok{)}
\end{Highlighting}
\end{Shaded}

\begin{verbatim}
array([0.7, 0.3])
\end{verbatim}

\subsection{Binomial}\label{binomial-1}

\begin{Shaded}
\begin{Highlighting}[]
\NormalTok{stats.binom.pmf(}\DecValTok{3}\NormalTok{, n}\OperatorTok{=}\DecValTok{10}\NormalTok{, p}\OperatorTok{=}\FloatTok{0.4}\NormalTok{), stats.binom.cdf(}\DecValTok{3}\NormalTok{, n}\OperatorTok{=}\DecValTok{10}\NormalTok{, p}\OperatorTok{=}\FloatTok{0.4}\NormalTok{)}
\end{Highlighting}
\end{Shaded}

\begin{verbatim}
(np.float64(0.21499084799999987), np.float64(0.3822806016))
\end{verbatim}

\subsection{Geometric}\label{geometric-1}

\begin{Shaded}
\begin{Highlighting}[]
\NormalTok{stats.geom.pmf(}\DecValTok{4}\NormalTok{, p}\OperatorTok{=}\FloatTok{0.2}\NormalTok{), stats.geom.cdf(}\DecValTok{4}\NormalTok{, p}\OperatorTok{=}\FloatTok{0.2}\NormalTok{)}
\end{Highlighting}
\end{Shaded}

\begin{verbatim}
(np.float64(0.10240000000000003), np.float64(0.5904))
\end{verbatim}

\subsection{Poisson}\label{poisson-1}

\begin{Shaded}
\begin{Highlighting}[]
\NormalTok{stats.poisson.pmf(}\DecValTok{5}\NormalTok{, mu}\OperatorTok{=}\DecValTok{3}\NormalTok{), stats.poisson.cdf(}\DecValTok{5}\NormalTok{, mu}\OperatorTok{=}\DecValTok{3}\NormalTok{)}
\end{Highlighting}
\end{Shaded}

\begin{verbatim}
(np.float64(0.10081881344492458), np.float64(0.9160820579686966))
\end{verbatim}

\subsection{Simulasi}\label{simulasi}

\begin{Shaded}
\begin{Highlighting}[]
\NormalTok{rng }\OperatorTok{=}\NormalTok{ np.random.default\_rng(}\DecValTok{0}\NormalTok{)}
\NormalTok{rng.binomial(n}\OperatorTok{=}\DecValTok{10}\NormalTok{, p}\OperatorTok{=}\FloatTok{0.3}\NormalTok{, size}\OperatorTok{=}\DecValTok{10}\NormalTok{)}
\end{Highlighting}
\end{Shaded}

\begin{verbatim}
array([3, 2, 1, 0, 4, 5, 3, 4, 3, 5])
\end{verbatim}

\section{3.13 Kesalahan Umum yang Harus
Dihindari}\label{kesalahan-umum-yang-harus-dihindari}

\begin{enumerate}
\def\labelenumi{\arabic{enumi}.}
\item
  \textbf{Salah memilih model}\\
  Misalnya memakai Binomial padahal yang ditanya adalah waktu tunggu
  sampai sukses pertama.
\item
  \textbf{Mencampuradukkan ``jumlah sukses'' dan ``percobaan sampai
  sukses''}\\
  Ini perbedaan Binomial vs Geometric.
\item
  \textbf{Menganggap Poisson untuk semua count data}\\
  Tidak. Poisson cocok jika mean dan variance kurang lebih sejalan. Bila
  variance jauh lebih besar, model lain mungkin lebih cocok.
\item
  \textbf{Lupa memeriksa asumsi independensi dan peluang tetap}\\
  Binomial dan Geometric bergantung kuat pada asumsi ini.
\item
  \textbf{Menghafal rumus tanpa memahami konteks}\\
  Tujuan bab ini justru agar Anda tahu kapan suatu distribusi relevan.
\end{enumerate}

\section{3.14 Menyimpulkan Bab Ini}\label{menyimpulkan-bab-ini-2}

Distribusi diskrit adalah alat yang sangat kuat untuk memodelkan banyak
situasi nyata yang berupa hitungan. Kita telah melihat beberapa keluarga
penting:

\begin{itemize}
\tightlist
\item
  \textbf{uniform diskrit} untuk nilai-nilai yang sama mungkin,
\item
  \textbf{Bernoulli} untuk satu percobaan dua hasil,
\item
  \textbf{Binomial} untuk jumlah sukses dari banyak percobaan,
\item
  \textbf{Geometric} untuk menunggu sukses pertama,
\item
  \textbf{Poisson} untuk jumlah kejadian dalam suatu interval.
\end{itemize}

Yang terpenting bukan hanya bisa menghitung peluangnya, tetapi juga bisa
mengatakan: - apa random variable-nya, - apa arti parameternya, -
mengapa model itu masuk akal, - dan keputusan apa yang didukung oleh
model itu.

\section{3.15 Ringkasan Poin Inti}\label{ringkasan-poin-inti-2}

\begin{enumerate}
\def\labelenumi{\arabic{enumi}.}
\tightlist
\item
  Distribusi diskrit menjelaskan bagaimana peluang tersebar pada
  nilai-nilai hitungan.
\item
  Histogram empiris bisa menjadi langkah awal membangun PMF kustom.
\item
  \textbf{Uniform diskrit} dipakai bila semua nilai sama mungkin.
\item
  \textbf{Bernoulli} memodelkan satu percobaan dua hasil.
\item
  \textbf{Binomial} memodelkan jumlah sukses dari \(n\) percobaan
  independen.
\item
  \textbf{Geometric} memodelkan jumlah percobaan sampai sukses pertama.
\item
  \textbf{Poisson} memodelkan jumlah kejadian dalam interval tetap.
\item
  Memilih distribusi yang tepat adalah kemampuan yang lebih penting
  daripada sekadar menghafal rumus.
\end{enumerate}

\section{3.16 Latihan Bab 3}\label{latihan-bab-3}

\subsection{A. Konseptual}\label{a.-konseptual-2}

\begin{enumerate}
\def\labelenumi{\arabic{enumi}.}
\tightlist
\item
  Jelaskan perbedaan Bernoulli, Binomial, dan Geometric.
\item
  Mengapa Poisson sering dipakai untuk jumlah kejadian per interval?
\item
  Kapan uniform diskrit masuk akal, dan kapan tidak?
\end{enumerate}

\subsection{B. Hitungan}\label{b.-hitungan-1}

\begin{enumerate}
\def\labelenumi{\arabic{enumi}.}
\tightlist
\item
  Sebuah email dibuka dengan peluang 0.2. Modelkan satu email sebagai
  Bernoulli dan hitung mean serta variance.
\item
  Dari 20 email independen, hitung peluang tepat 5 terbuka.
\item
  Jika peluang sukses 0.25, hitung peluang sukses pertama terjadi pada
  percobaan ke-4.
\item
  Jika rata-rata panggilan masuk 6 per jam, hitung peluang tepat 8
  panggilan masuk.
\end{enumerate}

\subsection{C. Python}\label{c.-python-1}

\begin{enumerate}
\def\labelenumi{\arabic{enumi}.}
\tightlist
\item
  Simulasikan 50.000 Bernoulli dengan \(p=0.4\), lalu bandingkan mean
  dan variance simulasi dengan teori.
\item
  Simulasikan Binomial(20, 0.1) dan bandingkan histogram dengan PMF
  teoritis.
\item
  Simulasikan Geometric(0.3) dan verifikasi mean-nya.
\item
  Simulasikan Poisson(5) dan bandingkan histogram dengan PMF teoritis.
\end{enumerate}

\subsection{D. Aplikatif}\label{d.-aplikatif-1}

\begin{enumerate}
\def\labelenumi{\arabic{enumi}.}
\tightlist
\item
  Modelkan masalah quality control sederhana dengan Binomial.
\item
  Modelkan kedatangan pelanggan per jam dengan Poisson.
\item
  Modelkan percobaan login sampai berhasil dengan Geometric.
\item
  Jelaskan distribusi mana yang Anda pilih dan mengapa.
\end{enumerate}

\section{3.17 Penutup Kecil}\label{penutup-kecil-2}

Kalau di bab sebelumnya kita baru belajar ``apa itu random variable'',
maka di bab ini kita mulai benar-benar melihat keluarga model yang bisa
dipakai untuk dunia nyata.

Di bab berikutnya, kita akan beralih ke \textbf{distribusi random
variable kontinu}. Di sana, fokus kita tidak lagi pada hitungan
kejadian, tetapi pada besaran seperti waktu, umur hidup, tinggi, jarak,
dan banyak pengukuran lain yang lebih nyaman dipandang kontinu.

\bookmarksetup{startatroot}

\chapter{Bab 4. Distribusi Random Variable
Kontinu}\label{bab-4.-distribusi-random-variable-kontinu}

\section{Tujuan Bab}\label{tujuan-bab-3}

Setelah mempelajari bab ini, mahasiswa diharapkan mampu:

\begin{enumerate}
\def\labelenumi{\arabic{enumi}.}
\tightlist
\item
  memahami perbedaan mendasar antara random variable diskrit dan
  kontinu,
\item
  membangun distribusi kontinu sederhana dari tabel selang atau
  pendekatan trapezoid,
\item
  mengenali dan menggunakan distribusi kontinu penting:

  \begin{itemize}
  \tightlist
  \item
    uniform kontinu,
  \item
    normal,
  \item
    gamma,
  \item
    eksponensial,
  \item
    erlang,
  \item
    weibull,
  \item
    pareto,
  \item
    chi-square,
  \end{itemize}
\item
  memahami hubungan antardistribusi,
\item
  menggunakan Python untuk simulasi, visualisasi PDF/CDF, dan
  perhitungan peluang,
\item
  menghubungkan distribusi kontinu dengan keputusan nyata di bidang
  teknik, layanan, bisnis, dan reliability.
\end{enumerate}

\section{Pembuka}\label{pembuka-3}

Di bab sebelumnya, kita banyak berurusan dengan besaran yang berupa
\textbf{hitungan}: - jumlah cacat, - jumlah pelanggan, - jumlah sukses,
- jumlah percobaan.

Sekarang kita beralih ke dunia yang sedikit berbeda: dunia
\textbf{pengukuran}.

Kita akan memodelkan besaran seperti: - waktu tunggu, - umur hidup
lampu, - tinggi badan, - error pengukuran, - besar permintaan, - lama
layanan, - dan waktu antar kedatangan.

Pada banyak situasi seperti ini, lebih masuk akal memodelkan nilai
sebagai \textbf{kontinu}. Nilainya tidak sekadar 0, 1, 2, 3, tetapi
berada pada rentang yang halus di garis bilangan.

\section{4.1 Quick Win: Mengapa PDF Bukan Peluang
Langsung?}\label{quick-win-mengapa-pdf-bukan-peluang-langsung}

Mari ambil contoh yang sederhana. Misalkan \(X\) adalah random variable
kontinu yang uniform pada interval {[}0, 1{]}. Kita tahu PDF-nya adalah:

\[
f(x)=1 \quad \text{untuk } 0 \le x \le 1
\]

Sekarang pertanyaannya: - berapa \(P(X=0.5)\)? - berapa
\(P(0.4 \le X \le 0.6)\)?

\begin{Shaded}
\begin{Highlighting}[]
\ImportTok{import}\NormalTok{ numpy }\ImportTok{as}\NormalTok{ np}
\ImportTok{import}\NormalTok{ matplotlib.pyplot }\ImportTok{as}\NormalTok{ plt}
\ImportTok{from}\NormalTok{ scipy }\ImportTok{import}\NormalTok{ stats}

\CommentTok{\#\# Uniform(0,1)}
\BuiltInTok{print}\NormalTok{(}\StringTok{"P(X=0.5) untuk model kontinu = 0"}\NormalTok{)}
\BuiltInTok{print}\NormalTok{(}\StringTok{"P(0.4 \textless{}= X \textless{}= 0.6) ="}\NormalTok{, }\FloatTok{0.6} \OperatorTok{{-}} \FloatTok{0.4}\NormalTok{)}
\end{Highlighting}
\end{Shaded}

\begin{verbatim}
P(X=0.5) untuk model kontinu = 0
P(0.4 <= X <= 0.6) = 0.19999999999999996
\end{verbatim}

\begin{Shaded}
\begin{Highlighting}[]
\NormalTok{x }\OperatorTok{=}\NormalTok{ np.linspace(}\OperatorTok{{-}}\FloatTok{0.2}\NormalTok{, }\FloatTok{1.2}\NormalTok{, }\DecValTok{400}\NormalTok{)}
\NormalTok{pdf }\OperatorTok{=}\NormalTok{ np.where((x }\OperatorTok{\textgreater{}=} \DecValTok{0}\NormalTok{) }\OperatorTok{\&}\NormalTok{ (x }\OperatorTok{\textless{}=} \DecValTok{1}\NormalTok{), }\FloatTok{1.0}\NormalTok{, }\FloatTok{0.0}\NormalTok{)}

\NormalTok{plt.figure(figsize}\OperatorTok{=}\NormalTok{(}\DecValTok{8}\NormalTok{,}\DecValTok{4}\NormalTok{))}
\NormalTok{plt.plot(x, pdf)}
\NormalTok{plt.fill\_between(x, }\DecValTok{0}\NormalTok{, pdf, where}\OperatorTok{=}\NormalTok{((x }\OperatorTok{\textgreater{}=} \FloatTok{0.4}\NormalTok{) }\OperatorTok{\&}\NormalTok{ (x }\OperatorTok{\textless{}=} \FloatTok{0.6}\NormalTok{)), alpha}\OperatorTok{=}\FloatTok{0.3}\NormalTok{)}
\NormalTok{plt.xlabel(}\StringTok{"x"}\NormalTok{)}
\NormalTok{plt.ylabel(}\StringTok{"f(x)"}\NormalTok{)}
\NormalTok{plt.title(}\StringTok{"Pada RV kontinu, peluang = luas area, bukan tinggi kurva"}\NormalTok{)}
\NormalTok{plt.grid(alpha}\OperatorTok{=}\FloatTok{0.3}\NormalTok{)}
\NormalTok{plt.show()}
\end{Highlighting}
\end{Shaded}

\pandocbounded{\includegraphics[keepaspectratio]{04-distribusi-kontinu_files/figure-pdf/cell-3-output-1.pdf}}

Dari \emph{quick win} ini kita memperoleh satu pelajaran besar:

\begin{quote}
Untuk random variable kontinu, peluang tidak datang dari tinggi titik,
tetapi dari \textbf{luas area di bawah kurva PDF}.
\end{quote}

\section{4.2 Distribusi Kontinu Kustom dari Tabel Selang /
Trapezoid}\label{distribusi-kontinu-kustom-dari-tabel-selang-trapezoid}

\subsection{K --- Konteks}\label{k-konteks-19}

Kadang kita tidak mulai dari distribusi terkenal seperti Normal atau
Exponential. Kita hanya punya perkiraan bentuk dari domain expert atau
dari tabel interval. Misalnya: - permintaan produk kemungkinan besar di
tengah, menurun di tepi, - usia layanan komponen diperkirakan lebih
padat pada interval tertentu, - waktu tempuh berada dalam rentang
tertentu dengan perubahan kerapatan yang kasar.

\subsection{M --- Model}\label{m-model-19}

Kita dapat membentuk PDF kustom dari titik-titik \((x_i, f_i)\), lalu
menggunakan interpolasi linier atau pendekatan trapezoid. Selama luas
totalnya dinormalisasi menjadi 1, fungsi itu dapat dipakai sebagai PDF.

\subsection{Q --- Questions}\label{q-questions-19}

\begin{enumerate}
\def\labelenumi{\arabic{enumi}.}
\tightlist
\item
  Bagaimana membangun PDF kontinu dari tabel?
\item
  Bagaimana memastikan luas total = 1?
\item
  Bagaimana menghitung peluang interval secara numerik?
\end{enumerate}

\subsection{A --- Apply}\label{a-apply-19}

\begin{Shaded}
\begin{Highlighting}[]
\CommentTok{\#\# Titik{-}titik pembentuk "density" mentah}
\NormalTok{x\_pts }\OperatorTok{=}\NormalTok{ np.array([}\DecValTok{0}\NormalTok{, }\DecValTok{2}\NormalTok{, }\DecValTok{4}\NormalTok{, }\DecValTok{6}\NormalTok{, }\DecValTok{8}\NormalTok{, }\DecValTok{10}\NormalTok{], dtype}\OperatorTok{=}\BuiltInTok{float}\NormalTok{)}
\NormalTok{y\_raw }\OperatorTok{=}\NormalTok{ np.array([}\FloatTok{0.0}\NormalTok{, }\FloatTok{0.1}\NormalTok{, }\FloatTok{0.25}\NormalTok{, }\FloatTok{0.2}\NormalTok{, }\FloatTok{0.08}\NormalTok{, }\FloatTok{0.0}\NormalTok{], dtype}\OperatorTok{=}\BuiltInTok{float}\NormalTok{)}

\CommentTok{\#\# Normalisasi agar integral = 1}
\NormalTok{area\_raw }\OperatorTok{=}\NormalTok{ np.trapezoid(y\_raw, x\_pts)}
\NormalTok{y }\OperatorTok{=}\NormalTok{ y\_raw }\OperatorTok{/}\NormalTok{ area\_raw}

\NormalTok{x\_grid }\OperatorTok{=}\NormalTok{ np.linspace(}\DecValTok{0}\NormalTok{, }\DecValTok{10}\NormalTok{, }\DecValTok{500}\NormalTok{)}
\NormalTok{y\_grid }\OperatorTok{=}\NormalTok{ np.interp(x\_grid, x\_pts, y)}
\end{Highlighting}
\end{Shaded}

\begin{Shaded}
\begin{Highlighting}[]
\NormalTok{plt.figure(figsize}\OperatorTok{=}\NormalTok{(}\DecValTok{8}\NormalTok{,}\DecValTok{4}\NormalTok{))}
\NormalTok{plt.plot(x\_grid, y\_grid, label}\OperatorTok{=}\StringTok{"PDF kustom (interpolasi linier)"}\NormalTok{)}
\NormalTok{plt.scatter(x\_pts, y, color}\OperatorTok{=}\StringTok{\textquotesingle{}red\textquotesingle{}}\NormalTok{)}
\NormalTok{plt.xlabel(}\StringTok{"x"}\NormalTok{)}
\NormalTok{plt.ylabel(}\StringTok{"f(x)"}\NormalTok{)}
\NormalTok{plt.title(}\StringTok{"PDF kontinu kustom dari tabel selang"}\NormalTok{)}
\NormalTok{plt.legend()}
\NormalTok{plt.grid(alpha}\OperatorTok{=}\FloatTok{0.3}\NormalTok{)}
\NormalTok{plt.show()}
\end{Highlighting}
\end{Shaded}

\pandocbounded{\includegraphics[keepaspectratio]{04-distribusi-kontinu_files/figure-pdf/cell-5-output-1.pdf}}

\begin{Shaded}
\begin{Highlighting}[]
\CommentTok{\#\# Peluang interval [3, 7]}
\NormalTok{mask }\OperatorTok{=}\NormalTok{ (x\_grid }\OperatorTok{\textgreater{}=} \DecValTok{3}\NormalTok{) }\OperatorTok{\&}\NormalTok{ (x\_grid }\OperatorTok{\textless{}=} \DecValTok{7}\NormalTok{)}
\NormalTok{prob\_3\_7 }\OperatorTok{=}\NormalTok{ np.trapezoid(y\_grid[mask], x\_grid[mask])}
\BuiltInTok{print}\NormalTok{(}\StringTok{"P(3 \textless{}= X \textless{}= 7) \textasciitilde{}"}\NormalTok{, prob\_3\_7)}

\CommentTok{\#\# Mean numerik}
\NormalTok{mean\_num }\OperatorTok{=}\NormalTok{ np.trapezoid(x\_grid }\OperatorTok{*}\NormalTok{ y\_grid, x\_grid)}
\NormalTok{var\_num }\OperatorTok{=}\NormalTok{ np.trapezoid((x\_grid }\OperatorTok{{-}}\NormalTok{ mean\_num)}\OperatorTok{**}\DecValTok{2} \OperatorTok{*}\NormalTok{ y\_grid, x\_grid)}
\BuiltInTok{print}\NormalTok{(}\StringTok{"Mean \textasciitilde{}"}\NormalTok{, mean\_num)}
\BuiltInTok{print}\NormalTok{(}\StringTok{"Variance \textasciitilde{}"}\NormalTok{, var\_num)}
\end{Highlighting}
\end{Shaded}

\begin{verbatim}
P(3 <= X <= 7) ~ 0.659204179902892
Mean ~ 4.825378851610921
Variance ~ 3.92183354068992
\end{verbatim}

\subsubsection{Interpretasi}\label{interpretasi-10}

Distribusi kontinu kustom penting ketika: - data atau pengetahuan domain
belum cukup untuk memilih distribusi tertutup tertentu, - kita tetap
ingin menghitung peluang, mean, variance, dan quantile, - pendekatan
numerik lebih realistis daripada memaksakan model yang belum tentu
cocok.

\section{4.3 Uniform Kontinu}\label{uniform-kontinu}

\subsection{K --- Konteks}\label{k-konteks-20}

Distribusi uniform kontinu cocok ketika nilai dianggap tersebar merata
pada sebuah interval. Misalnya: - waktu kedatangan acak di antara dua
titik waktu, - lokasi acak pada segmen garis, - nilai acak yang dipilih
secara merata dari suatu rentang.

\subsection{M --- Model}\label{m-model-20}

Jika: \[
X \sim \mathrm{Uniform}(a,b)
\] maka PDF-nya:

\[
f(x)=\frac{1}{b-a}, \qquad a \le x \le b
\]

dan nol di luar interval tersebut.

\subsection{Q --- Questions}\label{q-questions-20}

\begin{enumerate}
\def\labelenumi{\arabic{enumi}.}
\tightlist
\item
  Berapa peluang \(X\) berada pada sub-interval tertentu?
\item
  Berapa mean dan variance?
\item
  Kapan asumsi ``semua nilai di interval sama padat'' masuk akal?
\end{enumerate}

\subsection{A --- Apply}\label{a-apply-20}

\begin{Shaded}
\begin{Highlighting}[]
\NormalTok{a, b }\OperatorTok{=} \DecValTok{2}\NormalTok{, }\DecValTok{8}
\NormalTok{x }\OperatorTok{=}\NormalTok{ np.linspace(}\DecValTok{0}\NormalTok{, }\DecValTok{10}\NormalTok{, }\DecValTok{400}\NormalTok{)}
\NormalTok{pdf }\OperatorTok{=}\NormalTok{ np.where((x }\OperatorTok{\textgreater{}=}\NormalTok{ a) }\OperatorTok{\&}\NormalTok{ (x }\OperatorTok{\textless{}=}\NormalTok{ b), }\DecValTok{1}\OperatorTok{/}\NormalTok{(b}\OperatorTok{{-}}\NormalTok{a), }\FloatTok{0.0}\NormalTok{)}

\NormalTok{plt.figure(figsize}\OperatorTok{=}\NormalTok{(}\DecValTok{8}\NormalTok{,}\DecValTok{4}\NormalTok{))}
\NormalTok{plt.plot(x, pdf)}
\NormalTok{plt.fill\_between(x, }\DecValTok{0}\NormalTok{, pdf, where}\OperatorTok{=}\NormalTok{((x }\OperatorTok{\textgreater{}=} \DecValTok{3}\NormalTok{) }\OperatorTok{\&}\NormalTok{ (x }\OperatorTok{\textless{}=} \DecValTok{5}\NormalTok{)), alpha}\OperatorTok{=}\FloatTok{0.3}\NormalTok{)}
\NormalTok{plt.xlabel(}\StringTok{"x"}\NormalTok{)}
\NormalTok{plt.ylabel(}\StringTok{"f(x)"}\NormalTok{)}
\NormalTok{plt.title(}\StringTok{"PDF Uniform(2,8)"}\NormalTok{)}
\NormalTok{plt.grid(alpha}\OperatorTok{=}\FloatTok{0.3}\NormalTok{)}
\NormalTok{plt.show()}
\end{Highlighting}
\end{Shaded}

\pandocbounded{\includegraphics[keepaspectratio]{04-distribusi-kontinu_files/figure-pdf/cell-7-output-1.pdf}}

\begin{Shaded}
\begin{Highlighting}[]
\BuiltInTok{print}\NormalTok{(}\StringTok{"P(3 \textless{}= X \textless{}= 5) ="}\NormalTok{, (}\DecValTok{5}\OperatorTok{{-}}\DecValTok{3}\NormalTok{)}\OperatorTok{/}\NormalTok{(}\DecValTok{8}\OperatorTok{{-}}\DecValTok{2}\NormalTok{))}
\BuiltInTok{print}\NormalTok{(}\StringTok{"Mean ="}\NormalTok{, (a}\OperatorTok{+}\NormalTok{b)}\OperatorTok{/}\DecValTok{2}\NormalTok{)}
\BuiltInTok{print}\NormalTok{(}\StringTok{"Variance ="}\NormalTok{, (b}\OperatorTok{{-}}\NormalTok{a)}\OperatorTok{**}\DecValTok{2} \OperatorTok{/} \DecValTok{12}\NormalTok{)}
\end{Highlighting}
\end{Shaded}

\begin{verbatim}
P(3 <= X <= 5) = 0.3333333333333333
Mean = 5.0
Variance = 3.0
\end{verbatim}

\subsubsection{Hasil teoritis}\label{hasil-teoritis-5}

\[
E[X]=\frac{a+b}{2}
\]

\[
\mathrm{Var}(X)=\frac{(b-a)^2}{12}
\]

\subsection{Latihan singkat}\label{latihan-singkat-11}

\begin{enumerate}
\def\labelenumi{\arabic{enumi}.}
\tightlist
\item
  Jika \(X \sim \mathrm{Uniform}(10,20)\), hitung
  \(P(12 \le X \le 15)\).
\item
  Hitung mean dan variance-nya.
\item
  Beri satu contoh aplikasi nyata yang cocok dimodelkan uniform kontinu.
\end{enumerate}

\section{4.4 Normal Distribution}\label{normal-distribution}

\subsection{K --- Konteks}\label{k-konteks-21}

Normal adalah salah satu distribusi paling terkenal karena sering muncul
pada: - error pengukuran, - tinggi badan, - berat badan, - hasil
agregasi banyak pengaruh kecil, - noise sensor.

\subsection{M --- Model}\label{m-model-21}

Jika: \[
X \sim \mathcal{N}(\mu, \sigma^2)
\] maka PDF-nya:

\[
f(x)=\frac{1}{\sigma \sqrt{2\pi}}\exp\left(-\frac{(x-\mu)^2}{2\sigma^2}\right)
\]

Parameter: - \(\mu\): pusat distribusi, - \(\sigma\): skala penyebaran.

\subsection{Q --- Questions}\label{q-questions-21}

\begin{enumerate}
\def\labelenumi{\arabic{enumi}.}
\tightlist
\item
  Berapa peluang \(X\) kurang dari nilai tertentu?
\item
  Berapa peluang \(X\) berada dalam interval tertentu?
\item
  Bagaimana pengaruh \(\mu\) dan \(\sigma\) pada bentuk kurva?
\end{enumerate}

\subsection{A --- Apply}\label{a-apply-21}

\begin{Shaded}
\begin{Highlighting}[]
\NormalTok{mu, sigma }\OperatorTok{=} \DecValTok{100}\NormalTok{, }\DecValTok{15}
\NormalTok{x }\OperatorTok{=}\NormalTok{ np.linspace(}\DecValTok{40}\NormalTok{, }\DecValTok{160}\NormalTok{, }\DecValTok{500}\NormalTok{)}
\NormalTok{pdf }\OperatorTok{=}\NormalTok{ stats.norm.pdf(x, loc}\OperatorTok{=}\NormalTok{mu, scale}\OperatorTok{=}\NormalTok{sigma)}

\NormalTok{plt.figure(figsize}\OperatorTok{=}\NormalTok{(}\DecValTok{8}\NormalTok{,}\DecValTok{4}\NormalTok{))}
\NormalTok{plt.plot(x, pdf)}
\NormalTok{plt.xlabel(}\StringTok{"x"}\NormalTok{)}
\NormalTok{plt.ylabel(}\StringTok{"f(x)"}\NormalTok{)}
\NormalTok{plt.title(}\StringTok{"PDF Normal(100, 15\^{}2)"}\NormalTok{)}
\NormalTok{plt.grid(alpha}\OperatorTok{=}\FloatTok{0.3}\NormalTok{)}
\NormalTok{plt.show()}
\end{Highlighting}
\end{Shaded}

\pandocbounded{\includegraphics[keepaspectratio]{04-distribusi-kontinu_files/figure-pdf/cell-9-output-1.pdf}}

\begin{Shaded}
\begin{Highlighting}[]
\BuiltInTok{print}\NormalTok{(}\StringTok{"P(X \textless{} 90) ="}\NormalTok{, stats.norm.cdf(}\DecValTok{90}\NormalTok{, loc}\OperatorTok{=}\NormalTok{mu, scale}\OperatorTok{=}\NormalTok{sigma))}
\BuiltInTok{print}\NormalTok{(}\StringTok{"P(90 \textless{}= X \textless{}= 110) ="}\NormalTok{, stats.norm.cdf(}\DecValTok{110}\NormalTok{, loc}\OperatorTok{=}\NormalTok{mu, scale}\OperatorTok{=}\NormalTok{sigma) }\OperatorTok{{-}}\NormalTok{ stats.norm.cdf(}\DecValTok{90}\NormalTok{, loc}\OperatorTok{=}\NormalTok{mu, scale}\OperatorTok{=}\NormalTok{sigma))}
\end{Highlighting}
\end{Shaded}

\begin{verbatim}
P(X < 90) = 0.2524925375469229
P(90 <= X <= 110) = 0.4950149249061542
\end{verbatim}

\subsubsection{Hubungan dengan Binomial}\label{hubungan-dengan-binomial}

Jika: - \(X \sim \mathrm{Binomial}(n,p)\), - \(n\) besar,

maka Binomial dapat didekati oleh Normal dengan:

\[
\mu = np, \qquad \sigma^2 = np(1-p)
\]

Mari lihat contoh.

\begin{Shaded}
\begin{Highlighting}[]
\NormalTok{n, p }\OperatorTok{=} \DecValTok{100}\NormalTok{, }\FloatTok{0.4}
\NormalTok{mu\_bin }\OperatorTok{=}\NormalTok{ n}\OperatorTok{*}\NormalTok{p}
\NormalTok{sigma\_bin }\OperatorTok{=}\NormalTok{ np.sqrt(n}\OperatorTok{*}\NormalTok{p}\OperatorTok{*}\NormalTok{(}\DecValTok{1}\OperatorTok{{-}}\NormalTok{p))}

\NormalTok{x }\OperatorTok{=}\NormalTok{ np.arange(}\DecValTok{20}\NormalTok{, }\DecValTok{61}\NormalTok{)}
\NormalTok{pmf\_bin }\OperatorTok{=}\NormalTok{ stats.binom.pmf(x, n}\OperatorTok{=}\NormalTok{n, p}\OperatorTok{=}\NormalTok{p)}

\NormalTok{x\_cont }\OperatorTok{=}\NormalTok{ np.linspace(}\DecValTok{20}\NormalTok{, }\DecValTok{60}\NormalTok{, }\DecValTok{500}\NormalTok{)}
\NormalTok{pdf\_norm }\OperatorTok{=}\NormalTok{ stats.norm.pdf(x\_cont, loc}\OperatorTok{=}\NormalTok{mu\_bin, scale}\OperatorTok{=}\NormalTok{sigma\_bin)}

\NormalTok{plt.figure(figsize}\OperatorTok{=}\NormalTok{(}\DecValTok{8}\NormalTok{,}\DecValTok{4}\NormalTok{))}
\NormalTok{plt.stem(x, pmf\_bin, basefmt}\OperatorTok{=}\StringTok{" "}\NormalTok{, label}\OperatorTok{=}\StringTok{"Binomial"}\NormalTok{)}
\NormalTok{plt.plot(x\_cont, pdf\_norm, }\StringTok{\textquotesingle{}r\textquotesingle{}}\NormalTok{, label}\OperatorTok{=}\StringTok{"Normal approximation"}\NormalTok{)}
\NormalTok{plt.xlabel(}\StringTok{"x"}\NormalTok{)}
\NormalTok{plt.ylabel(}\StringTok{"Probability / Density"}\NormalTok{)}
\NormalTok{plt.title(}\StringTok{"Pendekatan Normal terhadap Binomial"}\NormalTok{)}
\NormalTok{plt.legend()}
\NormalTok{plt.grid(alpha}\OperatorTok{=}\FloatTok{0.3}\NormalTok{)}
\NormalTok{plt.show()}
\end{Highlighting}
\end{Shaded}

\pandocbounded{\includegraphics[keepaspectratio]{04-distribusi-kontinu_files/figure-pdf/cell-11-output-1.pdf}}

\subsubsection{Interpretasi keputusan}\label{interpretasi-keputusan-6}

Normal sering dipakai untuk: - menetapkan toleransi manufaktur, -
menghitung peluang deviasi pengukuran, - mengevaluasi performa relatif
terhadap target.

\subsection{Latihan singkat}\label{latihan-singkat-12}

\begin{enumerate}
\def\labelenumi{\arabic{enumi}.}
\tightlist
\item
  Jika \(X \sim \mathcal{N}(70,10^2)\), hitung \(P(X < 60)\).
\item
  Hitung \(P(65 \le X \le 80)\).
\item
  Simulasikan 10000 sampel dari distribusi ini dan bandingkan histogram
  dengan PDF.
\end{enumerate}

\section{4.5 Gamma Distribution}\label{gamma-distribution}

\subsection{K --- Konteks}\label{k-konteks-22}

Distribusi Gamma sangat berguna untuk memodelkan: - waktu tunggu sampai
beberapa kejadian terjadi, - lifetime positif dengan bentuk fleksibel, -
akumulasi beberapa waktu tunggu eksponensial.

\subsection{M --- Model}\label{m-model-22}

Salah satu parametrisasinya: \[
X \sim \mathrm{Gamma}(\alpha, \theta)
\]

dengan: - \(\alpha\) = shape, - \(\theta\) = scale.

PDF-nya:

\[
f(x)=\frac{1}{\Gamma(\alpha)\theta^\alpha}x^{\alpha-1}e^{-x/\theta}, \qquad x>0
\]

\subsection{Q --- Questions}\label{q-questions-22}

\begin{enumerate}
\def\labelenumi{\arabic{enumi}.}
\tightlist
\item
  Apa pengaruh shape dan scale?
\item
  Kapan Gamma lebih cocok daripada Exponential?
\item
  Bagaimana menghitung peluang waktu selesai sebelum batas tertentu?
\end{enumerate}

\subsection{A --- Apply}\label{a-apply-22}

\begin{Shaded}
\begin{Highlighting}[]
\NormalTok{alpha, theta }\OperatorTok{=} \DecValTok{3}\NormalTok{, }\DecValTok{2}
\NormalTok{x }\OperatorTok{=}\NormalTok{ np.linspace(}\DecValTok{0}\NormalTok{, }\DecValTok{30}\NormalTok{, }\DecValTok{500}\NormalTok{)}
\NormalTok{pdf }\OperatorTok{=}\NormalTok{ stats.gamma.pdf(x, a}\OperatorTok{=}\NormalTok{alpha, scale}\OperatorTok{=}\NormalTok{theta)}

\NormalTok{plt.figure(figsize}\OperatorTok{=}\NormalTok{(}\DecValTok{8}\NormalTok{,}\DecValTok{4}\NormalTok{))}
\NormalTok{plt.plot(x, pdf)}
\NormalTok{plt.xlabel(}\StringTok{"x"}\NormalTok{)}
\NormalTok{plt.ylabel(}\StringTok{"f(x)"}\NormalTok{)}
\NormalTok{plt.title(}\StringTok{"PDF Gamma(shape=3, scale=2)"}\NormalTok{)}
\NormalTok{plt.grid(alpha}\OperatorTok{=}\FloatTok{0.3}\NormalTok{)}
\NormalTok{plt.show()}
\end{Highlighting}
\end{Shaded}

\pandocbounded{\includegraphics[keepaspectratio]{04-distribusi-kontinu_files/figure-pdf/cell-12-output-1.pdf}}

\begin{Shaded}
\begin{Highlighting}[]
\BuiltInTok{print}\NormalTok{(}\StringTok{"Mean ="}\NormalTok{, alpha }\OperatorTok{*}\NormalTok{ theta)}
\BuiltInTok{print}\NormalTok{(}\StringTok{"Variance ="}\NormalTok{, alpha }\OperatorTok{*}\NormalTok{ theta}\OperatorTok{**}\DecValTok{2}\NormalTok{)}
\BuiltInTok{print}\NormalTok{(}\StringTok{"P(X \textless{}= 8) ="}\NormalTok{, stats.gamma.cdf(}\DecValTok{8}\NormalTok{, a}\OperatorTok{=}\NormalTok{alpha, scale}\OperatorTok{=}\NormalTok{theta))}
\end{Highlighting}
\end{Shaded}

\begin{verbatim}
Mean = 6
Variance = 12
P(X <= 8) = 0.7618966944464556
\end{verbatim}

\subsubsection{Hasil teoritis}\label{hasil-teoritis-6}

\[
E[X]=\alpha\theta
\]

\[
\mathrm{Var}(X)=\alpha\theta^2
\]

\subsection{Interpretasi}\label{interpretasi-11}

Gamma cocok untuk waktu positif yang tidak memoryless dan memiliki
bentuk yang bisa semakin ``membukit'' saat shape membesar.

\section{4.6 Exponential Distribution}\label{exponential-distribution}

\subsection{K --- Konteks}\label{k-konteks-23}

Distribusi Exponential sangat terkenal untuk memodelkan: - waktu antar
kedatangan pada Poisson process, - waktu sampai kegagalan pertama pada
sistem memoryless, - waktu tunggu hingga kejadian berikutnya.

\subsection{M --- Model}\label{m-model-23}

Jika: \[
X \sim \mathrm{Exponential}(\lambda)
\] maka PDF-nya:

\[
f(x)=\lambda e^{-\lambda x}, \qquad x \ge 0
\]

CDF-nya:

\[
F(x)=1-e^{-\lambda x}
\]

\subsection{Q --- Questions}\label{q-questions-23}

\begin{enumerate}
\def\labelenumi{\arabic{enumi}.}
\tightlist
\item
  Berapa peluang menunggu kurang dari \(t\)?
\item
  Mengapa distribusi ini memoryless?
\item
  Bagaimana hubungannya dengan Poisson process?
\end{enumerate}

\subsection{A --- Apply}\label{a-apply-23}

\begin{Shaded}
\begin{Highlighting}[]
\NormalTok{lam }\OperatorTok{=} \FloatTok{0.5}
\NormalTok{x }\OperatorTok{=}\NormalTok{ np.linspace(}\DecValTok{0}\NormalTok{, }\DecValTok{15}\NormalTok{, }\DecValTok{500}\NormalTok{)}
\NormalTok{pdf }\OperatorTok{=}\NormalTok{ stats.expon.pdf(x, scale}\OperatorTok{=}\DecValTok{1}\OperatorTok{/}\NormalTok{lam)}

\NormalTok{plt.figure(figsize}\OperatorTok{=}\NormalTok{(}\DecValTok{8}\NormalTok{,}\DecValTok{4}\NormalTok{))}
\NormalTok{plt.plot(x, pdf)}
\NormalTok{plt.xlabel(}\StringTok{"x"}\NormalTok{)}
\NormalTok{plt.ylabel(}\StringTok{"f(x)"}\NormalTok{)}
\NormalTok{plt.title(}\StringTok{"PDF Exponential(rate=0.5)"}\NormalTok{)}
\NormalTok{plt.grid(alpha}\OperatorTok{=}\FloatTok{0.3}\NormalTok{)}
\NormalTok{plt.show()}
\end{Highlighting}
\end{Shaded}

\pandocbounded{\includegraphics[keepaspectratio]{04-distribusi-kontinu_files/figure-pdf/cell-14-output-1.pdf}}

\begin{Shaded}
\begin{Highlighting}[]
\BuiltInTok{print}\NormalTok{(}\StringTok{"Mean ="}\NormalTok{, }\DecValTok{1}\OperatorTok{/}\NormalTok{lam)}
\BuiltInTok{print}\NormalTok{(}\StringTok{"Variance ="}\NormalTok{, }\DecValTok{1}\OperatorTok{/}\NormalTok{lam}\OperatorTok{**}\DecValTok{2}\NormalTok{)}
\BuiltInTok{print}\NormalTok{(}\StringTok{"P(X \textless{}= 3) ="}\NormalTok{, stats.expon.cdf(}\DecValTok{3}\NormalTok{, scale}\OperatorTok{=}\DecValTok{1}\OperatorTok{/}\NormalTok{lam))}
\BuiltInTok{print}\NormalTok{(}\StringTok{"P(X \textgreater{} 5) ="}\NormalTok{, }\DecValTok{1} \OperatorTok{{-}}\NormalTok{ stats.expon.cdf(}\DecValTok{5}\NormalTok{, scale}\OperatorTok{=}\DecValTok{1}\OperatorTok{/}\NormalTok{lam))}
\end{Highlighting}
\end{Shaded}

\begin{verbatim}
Mean = 2.0
Variance = 4.0
P(X <= 3) = 0.7768698398515702
P(X > 5) = 0.08208499862389884
\end{verbatim}

\subsubsection{Sifat memoryless}\label{sifat-memoryless-1}

\[
P(X>s+t \mid X>s)=P(X>t)
\]

Distribusi Exponential adalah versi kontinu dari ide memoryless yang di
dunia diskrit dimiliki distribusi Geometric.

\subsubsection{Hubungan dengan Poisson
process}\label{hubungan-dengan-poisson-process}

Jika jumlah kejadian per interval mengikuti Poisson process dengan rate
\(\lambda\), maka waktu antar kedatangan mengikuti
Exponential(\(\lambda\)).

Mari simulasikan.

\begin{Shaded}
\begin{Highlighting}[]
\NormalTok{rng }\OperatorTok{=}\NormalTok{ np.random.default\_rng(}\DecValTok{123}\NormalTok{)}
\NormalTok{lam }\OperatorTok{=} \FloatTok{2.0}  \CommentTok{\#\# rata{-}rata 2 kejadian per satuan waktu}

\NormalTok{interarrival }\OperatorTok{=}\NormalTok{ rng.exponential(scale}\OperatorTok{=}\DecValTok{1}\OperatorTok{/}\NormalTok{lam, size}\OperatorTok{=}\DecValTok{50000}\NormalTok{)}

\NormalTok{x }\OperatorTok{=}\NormalTok{ np.linspace(}\DecValTok{0}\NormalTok{, np.percentile(interarrival, }\FloatTok{99.5}\NormalTok{), }\DecValTok{400}\NormalTok{)}
\NormalTok{pdf }\OperatorTok{=}\NormalTok{ stats.expon.pdf(x, scale}\OperatorTok{=}\DecValTok{1}\OperatorTok{/}\NormalTok{lam)}

\NormalTok{plt.figure(figsize}\OperatorTok{=}\NormalTok{(}\DecValTok{8}\NormalTok{,}\DecValTok{4}\NormalTok{))}
\NormalTok{plt.hist(interarrival, bins}\OperatorTok{=}\DecValTok{60}\NormalTok{, density}\OperatorTok{=}\VariableTok{True}\NormalTok{, alpha}\OperatorTok{=}\FloatTok{0.6}\NormalTok{, label}\OperatorTok{=}\StringTok{"Simulasi"}\NormalTok{)}
\NormalTok{plt.plot(x, pdf, label}\OperatorTok{=}\StringTok{"PDF teoritis"}\NormalTok{)}
\NormalTok{plt.xlabel(}\StringTok{"Interarrival time"}\NormalTok{)}
\NormalTok{plt.ylabel(}\StringTok{"Density"}\NormalTok{)}
\NormalTok{plt.title(}\StringTok{"Interarrival time pada Poisson process \textasciitilde{} Exponential"}\NormalTok{)}
\NormalTok{plt.legend()}
\NormalTok{plt.grid(alpha}\OperatorTok{=}\FloatTok{0.3}\NormalTok{)}
\NormalTok{plt.show()}
\end{Highlighting}
\end{Shaded}

\pandocbounded{\includegraphics[keepaspectratio]{04-distribusi-kontinu_files/figure-pdf/cell-16-output-1.pdf}}

\subsection{Latihan singkat}\label{latihan-singkat-13}

\begin{enumerate}
\def\labelenumi{\arabic{enumi}.}
\tightlist
\item
  Jika rata-rata waktu tunggu 10 menit, berapa rate \(\lambda\)?
\item
  Hitung peluang waktu tunggu kurang dari 5 menit.
\item
  Verifikasi sifat memoryless dengan simulasi.
\end{enumerate}

\section{4.7 Erlang Distribution}\label{erlang-distribution}

\subsection{K --- Konteks}\label{k-konteks-24}

Erlang muncul ketika kita menjumlahkan beberapa waktu tunggu
Eksponensial independen dengan rate yang sama. Misalnya: - total waktu
layanan yang melewati beberapa tahap berurutan, - waktu sampai \(k\)
kejadian dalam Poisson process.

\subsection{M --- Model}\label{m-model-24}

Jika: \[
X = X_1 + X_2 + \cdots + X_k
\] dengan tiap \(X_i \sim \text{Exponential}(\lambda)\) independen,
maka:

\[
X \sim \text{Erlang}(k,\lambda)
\]

Erlang adalah kasus khusus Gamma dengan shape integer.

\subsection{Q --- Questions}\label{q-questions-24}

\begin{enumerate}
\def\labelenumi{\arabic{enumi}.}
\tightlist
\item
  Bagaimana bentuk distribusinya dibanding Exponential?
\item
  Bagaimana pengaruh \(k\)?
\item
  Apa arti praktis ``jumlah beberapa tahap layanan''?
\end{enumerate}

\subsection{A --- Apply}\label{a-apply-24}

\begin{Shaded}
\begin{Highlighting}[]
\NormalTok{lam }\OperatorTok{=} \FloatTok{0.5}
\NormalTok{k }\OperatorTok{=} \DecValTok{3}

\NormalTok{x }\OperatorTok{=}\NormalTok{ np.linspace(}\DecValTok{0}\NormalTok{, }\DecValTok{25}\NormalTok{, }\DecValTok{500}\NormalTok{)}
\NormalTok{pdf }\OperatorTok{=}\NormalTok{ stats.gamma.pdf(x, a}\OperatorTok{=}\NormalTok{k, scale}\OperatorTok{=}\DecValTok{1}\OperatorTok{/}\NormalTok{lam)}

\NormalTok{plt.figure(figsize}\OperatorTok{=}\NormalTok{(}\DecValTok{8}\NormalTok{,}\DecValTok{4}\NormalTok{))}
\NormalTok{plt.plot(x, pdf)}
\NormalTok{plt.xlabel(}\StringTok{"x"}\NormalTok{)}
\NormalTok{plt.ylabel(}\StringTok{"f(x)"}\NormalTok{)}
\NormalTok{plt.title(}\StringTok{"PDF Erlang(k=3, rate=0.5)"}\NormalTok{)}
\NormalTok{plt.grid(alpha}\OperatorTok{=}\FloatTok{0.3}\NormalTok{)}
\NormalTok{plt.show()}
\end{Highlighting}
\end{Shaded}

\pandocbounded{\includegraphics[keepaspectratio]{04-distribusi-kontinu_files/figure-pdf/cell-17-output-1.pdf}}

\begin{Shaded}
\begin{Highlighting}[]
\BuiltInTok{print}\NormalTok{(}\StringTok{"Mean ="}\NormalTok{, k}\OperatorTok{/}\NormalTok{lam)}
\BuiltInTok{print}\NormalTok{(}\StringTok{"Variance ="}\NormalTok{, k}\OperatorTok{/}\NormalTok{(lam}\OperatorTok{**}\DecValTok{2}\NormalTok{))}
\BuiltInTok{print}\NormalTok{(}\StringTok{"P(X \textless{}= 8) ="}\NormalTok{, stats.gamma.cdf(}\DecValTok{8}\NormalTok{, a}\OperatorTok{=}\NormalTok{k, scale}\OperatorTok{=}\DecValTok{1}\OperatorTok{/}\NormalTok{lam))}
\end{Highlighting}
\end{Shaded}

\begin{verbatim}
Mean = 6.0
Variance = 12.0
P(X <= 8) = 0.7618966944464556
\end{verbatim}

\subsubsection{Interpretasi}\label{interpretasi-12}

Jika suatu proses terdiri dari 3 tahap memoryless independen, maka total
waktunya tidak lagi Exponential, tetapi Erlang-3.

\section{4.8 Weibull Distribution}\label{weibull-distribution}

\subsection{K --- Konteks}\label{k-konteks-25}

Weibull sangat populer dalam reliability engineering karena dapat
memodelkan berbagai pola kegagalan: - kegagalan awal, - kegagalan acak,
- kegagalan akibat keausan.

\subsection{M --- Model}\label{m-model-25}

Jika: \[
X \sim \text{Weibull}(k,\lambda)
\] dengan: - \(k\) = shape, - \(\lambda\) = scale,

maka PDF-nya adalah:

\[
f(x)=\frac{k}{\lambda}\left(\frac{x}{\lambda}\right)^{k-1}e^{-(x/\lambda)^k}, \qquad x \ge 0
\]

\subsection{Q --- Questions}\label{q-questions-25}

\begin{enumerate}
\def\labelenumi{\arabic{enumi}.}
\tightlist
\item
  Bagaimana shape \(k\) memengaruhi hazard?
\item
  Kapan Weibull lebih realistis daripada Exponential?
\item
  Bagaimana peluang komponen gagal sebelum waktu tertentu?
\end{enumerate}

\subsection{A --- Apply}\label{a-apply-25}

\begin{Shaded}
\begin{Highlighting}[]
\NormalTok{k\_shape }\OperatorTok{=} \FloatTok{2.0}
\NormalTok{lam\_scale }\OperatorTok{=} \FloatTok{10.0}

\NormalTok{x }\OperatorTok{=}\NormalTok{ np.linspace(}\DecValTok{0}\NormalTok{, }\DecValTok{30}\NormalTok{, }\DecValTok{500}\NormalTok{)}
\NormalTok{pdf }\OperatorTok{=}\NormalTok{ stats.weibull\_min.pdf(x, c}\OperatorTok{=}\NormalTok{k\_shape, scale}\OperatorTok{=}\NormalTok{lam\_scale)}

\NormalTok{plt.figure(figsize}\OperatorTok{=}\NormalTok{(}\DecValTok{8}\NormalTok{,}\DecValTok{4}\NormalTok{))}
\NormalTok{plt.plot(x, pdf)}
\NormalTok{plt.xlabel(}\StringTok{"x"}\NormalTok{)}
\NormalTok{plt.ylabel(}\StringTok{"f(x)"}\NormalTok{)}
\NormalTok{plt.title(}\StringTok{"PDF Weibull(shape=2, scale=10)"}\NormalTok{)}
\NormalTok{plt.grid(alpha}\OperatorTok{=}\FloatTok{0.3}\NormalTok{)}
\NormalTok{plt.show()}
\end{Highlighting}
\end{Shaded}

\pandocbounded{\includegraphics[keepaspectratio]{04-distribusi-kontinu_files/figure-pdf/cell-19-output-1.pdf}}

\begin{Shaded}
\begin{Highlighting}[]
\BuiltInTok{print}\NormalTok{(}\StringTok{"P(X \textless{}= 8) ="}\NormalTok{, stats.weibull\_min.cdf(}\DecValTok{8}\NormalTok{, c}\OperatorTok{=}\NormalTok{k\_shape, scale}\OperatorTok{=}\NormalTok{lam\_scale))}
\end{Highlighting}
\end{Shaded}

\begin{verbatim}
P(X <= 8) = 0.47270757595695145
\end{verbatim}

\subsubsection{Interpretasi}\label{interpretasi-13}

Weibull lebih fleksibel daripada Exponential. Ia sangat berguna saat
laju kegagalan tidak konstan.

\section{4.9 Pareto Distribution}\label{pareto-distribution}

\subsection{K --- Konteks}\label{k-konteks-26}

Pareto dipakai untuk fenomena heavy-tail, misalnya: - kekayaan, - ukuran
file, - kerugian ekstrem, - pendapatan segelintir pihak yang sangat
besar.

\subsection{M --- Model}\label{m-model-26}

Jika: \[
X \sim \text{Pareto}(x_m,\alpha)
\] maka untuk \(x \ge x_m\):

\[
f(x)=\alpha \frac{x_m^\alpha}{x^{\alpha+1}}
\]

\subsection{Q --- Questions}\label{q-questions-26}

\begin{enumerate}
\def\labelenumi{\arabic{enumi}.}
\tightlist
\item
  Mengapa Pareto menghasilkan peluang ekstrem yang lebih besar?
\item
  Kapan model heavy-tail penting dalam keputusan?
\item
  Mengapa mean/variance bisa bermasalah untuk parameter tertentu?
\end{enumerate}

\subsection{A --- Apply}\label{a-apply-26}

\begin{Shaded}
\begin{Highlighting}[]
\NormalTok{alpha }\OperatorTok{=} \FloatTok{3.0}
\NormalTok{xm }\OperatorTok{=} \FloatTok{1.0}

\NormalTok{x }\OperatorTok{=}\NormalTok{ np.linspace(}\DecValTok{1}\NormalTok{, }\DecValTok{10}\NormalTok{, }\DecValTok{500}\NormalTok{)}
\NormalTok{pdf }\OperatorTok{=}\NormalTok{ stats.pareto.pdf(x, b}\OperatorTok{=}\NormalTok{alpha, scale}\OperatorTok{=}\NormalTok{xm)}

\NormalTok{plt.figure(figsize}\OperatorTok{=}\NormalTok{(}\DecValTok{8}\NormalTok{,}\DecValTok{4}\NormalTok{))}
\NormalTok{plt.plot(x, pdf)}
\NormalTok{plt.xlabel(}\StringTok{"x"}\NormalTok{)}
\NormalTok{plt.ylabel(}\StringTok{"f(x)"}\NormalTok{)}
\NormalTok{plt.title(}\StringTok{"PDF Pareto"}\NormalTok{)}
\NormalTok{plt.grid(alpha}\OperatorTok{=}\FloatTok{0.3}\NormalTok{)}
\NormalTok{plt.show()}
\end{Highlighting}
\end{Shaded}

\pandocbounded{\includegraphics[keepaspectratio]{04-distribusi-kontinu_files/figure-pdf/cell-21-output-1.pdf}}

\begin{Shaded}
\begin{Highlighting}[]
\BuiltInTok{print}\NormalTok{(}\StringTok{"P(X \textgreater{} 3) ="}\NormalTok{, }\DecValTok{1} \OperatorTok{{-}}\NormalTok{ stats.pareto.cdf(}\DecValTok{3}\NormalTok{, b}\OperatorTok{=}\NormalTok{alpha, scale}\OperatorTok{=}\NormalTok{xm))}
\end{Highlighting}
\end{Shaded}

\begin{verbatim}
P(X > 3) = 0.03703703703703698
\end{verbatim}

\subsubsection{Interpretasi}\label{interpretasi-14}

Pareto penting ketika kejadian ekstrem tidak boleh diabaikan. Dalam
pengambilan keputusan risiko, tail tebal bisa sangat menentukan.

\section{4.10 Chi-Square Distribution}\label{chi-square-distribution}

\subsection{K --- Konteks}\label{k-konteks-27}

Chi-Square muncul dalam: - inferensi statistik, - goodness-of-fit, -
analisis varians, - penjumlahan kuadrat peubah normal baku.

\subsection{M --- Model}\label{m-model-27}

Jika: \[
Z_1, Z_2, \dots, Z_k \sim \mathcal{N}(0,1)
\] independen, maka:

\[
X = Z_1^2 + Z_2^2 + \cdots + Z_k^2
\sim \chi^2(k)
\]

Distribusi ini adalah kasus khusus Gamma.

\subsection{Q --- Questions}\label{q-questions-27}

\begin{enumerate}
\def\labelenumi{\arabic{enumi}.}
\tightlist
\item
  Mengapa bentuknya hanya pada nilai non-negatif?
\item
  Apa hubungan Chi-Square dengan Normal?
\item
  Mengapa distribusi ini penting di statistika?
\end{enumerate}

\subsection{A --- Apply}\label{a-apply-27}

\begin{Shaded}
\begin{Highlighting}[]
\NormalTok{df }\OperatorTok{=} \DecValTok{5}
\NormalTok{x }\OperatorTok{=}\NormalTok{ np.linspace(}\DecValTok{0}\NormalTok{, }\DecValTok{20}\NormalTok{, }\DecValTok{500}\NormalTok{)}
\NormalTok{pdf }\OperatorTok{=}\NormalTok{ stats.chi2.pdf(x, df}\OperatorTok{=}\NormalTok{df)}

\NormalTok{plt.figure(figsize}\OperatorTok{=}\NormalTok{(}\DecValTok{8}\NormalTok{,}\DecValTok{4}\NormalTok{))}
\NormalTok{plt.plot(x, pdf)}
\NormalTok{plt.xlabel(}\StringTok{"x"}\NormalTok{)}
\NormalTok{plt.ylabel(}\StringTok{"f(x)"}\NormalTok{)}
\NormalTok{plt.title(}\StringTok{"PDF Chi{-}Square(df=5)"}\NormalTok{)}
\NormalTok{plt.grid(alpha}\OperatorTok{=}\FloatTok{0.3}\NormalTok{)}
\NormalTok{plt.show()}
\end{Highlighting}
\end{Shaded}

\pandocbounded{\includegraphics[keepaspectratio]{04-distribusi-kontinu_files/figure-pdf/cell-23-output-1.pdf}}

\begin{Shaded}
\begin{Highlighting}[]
\BuiltInTok{print}\NormalTok{(}\StringTok{"Mean ="}\NormalTok{, df)}
\BuiltInTok{print}\NormalTok{(}\StringTok{"Variance ="}\NormalTok{, }\DecValTok{2}\OperatorTok{*}\NormalTok{df)}
\BuiltInTok{print}\NormalTok{(}\StringTok{"P(X \textless{}= 7) ="}\NormalTok{, stats.chi2.cdf(}\DecValTok{7}\NormalTok{, df}\OperatorTok{=}\NormalTok{df))}
\end{Highlighting}
\end{Shaded}

\begin{verbatim}
Mean = 5
Variance = 10
P(X <= 7) = 0.7793596920632894
\end{verbatim}

\section{4.11 Hubungan Antardistribusi}\label{hubungan-antardistribusi}

Bagian ini penting karena membantu mahasiswa melihat probabilitas
sebagai jaringan ide, bukan daftar rumus terpisah.

\subsection{Exponential ↔ Poisson
Process}\label{exponential-poisson-process}

Jika banyak kejadian dalam interval mengikuti Poisson process dengan
rate \(\lambda\), maka waktu antar kedatangan mengikuti
Exponential(\(\lambda\)).

\subsection{Erlang ↔ Exponential}\label{erlang-exponential}

Erlang adalah jumlah beberapa random variable Exponential independen
dengan rate sama.

\subsection{Gamma ↔ Erlang}\label{gamma-erlang}

Erlang adalah kasus khusus Gamma dengan parameter shape berupa bilangan
bulat positif.

\subsection{Normal ↔ Binomial}\label{normal-binomial}

Binomial dengan \(n\) besar sering dapat didekati oleh Normal.

\subsection{Chi-Square ↔ Gamma}\label{chi-square-gamma}

Chi-Square adalah kasus khusus Gamma.

Mari buat diagram singkat dengan Mermaid.

\begin{Shaded}
\begin{Highlighting}[]
\NormalTok{flowchart LR}
\NormalTok{    PoissonProcess["Poisson Process"] {-}{-}\textgreater{} Exponential["Exponential (interarrival time)"]}
\NormalTok{    Exponential {-}{-}\textgreater{} Erlang["Erlang = jumlah k Exponential iid"]}
\NormalTok{    Erlang {-}{-}\textgreater{} Gamma["Gamma (generalisasi Erlang)"]}
\NormalTok{    Binomial["Binomial (n besar)"] {-}{-}\textgreater{} Normal["Normal approximation"]}
\NormalTok{    Gamma {-}{-}\textgreater{} ChiSquare["Chi{-}Square (kasus khusus Gamma)"]}
\end{Highlighting}
\end{Shaded}

\includegraphics[width=15.07in,height=2.06in]{04-distribusi-kontinu_files/figure-latex/mermaid-figure-1.png}

\section{4.12 Mini-Case KMQA 1 --- Garansi
Produk}\label{mini-case-kmqa-1-garansi-produk}

\subsection{K --- Konteks}\label{k-konteks-28}

Umur hidup lampu memengaruhi kebijakan garansi. Jika garansi terlalu
panjang, biaya klaim meningkat. Jika terlalu pendek, pelanggan kecewa.

\subsection{M --- Model}\label{m-model-28}

Misalkan umur hidup \(X \sim \mathcal{N}(900,50^2)\).

\subsection{Q --- Questions}\label{q-questions-28}

\begin{enumerate}
\def\labelenumi{\arabic{enumi}.}
\tightlist
\item
  Berapa peluang rusak sebelum 800 jam?
\item
  Berapa jika garansi 850 atau 900 jam?
\item
  Kebijakan mana yang lebih konservatif?
\end{enumerate}

\subsection{A --- Apply}\label{a-apply-28}

\begin{Shaded}
\begin{Highlighting}[]
\NormalTok{mu, sigma }\OperatorTok{=} \DecValTok{900}\NormalTok{, }\DecValTok{50}
\ControlFlowTok{for}\NormalTok{ T }\KeywordTok{in}\NormalTok{ [}\DecValTok{800}\NormalTok{, }\DecValTok{850}\NormalTok{, }\DecValTok{900}\NormalTok{]:}
    \BuiltInTok{print}\NormalTok{(}\SpecialStringTok{f"P(X \textless{} }\SpecialCharTok{\{}\NormalTok{T}\SpecialCharTok{\}}\SpecialStringTok{) ="}\NormalTok{, stats.norm.cdf(T, loc}\OperatorTok{=}\NormalTok{mu, scale}\OperatorTok{=}\NormalTok{sigma))}
\end{Highlighting}
\end{Shaded}

\begin{verbatim}
P(X < 800) = 0.0227501319481792
P(X < 850) = 0.15865525393145707
P(X < 900) = 0.5
\end{verbatim}

\subsubsection{Interpretasi}\label{interpretasi-15}

CDF langsung memberi estimasi proporsi klaim garansi.

\section{4.13 Mini-Case KMQA 2 --- Waktu Antar Kedatangan
Pasien}\label{mini-case-kmqa-2-waktu-antar-kedatangan-pasien}

\subsection{K --- Konteks}\label{k-konteks-29}

Sebuah klinik menerima pasien secara acak. Kita ingin memodelkan selang
waktu antar kedatangan.

\subsection{M --- Model}\label{m-model-29}

Jika kedatangan mengikuti Poisson process dengan rate \(\lambda\), maka
interarrival time: \[
X \sim \mathrm{Exponential}(\lambda)
\]

\subsection{Q --- Questions}\label{q-questions-29}

\begin{enumerate}
\def\labelenumi{\arabic{enumi}.}
\tightlist
\item
  Berapa peluang pasien berikut datang dalam 10 menit?
\item
  Berapa rata-rata selang tunggu?
\end{enumerate}

\subsection{A --- Apply}\label{a-apply-29}

\begin{Shaded}
\begin{Highlighting}[]
\CommentTok{\#\# rate 2 pasien/jam =\textgreater{} 1 pasien setiap 30 menit rata{-}rata}
\NormalTok{lam\_per\_hour }\OperatorTok{=} \DecValTok{2}
\NormalTok{lam\_per\_minute }\OperatorTok{=}\NormalTok{ lam\_per\_hour }\OperatorTok{/} \DecValTok{60}

\BuiltInTok{print}\NormalTok{(}\StringTok{"Rata{-}rata waktu tunggu (menit) ="}\NormalTok{, }\DecValTok{1}\OperatorTok{/}\NormalTok{lam\_per\_minute)}
\BuiltInTok{print}\NormalTok{(}\StringTok{"P(datang dalam 10 menit) ="}\NormalTok{, stats.expon.cdf(}\DecValTok{10}\NormalTok{, scale}\OperatorTok{=}\DecValTok{1}\OperatorTok{/}\NormalTok{lam\_per\_minute))}
\end{Highlighting}
\end{Shaded}

\begin{verbatim}
Rata-rata waktu tunggu (menit) = 30.0
P(datang dalam 10 menit) = 0.28346868942621073
\end{verbatim}

\section{4.14 Mini-Case KMQA 3 --- Tiga Sub-Layanan
Berantai}\label{mini-case-kmqa-3-tiga-sub-layanan-berantai}

\subsection{K --- Konteks}\label{k-konteks-30}

Sebuah layanan terdiri dari tiga tahap berurutan, masing-masing
memoryless dan independen.

\subsection{M --- Model}\label{m-model-30}

Masing-masing tahap \(X_i \sim \mathrm{Exponential}(\lambda)\), maka
total: \[
T = X_1 + X_2 + X_3 \sim \text{Erlang}(3,\lambda)
\]

\subsection{Q --- Questions}\label{q-questions-30}

\begin{enumerate}
\def\labelenumi{\arabic{enumi}.}
\tightlist
\item
  Berapa mean total durasi?
\item
  Mengapa totalnya tidak lagi Exponential?
\item
  Bagaimana histogram Monte Carlo dibanding PDF Erlang?
\end{enumerate}

\subsection{A --- Apply}\label{a-apply-30}

\begin{Shaded}
\begin{Highlighting}[]
\NormalTok{rng }\OperatorTok{=}\NormalTok{ np.random.default\_rng(}\DecValTok{42}\NormalTok{)}
\NormalTok{lam }\OperatorTok{=} \FloatTok{0.5}
\NormalTok{n }\OperatorTok{=} \DecValTok{100000}

\NormalTok{X1 }\OperatorTok{=}\NormalTok{ rng.exponential(scale}\OperatorTok{=}\DecValTok{1}\OperatorTok{/}\NormalTok{lam, size}\OperatorTok{=}\NormalTok{n)}
\NormalTok{X2 }\OperatorTok{=}\NormalTok{ rng.exponential(scale}\OperatorTok{=}\DecValTok{1}\OperatorTok{/}\NormalTok{lam, size}\OperatorTok{=}\NormalTok{n)}
\NormalTok{X3 }\OperatorTok{=}\NormalTok{ rng.exponential(scale}\OperatorTok{=}\DecValTok{1}\OperatorTok{/}\NormalTok{lam, size}\OperatorTok{=}\NormalTok{n)}
\NormalTok{T }\OperatorTok{=}\NormalTok{ X1 }\OperatorTok{+}\NormalTok{ X2 }\OperatorTok{+}\NormalTok{ X3}

\NormalTok{x }\OperatorTok{=}\NormalTok{ np.linspace(}\DecValTok{0}\NormalTok{, np.percentile(T, }\FloatTok{99.5}\NormalTok{), }\DecValTok{400}\NormalTok{)}
\NormalTok{pdf }\OperatorTok{=}\NormalTok{ stats.gamma.pdf(x, a}\OperatorTok{=}\DecValTok{3}\NormalTok{, scale}\OperatorTok{=}\DecValTok{1}\OperatorTok{/}\NormalTok{lam)}

\NormalTok{plt.figure(figsize}\OperatorTok{=}\NormalTok{(}\DecValTok{8}\NormalTok{,}\DecValTok{4}\NormalTok{))}
\NormalTok{plt.hist(T, bins}\OperatorTok{=}\DecValTok{60}\NormalTok{, density}\OperatorTok{=}\VariableTok{True}\NormalTok{, alpha}\OperatorTok{=}\FloatTok{0.6}\NormalTok{, label}\OperatorTok{=}\StringTok{"Monte Carlo"}\NormalTok{)}
\NormalTok{plt.plot(x, pdf, label}\OperatorTok{=}\StringTok{"PDF Erlang{-}3"}\NormalTok{)}
\NormalTok{plt.xlabel(}\StringTok{"Durasi total"}\NormalTok{)}
\NormalTok{plt.ylabel(}\StringTok{"Density"}\NormalTok{)}
\NormalTok{plt.title(}\StringTok{"Jumlah 3 Exponential iid menghasilkan Erlang{-}3"}\NormalTok{)}
\NormalTok{plt.legend()}
\NormalTok{plt.grid(alpha}\OperatorTok{=}\FloatTok{0.3}\NormalTok{)}
\NormalTok{plt.show()}
\end{Highlighting}
\end{Shaded}

\pandocbounded{\includegraphics[keepaspectratio]{04-distribusi-kontinu_files/figure-pdf/cell-27-output-1.pdf}}

\section{4.15 Python Toolbox untuk Distribusi
Kontinu}\label{python-toolbox-untuk-distribusi-kontinu}

\begin{Shaded}
\begin{Highlighting}[]
\ImportTok{from}\NormalTok{ scipy }\ImportTok{import}\NormalTok{ stats}
\ImportTok{import}\NormalTok{ numpy }\ImportTok{as}\NormalTok{ np}
\ImportTok{import}\NormalTok{ matplotlib.pyplot }\ImportTok{as}\NormalTok{ plt}
\end{Highlighting}
\end{Shaded}

\subsection{Normal}\label{normal}

\begin{Shaded}
\begin{Highlighting}[]
\NormalTok{stats.norm.pdf(}\DecValTok{0}\NormalTok{, loc}\OperatorTok{=}\DecValTok{0}\NormalTok{, scale}\OperatorTok{=}\DecValTok{1}\NormalTok{), stats.norm.cdf(}\FloatTok{1.96}\NormalTok{, loc}\OperatorTok{=}\DecValTok{0}\NormalTok{, scale}\OperatorTok{=}\DecValTok{1}\NormalTok{)}
\end{Highlighting}
\end{Shaded}

\begin{verbatim}
(np.float64(0.3989422804014327), np.float64(0.9750021048517795))
\end{verbatim}

\subsection{Exponential}\label{exponential}

\begin{Shaded}
\begin{Highlighting}[]
\NormalTok{stats.expon.pdf(}\DecValTok{2}\NormalTok{, scale}\OperatorTok{=}\DecValTok{3}\NormalTok{), stats.expon.cdf(}\DecValTok{2}\NormalTok{, scale}\OperatorTok{=}\DecValTok{3}\NormalTok{)}
\end{Highlighting}
\end{Shaded}

\begin{verbatim}
(np.float64(0.17113903967753066), np.float64(0.486582880967408))
\end{verbatim}

\subsection{Gamma}\label{gamma}

\begin{Shaded}
\begin{Highlighting}[]
\NormalTok{stats.gamma.pdf(}\DecValTok{4}\NormalTok{, a}\OperatorTok{=}\DecValTok{3}\NormalTok{, scale}\OperatorTok{=}\DecValTok{2}\NormalTok{), stats.gamma.cdf(}\DecValTok{4}\NormalTok{, a}\OperatorTok{=}\DecValTok{3}\NormalTok{, scale}\OperatorTok{=}\DecValTok{2}\NormalTok{)}
\end{Highlighting}
\end{Shaded}

\begin{verbatim}
(np.float64(0.1353352832366127), np.float64(0.32332358381693654))
\end{verbatim}

\subsection{Weibull}\label{weibull}

\begin{Shaded}
\begin{Highlighting}[]
\NormalTok{stats.weibull\_min.pdf(}\DecValTok{4}\NormalTok{, c}\OperatorTok{=}\DecValTok{2}\NormalTok{, scale}\OperatorTok{=}\DecValTok{10}\NormalTok{), stats.weibull\_min.cdf(}\DecValTok{4}\NormalTok{, c}\OperatorTok{=}\DecValTok{2}\NormalTok{, scale}\OperatorTok{=}\DecValTok{10}\NormalTok{)}
\end{Highlighting}
\end{Shaded}

\begin{verbatim}
(np.float64(0.06817150311729692), np.float64(0.14785621103378868))
\end{verbatim}

\subsection{Pareto}\label{pareto}

\begin{Shaded}
\begin{Highlighting}[]
\NormalTok{stats.pareto.pdf(}\DecValTok{2}\NormalTok{, b}\OperatorTok{=}\DecValTok{3}\NormalTok{, scale}\OperatorTok{=}\DecValTok{1}\NormalTok{), stats.pareto.cdf(}\DecValTok{2}\NormalTok{, b}\OperatorTok{=}\DecValTok{3}\NormalTok{, scale}\OperatorTok{=}\DecValTok{1}\NormalTok{)}
\end{Highlighting}
\end{Shaded}

\begin{verbatim}
(np.float64(0.1875), np.float64(0.875))
\end{verbatim}

\subsection{Chi-Square}\label{chi-square}

\begin{Shaded}
\begin{Highlighting}[]
\NormalTok{stats.chi2.pdf(}\DecValTok{5}\NormalTok{, df}\OperatorTok{=}\DecValTok{4}\NormalTok{), stats.chi2.cdf(}\DecValTok{5}\NormalTok{, df}\OperatorTok{=}\DecValTok{4}\NormalTok{)}
\end{Highlighting}
\end{Shaded}

\begin{verbatim}
(np.float64(0.10260624827987348), np.float64(0.7127025048163542))
\end{verbatim}

\section{4.16 Kesalahan Umum yang Harus
Dihindari}\label{kesalahan-umum-yang-harus-dihindari-1}

\begin{enumerate}
\def\labelenumi{\arabic{enumi}.}
\item
  \textbf{Menganggap tinggi PDF adalah peluang}\\
  Tidak. Peluang adalah area di bawah kurva.
\item
  \textbf{Menghitung \(P(X=x)\) untuk kontinu seolah-olah positif}\\
  Dalam model kontinu, peluang titik tunggal adalah nol.
\item
  \textbf{Memakai Exponential untuk semua lifetime}\\
  Exponential hanya cocok bila hazard rate konstan / memoryless.
\item
  \textbf{Memakai Normal untuk data yang sangat skewed atau bernilai
  negatif-padahal tidak mungkin}\\
  Pilih model sesuai konteks.
\item
  \textbf{Menghafal distribusi tanpa memahami makna parameter}\\
  Yang penting bukan sekadar nama distribusinya, tetapi arti shape,
  scale, mean, variance, dan implikasinya.
\end{enumerate}

\section{4.17 Menyimpulkan Bab Ini}\label{menyimpulkan-bab-ini-3}

Distribusi kontinu memberi kita bahasa untuk memodelkan banyak besaran
dunia nyata yang berupa pengukuran, waktu, umur hidup, dan nilai-nilai
positif yang berubah halus.

Kita telah melihat: - \textbf{uniform kontinu} untuk interval yang
merata, - \textbf{normal} untuk fenomena simetris dan agregatif, -
\textbf{gamma} untuk waktu tunggu positif yang fleksibel, -
\textbf{exponential} untuk memoryless waiting time, - \textbf{erlang}
untuk jumlah beberapa exponential, - \textbf{weibull} untuk reliability
yang lebih fleksibel, - \textbf{pareto} untuk heavy-tail, -
\textbf{chi-square} untuk inferensi statistik dan hubungan dengan
normal.

Di samping itu, hubungan antardistribusi membantu kita melihat
probabilitas sebagai sistem ide yang saling terhubung.

\section{4.18 Ringkasan Poin Inti}\label{ringkasan-poin-inti-3}

\begin{enumerate}
\def\labelenumi{\arabic{enumi}.}
\tightlist
\item
  Untuk random variable kontinu, peluang diperoleh dari \textbf{area} di
  bawah PDF.
\item
  Distribusi kontinu kustom bisa dibangun dari tabel atau interpolasi
  numerik.
\item
  \textbf{Uniform kontinu} cocok untuk interval merata.
\item
  \textbf{Normal} penting untuk banyak fenomena simetris dan sebagai
  pendekatan Binomial.
\item
  \textbf{Exponential} memodelkan waiting time memoryless dan terkait
  erat dengan Poisson process.
\item
  \textbf{Gamma} dan \textbf{Erlang} memodelkan akumulasi waiting time.
\item
  \textbf{Weibull} sangat penting dalam reliability engineering.
\item
  \textbf{Pareto} penting untuk fenomena ekstrem dan heavy-tail.
\item
  \textbf{Chi-Square} adalah distribusi penting dalam statistika
  inferensial.
\item
  Memahami hubungan antardistribusi membantu membangun intuisi yang
  lebih dalam.
\end{enumerate}

\section{4.19 Latihan Bab 4}\label{latihan-bab-4}

\subsection{A. Konseptual}\label{a.-konseptual-3}

\begin{enumerate}
\def\labelenumi{\arabic{enumi}.}
\tightlist
\item
  Mengapa peluang titik tunggal pada RV kontinu bernilai nol?
\item
  Apa perbedaan utama antara PDF dan CDF?
\item
  Mengapa Exponential disebut memoryless?
\item
  Mengapa Weibull sering lebih realistis daripada Exponential untuk data
  kegagalan?
\end{enumerate}

\subsection{B. Hitungan}\label{b.-hitungan-2}

\begin{enumerate}
\def\labelenumi{\arabic{enumi}.}
\tightlist
\item
  Jika \(X \sim \mathrm{Uniform}(2,8)\), hitung \(P(3 \le X \le 5)\).
\item
  Jika \(X \sim \mathcal{N}(100,15^2)\), hitung \(P(X<90)\).
\item
  Jika \(X \sim \mathrm{Exponential}(\lambda=0.5)\), hitung
  \(P(X \le 3)\).
\item
  Jika \(X \sim \mathrm{Gamma}(3,2)\), hitung \(E[X]\) dan
  \(\mathrm{Var}(X)\).
\end{enumerate}

\subsection{C. Python}\label{c.-python-2}

\begin{enumerate}
\def\labelenumi{\arabic{enumi}.}
\tightlist
\item
  Simulasikan 10000 sampel Uniform(0,1) dan bandingkan histogram dengan
  PDF.
\item
  Simulasikan Normal dan bandingkan histogram dengan PDF.
\item
  Simulasikan interarrival time Poisson process dan bandingkan dengan
  Exponential teoritis.
\item
  Simulasikan jumlah tiga Exponential dan bandingkan dengan Erlang
  teoritis.
\end{enumerate}

\subsection{D. Aplikatif}\label{d.-aplikatif-2}

\begin{enumerate}
\def\labelenumi{\arabic{enumi}.}
\tightlist
\item
  Pilih model yang paling cocok untuk:

  \begin{itemize}
  \tightlist
  \item
    tinggi badan,
  \item
    waktu antar kedatangan pelanggan,
  \item
    lifetime komponen dengan keausan,
  \item
    kerugian ekstrem.
  \end{itemize}
\item
  Jelaskan alasan pemilihan model tersebut.
\item
  Dalam kasus garansi produk, ukuran apa yang paling penting: mean umur,
  variance umur, atau peluang rusak sebelum T?
\end{enumerate}

\section{4.20 Penutup Kecil}\label{penutup-kecil-3}

Kalau bab sebelumnya mengajarkan kita menghitung kejadian yang berupa
hitungan, maka bab ini mengajarkan kita berpikir tentang besaran yang
mengalir: waktu, umur hidup, pengukuran, dan intensitas.

Di bab berikutnya, kita akan memperluas pandangan lagi: dunia nyata
jarang hanya punya satu peubah acak. Kita akan mulai melihat
\textbf{lebih dari satu random variable sekaligus}, hubungan di
antaranya, serta apa yang terjadi ketika satu random variable menjadi
fungsi dari yang lain.

\bookmarksetup{startatroot}

\chapter{Bab 5. Random Variable Multivariat dan Fungsi Random
Variable}\label{bab-5.-random-variable-multivariat-dan-fungsi-random-variable}

\section{Tujuan Bab}\label{tujuan-bab-4}

Setelah mempelajari bab ini, mahasiswa diharapkan mampu:

\begin{enumerate}
\def\labelenumi{\arabic{enumi}.}
\tightlist
\item
  memahami bahwa banyak masalah nyata melibatkan \textbf{lebih dari satu
  random variable},
\item
  membedakan distribusi \textbf{joint}, \textbf{marginal}, dan
  \textbf{conditional},
\item
  memahami makna \textbf{independensi} pada random variable multivariat,
\item
  menghitung dan menafsirkan \textbf{kovariansi} dan \textbf{korelasi},
\item
  memahami bahwa fungsi dari random variable juga merupakan random
  variable,
\item
  menggunakan Python untuk simulasi, visualisasi, dan eksplorasi
  hubungan antar peubah acak,
\item
  menghubungkan konsep-konsep ini dengan pengambilan keputusan di bidang
  teknik, bisnis, dan layanan.
\end{enumerate}

\section{Pembuka}\label{pembuka-4}

Sampai bab sebelumnya, kita sering berbicara seolah-olah dunia hanya
memiliki \textbf{satu} random variable pada satu waktu: - jumlah cacat,
- waktu tunggu, - umur hidup, - profit, - jumlah pelanggan.

Namun kenyataan jauh lebih kaya. Dalam dunia nyata, keputusan hampir
selalu melibatkan \textbf{lebih dari satu peubah acak}.

Contoh: - tinggi dan berat badan, - permintaan dan pendapatan, - jumlah
pelanggan dan waktu layanan, - suhu dan konsumsi listrik, - arus dan
daya, - umur komponen dan biaya klaim.

Begitu ada lebih dari satu random variable, muncul pertanyaan-pertanyaan
baru: - apakah mereka saling memengaruhi? - apakah keduanya independen?
- kalau satu naik, apakah yang lain cenderung naik? - bagaimana
distribusi salah satunya bila yang lain diketahui? - bagaimana
distribusi output jika output adalah fungsi dari input yang acak?

Bab ini menjawab pertanyaan-pertanyaan itu.

\section{5.1 Quick Win: Dua Peubah, Satu
Plot}\label{quick-win-dua-peubah-satu-plot}

Mari mulai dengan simulasi dua random variable yang berkorelasi positif.
Misalnya: - \(X\) = jam belajar, - \(Y\) = skor kuis.

Ini tentu penyederhanaan, tetapi cukup baik untuk melihat ide dasar.

\begin{Shaded}
\begin{Highlighting}[]
\ImportTok{import}\NormalTok{ numpy }\ImportTok{as}\NormalTok{ np}
\ImportTok{import}\NormalTok{ matplotlib.pyplot }\ImportTok{as}\NormalTok{ plt}
\ImportTok{from}\NormalTok{ scipy }\ImportTok{import}\NormalTok{ stats}

\NormalTok{rng }\OperatorTok{=}\NormalTok{ np.random.default\_rng(}\DecValTok{42}\NormalTok{)}

\NormalTok{mean }\OperatorTok{=}\NormalTok{ [}\DecValTok{5}\NormalTok{, }\DecValTok{70}\NormalTok{]}
\NormalTok{cov }\OperatorTok{=}\NormalTok{ [[}\FloatTok{1.5}\NormalTok{, }\DecValTok{8}\NormalTok{],}
\NormalTok{       [}\DecValTok{8}\NormalTok{, }\DecValTok{80}\NormalTok{]]}

\NormalTok{data }\OperatorTok{=}\NormalTok{ rng.multivariate\_normal(mean, cov, size}\OperatorTok{=}\DecValTok{1000}\NormalTok{)}
\NormalTok{X }\OperatorTok{=}\NormalTok{ data[:, }\DecValTok{0}\NormalTok{]}
\NormalTok{Y }\OperatorTok{=}\NormalTok{ data[:, }\DecValTok{1}\NormalTok{]}

\NormalTok{plt.figure(figsize}\OperatorTok{=}\NormalTok{(}\DecValTok{6}\NormalTok{,}\DecValTok{5}\NormalTok{))}
\NormalTok{plt.scatter(X, Y, alpha}\OperatorTok{=}\FloatTok{0.4}\NormalTok{, s}\OperatorTok{=}\DecValTok{15}\NormalTok{)}
\NormalTok{plt.xlabel(}\StringTok{"Jam belajar"}\NormalTok{)}
\NormalTok{plt.ylabel(}\StringTok{"Skor kuis"}\NormalTok{)}
\NormalTok{plt.title(}\StringTok{"Dua random variable dengan kecenderungan positif"}\NormalTok{)}
\NormalTok{plt.grid(alpha}\OperatorTok{=}\FloatTok{0.3}\NormalTok{)}
\NormalTok{plt.show()}
\end{Highlighting}
\end{Shaded}

\pandocbounded{\includegraphics[keepaspectratio]{05-random-variable-multivariat-dan-fungsi_files/figure-pdf/cell-2-output-1.pdf}}

\begin{Shaded}
\begin{Highlighting}[]
\BuiltInTok{print}\NormalTok{(}\StringTok{"Mean X ="}\NormalTok{, X.mean())}
\BuiltInTok{print}\NormalTok{(}\StringTok{"Mean Y ="}\NormalTok{, Y.mean())}
\BuiltInTok{print}\NormalTok{(}\StringTok{"Cov(X,Y) ="}\NormalTok{, np.cov(X, Y, ddof}\OperatorTok{=}\DecValTok{0}\NormalTok{)[}\DecValTok{0}\NormalTok{,}\DecValTok{1}\NormalTok{])}
\BuiltInTok{print}\NormalTok{(}\StringTok{"Corr(X,Y) ="}\NormalTok{, np.corrcoef(X, Y)[}\DecValTok{0}\NormalTok{,}\DecValTok{1}\NormalTok{])}
\end{Highlighting}
\end{Shaded}

\begin{verbatim}
Mean X = 5.0966118178007385
Mean Y = 70.63531660231106
Cov(X,Y) = 8.268878410266762
Corr(X,Y) = 0.7409518158857897
\end{verbatim}

Dari \emph{quick win} ini, kita langsung melihat dua hal: 1. dunia nyata
sering lebih masuk akal bila dimodelkan dengan lebih dari satu peubah
acak, 2. hubungan antar peubah dapat diringkas dengan konsep seperti
\textbf{covariance} dan \textbf{correlation}.

\section{5.2 Joint Distribution: Distribusi
Bersama}\label{joint-distribution-distribusi-bersama}

\subsection{K --- Konteks}\label{k-konteks-31}

Misalkan: - \(X\) = jumlah pelanggan datang per jam, - \(Y\) = total
pendapatan per jam.

Atau: - \(X\) = jumlah produk cacat dari supplier A, - \(Y\) = jumlah
produk cacat dari supplier B.

Kita tidak cukup hanya tahu distribusi \(X\) sendiri atau \(Y\) sendiri.
Kita ingin tahu \textbf{bagaimana keduanya muncul bersama}.

\subsection{M --- Model}\label{m-model-31}

Distribusi bersama atau \textbf{joint distribution} mendeskripsikan
peluang dua random variable secara simultan.

\subsubsection{Untuk diskrit}\label{untuk-diskrit-1}

\[
p_{X,Y}(x,y)=P(X=x, Y=y)
\]

\subsubsection{Untuk kontinu}\label{untuk-kontinu-1}

\[
f_{X,Y}(x,y)
\]

Joint distribution memberi informasi lengkap tentang pasangan \((X,Y)\).

\subsection{Q --- Questions}\label{q-questions-31}

\begin{enumerate}
\def\labelenumi{\arabic{enumi}.}
\tightlist
\item
  Berapa peluang \(X\) dan \(Y\) mengambil nilai tertentu secara
  bersamaan?
\item
  Apakah nilai besar pada \(X\) sering disertai nilai besar pada \(Y\)?
\item
  Bagaimana menurunkan distribusi salah satu peubah dari distribusi
  bersama?
\end{enumerate}

\subsection{A --- Apply (diskrit
sederhana)}\label{a-apply-diskrit-sederhana}

Misalkan dua koin fair dilempar. Definisikan: - \(X\) = jumlah Head, -
\(Y\) = indikator apakah koin pertama Head (1) atau tidak (0).

Mari bangun joint distribution-nya.

\begin{Shaded}
\begin{Highlighting}[]
\ImportTok{import}\NormalTok{ itertools}

\NormalTok{outcomes }\OperatorTok{=} \BuiltInTok{list}\NormalTok{(itertools.product([}\StringTok{\textquotesingle{}H\textquotesingle{}}\NormalTok{, }\StringTok{\textquotesingle{}T\textquotesingle{}}\NormalTok{], repeat}\OperatorTok{=}\DecValTok{2}\NormalTok{))}

\KeywordTok{def}\NormalTok{ X\_fun(o):}
    \ControlFlowTok{return}\NormalTok{ o.count(}\StringTok{\textquotesingle{}H\textquotesingle{}}\NormalTok{)}

\KeywordTok{def}\NormalTok{ Y\_fun(o):}
    \ControlFlowTok{return} \DecValTok{1} \ControlFlowTok{if}\NormalTok{ o[}\DecValTok{0}\NormalTok{] }\OperatorTok{==} \StringTok{\textquotesingle{}H\textquotesingle{}} \ControlFlowTok{else} \DecValTok{0}

\NormalTok{pairs }\OperatorTok{=}\NormalTok{ [(X\_fun(o), Y\_fun(o)) }\ControlFlowTok{for}\NormalTok{ o }\KeywordTok{in}\NormalTok{ outcomes]}
\NormalTok{pairs}
\end{Highlighting}
\end{Shaded}

\begin{verbatim}
[(2, 1), (1, 1), (1, 0), (0, 0)]
\end{verbatim}

\begin{Shaded}
\begin{Highlighting}[]
\ImportTok{from}\NormalTok{ collections }\ImportTok{import}\NormalTok{ Counter}

\NormalTok{counter }\OperatorTok{=}\NormalTok{ Counter(pairs)}
\NormalTok{joint }\OperatorTok{=}\NormalTok{ \{k: v}\OperatorTok{/}\BuiltInTok{len}\NormalTok{(outcomes) }\ControlFlowTok{for}\NormalTok{ k,v }\KeywordTok{in}\NormalTok{ counter.items()\}}
\NormalTok{joint}
\end{Highlighting}
\end{Shaded}

\begin{verbatim}
{(2, 1): 0.25, (1, 1): 0.25, (1, 0): 0.25, (0, 0): 0.25}
\end{verbatim}

Joint distribution ini memberi peluang bersama, misalnya: \[
P(X=2, Y=1)
\]

yang berarti: - jumlah Head = 2, - dan koin pertama = Head.

\section{5.3 Bivariat Diskrit}\label{bivariat-diskrit}

Mari buat contoh yang lebih sistematis.

\subsection{Contoh}\label{contoh}

Misalkan joint PMF didefinisikan sebagai tabel berikut:

\[
p_{X,Y}(x,y)
\]

dengan \(X \in \{0,1,2\}\) dan \(Y \in \{0,1\}\):

{\def\LTcaptype{none} % do not increment counter
\begin{longtable}[]{@{}lrr@{}}
\toprule\noalign{}
\(x \backslash y\) & 0 & 1 \\
\midrule\noalign{}
\endhead
\bottomrule\noalign{}
\endlastfoot
0 & 0.10 & 0.20 \\
1 & 0.15 & 0.25 \\
2 & 0.10 & 0.20 \\
\end{longtable}
}

Total = 1.

\subsection{Python representation}\label{python-representation}

\begin{Shaded}
\begin{Highlighting}[]
\NormalTok{joint\_table }\OperatorTok{=}\NormalTok{ np.array([}
\NormalTok{    [}\FloatTok{0.10}\NormalTok{, }\FloatTok{0.20}\NormalTok{],}
\NormalTok{    [}\FloatTok{0.15}\NormalTok{, }\FloatTok{0.25}\NormalTok{],}
\NormalTok{    [}\FloatTok{0.10}\NormalTok{, }\FloatTok{0.20}\NormalTok{]}
\NormalTok{])}

\BuiltInTok{print}\NormalTok{(}\StringTok{"Total probability ="}\NormalTok{, joint\_table.}\BuiltInTok{sum}\NormalTok{())}
\end{Highlighting}
\end{Shaded}

\begin{verbatim}
Total probability = 1.0
\end{verbatim}

\subsubsection{Visualisasi heatmap}\label{visualisasi-heatmap}

\begin{Shaded}
\begin{Highlighting}[]
\NormalTok{plt.figure(figsize}\OperatorTok{=}\NormalTok{(}\DecValTok{5}\NormalTok{,}\DecValTok{4}\NormalTok{))}
\NormalTok{plt.imshow(joint\_table, cmap}\OperatorTok{=}\StringTok{\textquotesingle{}Blues\textquotesingle{}}\NormalTok{, origin}\OperatorTok{=}\StringTok{\textquotesingle{}lower\textquotesingle{}}\NormalTok{)}
\NormalTok{plt.colorbar(label}\OperatorTok{=}\StringTok{\textquotesingle{}Probability\textquotesingle{}}\NormalTok{)}
\NormalTok{plt.xticks([}\DecValTok{0}\NormalTok{,}\DecValTok{1}\NormalTok{], [}\DecValTok{0}\NormalTok{,}\DecValTok{1}\NormalTok{])}
\NormalTok{plt.yticks([}\DecValTok{0}\NormalTok{,}\DecValTok{1}\NormalTok{,}\DecValTok{2}\NormalTok{], [}\DecValTok{0}\NormalTok{,}\DecValTok{1}\NormalTok{,}\DecValTok{2}\NormalTok{])}
\NormalTok{plt.xlabel(}\StringTok{"Y"}\NormalTok{)}
\NormalTok{plt.ylabel(}\StringTok{"X"}\NormalTok{)}
\NormalTok{plt.title(}\StringTok{"Joint PMF bivariat diskrit"}\NormalTok{)}
\NormalTok{plt.show()}
\end{Highlighting}
\end{Shaded}

\pandocbounded{\includegraphics[keepaspectratio]{05-random-variable-multivariat-dan-fungsi_files/figure-pdf/cell-7-output-1.pdf}}

\section{5.4 Marginal Distribution}\label{marginal-distribution}

\subsection{K --- Konteks}\label{k-konteks-32}

Kadang kita hanya ingin distribusi satu peubah, tanpa memperhatikan
peubah lainnya.

\subsection{M --- Model}\label{m-model-32}

Dari joint distribution, kita bisa memperoleh distribusi
\textbf{marginal}.

\subsubsection{Untuk diskrit}\label{untuk-diskrit-2}

\[
p_X(x)=\sum_y p_{X,Y}(x,y)
\]

\[
p_Y(y)=\sum_x p_{X,Y}(x,y)
\]

\subsection{Q --- Questions}\label{q-questions-32}

\begin{enumerate}
\def\labelenumi{\arabic{enumi}.}
\tightlist
\item
  Bagaimana ``mengabaikan'' peubah lain secara formal?
\item
  Apa arti marginal dalam keputusan?
\end{enumerate}

\subsection{A --- Apply}\label{a-apply-31}

\begin{Shaded}
\begin{Highlighting}[]
\NormalTok{pX }\OperatorTok{=}\NormalTok{ joint\_table.}\BuiltInTok{sum}\NormalTok{(axis}\OperatorTok{=}\DecValTok{1}\NormalTok{)  }\CommentTok{\#\# jumlahkan sepanjang kolom}
\NormalTok{pY }\OperatorTok{=}\NormalTok{ joint\_table.}\BuiltInTok{sum}\NormalTok{(axis}\OperatorTok{=}\DecValTok{0}\NormalTok{)  }\CommentTok{\#\# jumlahkan sepanjang baris}

\BuiltInTok{print}\NormalTok{(}\StringTok{"Marginal pX ="}\NormalTok{, pX)}
\BuiltInTok{print}\NormalTok{(}\StringTok{"Marginal pY ="}\NormalTok{, pY)}
\end{Highlighting}
\end{Shaded}

\begin{verbatim}
Marginal pX = [0.3 0.4 0.3]
Marginal pY = [0.35 0.65]
\end{verbatim}

\begin{Shaded}
\begin{Highlighting}[]
\NormalTok{plt.figure(figsize}\OperatorTok{=}\NormalTok{(}\DecValTok{8}\NormalTok{,}\DecValTok{3}\NormalTok{))}

\NormalTok{plt.subplot(}\DecValTok{1}\NormalTok{,}\DecValTok{2}\NormalTok{,}\DecValTok{1}\NormalTok{)}
\NormalTok{plt.bar([}\DecValTok{0}\NormalTok{,}\DecValTok{1}\NormalTok{,}\DecValTok{2}\NormalTok{], pX)}
\NormalTok{plt.title(}\StringTok{"Marginal pX"}\NormalTok{)}
\NormalTok{plt.xlabel(}\StringTok{"x"}\NormalTok{)}
\NormalTok{plt.ylabel(}\StringTok{"P(X=x)"}\NormalTok{)}
\NormalTok{plt.grid(alpha}\OperatorTok{=}\FloatTok{0.3}\NormalTok{)}

\NormalTok{plt.subplot(}\DecValTok{1}\NormalTok{,}\DecValTok{2}\NormalTok{,}\DecValTok{2}\NormalTok{)}
\NormalTok{plt.bar([}\DecValTok{0}\NormalTok{,}\DecValTok{1}\NormalTok{], pY)}
\NormalTok{plt.title(}\StringTok{"Marginal pY"}\NormalTok{)}
\NormalTok{plt.xlabel(}\StringTok{"y"}\NormalTok{)}
\NormalTok{plt.ylabel(}\StringTok{"P(Y=y)"}\NormalTok{)}
\NormalTok{plt.grid(alpha}\OperatorTok{=}\FloatTok{0.3}\NormalTok{)}

\NormalTok{plt.tight\_layout()}
\NormalTok{plt.show()}
\end{Highlighting}
\end{Shaded}

\pandocbounded{\includegraphics[keepaspectratio]{05-random-variable-multivariat-dan-fungsi_files/figure-pdf/cell-9-output-1.pdf}}

\subsubsection{Interpretasi}\label{interpretasi-16}

Marginal distribution memberi gambaran perilaku satu peubah ketika
peubah lain tidak diperhatikan secara eksplisit.

\section{5.5 Conditional Distribution}\label{conditional-distribution}

\subsection{K --- Konteks}\label{k-konteks-33}

Dalam keputusan nyata, sering kali kita sudah mengetahui sebagian
informasi. Misalnya: - distribusi pendapatan per jam \textbf{dengan
syarat} jumlah pelanggan = 10, - distribusi nilai ujian \textbf{dengan
syarat} jam belajar \textgreater{} 5, - distribusi demand \textbf{dengan
syarat} hari hujan.

\subsection{M --- Model}\label{m-model-33}

\subsubsection{Untuk diskrit}\label{untuk-diskrit-3}

Jika \(P(Y=y)>0\), maka:

\[
P(X=x \mid Y=y)=\frac{P(X=x, Y=y)}{P(Y=y)}
\]

\subsection{Q --- Questions}\label{q-questions-33}

\begin{enumerate}
\def\labelenumi{\arabic{enumi}.}
\tightlist
\item
  Bagaimana distribusi satu peubah berubah jika yang lain diketahui?
\item
  Mengapa conditional distribution penting untuk inferensi dan
  keputusan?
\end{enumerate}

\subsection{A --- Apply}\label{a-apply-32}

Hitung \(P(X=x \mid Y=1)\) dari tabel joint sebelumnya.

\begin{Shaded}
\begin{Highlighting}[]
\CommentTok{\#\# Conditional p(X|Y=1)}
\NormalTok{cond\_X\_given\_Y1 }\OperatorTok{=}\NormalTok{ joint\_table[:, }\DecValTok{1}\NormalTok{] }\OperatorTok{/}\NormalTok{ pY[}\DecValTok{1}\NormalTok{]}
\BuiltInTok{print}\NormalTok{(}\StringTok{"P(X=x | Y=1) ="}\NormalTok{, cond\_X\_given\_Y1)}
\BuiltInTok{print}\NormalTok{(}\StringTok{"Jumlah ="}\NormalTok{, cond\_X\_given\_Y1.}\BuiltInTok{sum}\NormalTok{())}
\end{Highlighting}
\end{Shaded}

\begin{verbatim}
P(X=x | Y=1) = [0.30769231 0.38461538 0.30769231]
Jumlah = 1.0
\end{verbatim}

\begin{Shaded}
\begin{Highlighting}[]
\NormalTok{plt.figure(figsize}\OperatorTok{=}\NormalTok{(}\DecValTok{6}\NormalTok{,}\DecValTok{4}\NormalTok{))}
\NormalTok{plt.bar([}\DecValTok{0}\NormalTok{,}\DecValTok{1}\NormalTok{,}\DecValTok{2}\NormalTok{], cond\_X\_given\_Y1)}
\NormalTok{plt.xlabel(}\StringTok{"x"}\NormalTok{)}
\NormalTok{plt.ylabel(}\StringTok{"P(X=x | Y=1)"}\NormalTok{)}
\NormalTok{plt.title(}\StringTok{"Conditional distribution X | Y=1"}\NormalTok{)}
\NormalTok{plt.grid(alpha}\OperatorTok{=}\FloatTok{0.3}\NormalTok{)}
\NormalTok{plt.show()}
\end{Highlighting}
\end{Shaded}

\pandocbounded{\includegraphics[keepaspectratio]{05-random-variable-multivariat-dan-fungsi_files/figure-pdf/cell-11-output-1.pdf}}

\subsubsection{Interpretasi}\label{interpretasi-17}

Conditional distribution menunjukkan bahwa ketika informasi tambahan
tersedia, distribusi yang relevan untuk keputusan juga berubah.

\section{5.6 Independensi}\label{independensi}

\subsection{K --- Konteks}\label{k-konteks-34}

Kadang dua peubah tampak tidak saling memengaruhi. Kadang justru sangat
terkait. Membedakan dua situasi ini penting.

\subsection{M --- Model}\label{m-model-34}

Dua random variable \(X\) dan \(Y\) dikatakan \textbf{independen} jika:

\subsubsection{Untuk diskrit}\label{untuk-diskrit-4}

\[
P(X=x, Y=y)=P(X=x)P(Y=y)
\]

untuk semua \(x,y\).

\subsubsection{Untuk kontinu}\label{untuk-kontinu-2}

\[
f_{X,Y}(x,y)=f_X(x)f_Y(y)
\]

\subsection{Q --- Questions}\label{q-questions-34}

\begin{enumerate}
\def\labelenumi{\arabic{enumi}.}
\tightlist
\item
  Bagaimana menguji independensi dari joint distribution?
\item
  Apakah korelasi nol berarti pasti independen?
\item
  Kapan independensi terlalu kuat sebagai asumsi?
\end{enumerate}

\subsection{A --- Apply}\label{a-apply-33}

Mari cek apakah joint\_table tadi independen.

\begin{Shaded}
\begin{Highlighting}[]
\NormalTok{product\_table }\OperatorTok{=}\NormalTok{ np.outer(pX, pY)}

\BuiltInTok{print}\NormalTok{(}\StringTok{"Joint table:}\CharTok{\textbackslash{}n}\StringTok{"}\NormalTok{, joint\_table)}
\BuiltInTok{print}\NormalTok{(}\StringTok{"}\CharTok{\textbackslash{}n}\StringTok{Product of marginals:}\CharTok{\textbackslash{}n}\StringTok{"}\NormalTok{, product\_table)}
\BuiltInTok{print}\NormalTok{(}\StringTok{"}\CharTok{\textbackslash{}n}\StringTok{Apakah sama? "}\NormalTok{, np.allclose(joint\_table, product\_table))}
\end{Highlighting}
\end{Shaded}

\begin{verbatim}
Joint table:
 [[0.1  0.2 ]
 [0.15 0.25]
 [0.1  0.2 ]]

Product of marginals:
 [[0.105 0.195]
 [0.14  0.26 ]
 [0.105 0.195]]

Apakah sama?  False
\end{verbatim}

\subsubsection{Interpretasi}\label{interpretasi-18}

Jika joint distribution tidak sama dengan hasil kali marginals, maka
\(X\) dan \(Y\) tidak independen.

\subsection{Catatan penting}\label{catatan-penting-1}

\begin{itemize}
\tightlist
\item
  \textbf{Independen ⇒ korelasi nol} (dalam banyak kasus dengan momen
  terbatas),
\item
  tetapi \textbf{korelasi nol tidak selalu berarti independen}.
\end{itemize}

Nanti kita lihat contoh kecilnya.

\section{5.7 Bivariat Kontinu}\label{bivariat-kontinu}

\subsection{K --- Konteks}\label{k-konteks-35}

Banyak pasangan peubah acak lebih alami dimodelkan kontinu, misalnya: -
tinggi dan berat badan, - suhu dan konsumsi listrik, - waktu proses
tahap 1 dan tahap 2.

\subsection{M --- Model}\label{m-model-35}

Joint density \(f_{X,Y}(x,y)\) memberi kerapatan peluang pasangan
\((X,Y)\).

Peluang pada daerah \(A\) adalah: \[
P((X,Y)\in A)=\iint_A f_{X,Y}(x,y)\,dx\,dy
\]

\subsection{Q --- Questions}\label{q-questions-35}

\begin{enumerate}
\def\labelenumi{\arabic{enumi}.}
\tightlist
\item
  Bagaimana membayangkan distribusi bersama kontinu?
\item
  Bagaimana memperoleh marginal?
\item
  Bagaimana memvisualisasikannya?
\end{enumerate}

\subsection{A --- Apply}\label{a-apply-34}

Mari ambil contoh Normal bivariat.

\begin{Shaded}
\begin{Highlighting}[]
\NormalTok{mean }\OperatorTok{=}\NormalTok{ [}\DecValTok{0}\NormalTok{, }\DecValTok{0}\NormalTok{]}
\NormalTok{cov }\OperatorTok{=}\NormalTok{ [[}\DecValTok{1}\NormalTok{, }\FloatTok{0.8}\NormalTok{],}
\NormalTok{       [}\FloatTok{0.8}\NormalTok{, }\DecValTok{1}\NormalTok{]]}

\NormalTok{data }\OperatorTok{=}\NormalTok{ rng.multivariate\_normal(mean, cov, size}\OperatorTok{=}\DecValTok{5000}\NormalTok{)}
\NormalTok{X }\OperatorTok{=}\NormalTok{ data[:,}\DecValTok{0}\NormalTok{]}
\NormalTok{Y }\OperatorTok{=}\NormalTok{ data[:,}\DecValTok{1}\NormalTok{]}

\NormalTok{plt.figure(figsize}\OperatorTok{=}\NormalTok{(}\DecValTok{6}\NormalTok{,}\DecValTok{5}\NormalTok{))}
\NormalTok{plt.scatter(X, Y, alpha}\OperatorTok{=}\FloatTok{0.2}\NormalTok{, s}\OperatorTok{=}\DecValTok{8}\NormalTok{)}
\NormalTok{plt.xlabel(}\StringTok{"X"}\NormalTok{)}
\NormalTok{plt.ylabel(}\StringTok{"Y"}\NormalTok{)}
\NormalTok{plt.title(}\StringTok{"Scatterplot Normal bivariat"}\NormalTok{)}
\NormalTok{plt.grid(alpha}\OperatorTok{=}\FloatTok{0.3}\NormalTok{)}
\NormalTok{plt.show()}
\end{Highlighting}
\end{Shaded}

\pandocbounded{\includegraphics[keepaspectratio]{05-random-variable-multivariat-dan-fungsi_files/figure-pdf/cell-13-output-1.pdf}}

Marginalnya dapat dilihat dari histogram tiap komponen.

\begin{Shaded}
\begin{Highlighting}[]
\NormalTok{plt.figure(figsize}\OperatorTok{=}\NormalTok{(}\DecValTok{8}\NormalTok{,}\DecValTok{3}\NormalTok{))}

\NormalTok{plt.subplot(}\DecValTok{1}\NormalTok{,}\DecValTok{2}\NormalTok{,}\DecValTok{1}\NormalTok{)}
\NormalTok{plt.hist(X, bins}\OperatorTok{=}\DecValTok{40}\NormalTok{, density}\OperatorTok{=}\VariableTok{True}\NormalTok{, alpha}\OperatorTok{=}\FloatTok{0.7}\NormalTok{)}
\NormalTok{plt.title(}\StringTok{"Marginal X"}\NormalTok{)}
\NormalTok{plt.grid(alpha}\OperatorTok{=}\FloatTok{0.3}\NormalTok{)}

\NormalTok{plt.subplot(}\DecValTok{1}\NormalTok{,}\DecValTok{2}\NormalTok{,}\DecValTok{2}\NormalTok{)}
\NormalTok{plt.hist(Y, bins}\OperatorTok{=}\DecValTok{40}\NormalTok{, density}\OperatorTok{=}\VariableTok{True}\NormalTok{, alpha}\OperatorTok{=}\FloatTok{0.7}\NormalTok{)}
\NormalTok{plt.title(}\StringTok{"Marginal Y"}\NormalTok{)}
\NormalTok{plt.grid(alpha}\OperatorTok{=}\FloatTok{0.3}\NormalTok{)}

\NormalTok{plt.tight\_layout()}
\NormalTok{plt.show()}
\end{Highlighting}
\end{Shaded}

\pandocbounded{\includegraphics[keepaspectratio]{05-random-variable-multivariat-dan-fungsi_files/figure-pdf/cell-14-output-1.pdf}}

\section{5.8 Kovariansi}\label{kovariansi}

\subsection{K --- Konteks}\label{k-konteks-36}

Kita sering ingin tahu apakah dua peubah cenderung bergerak bersama.
Misalnya: - semakin banyak pelanggan, semakin besar pendapatan, -
semakin lama hujan, semakin kecil pengunjung, - semakin besar arus,
semakin besar daya.

\subsection{M --- Model}\label{m-model-36}

Kovariansi didefinisikan sebagai:

\[
\mathrm{Cov}(X,Y)=E[(X-E[X])(Y-E[Y])]
\]

Bentuk ekuivalennya: \[
\mathrm{Cov}(X,Y)=E[XY]-E[X]E[Y]
\]

\subsection{Q --- Questions}\label{q-questions-36}

\begin{enumerate}
\def\labelenumi{\arabic{enumi}.}
\tightlist
\item
  Apa arti kovariansi positif, negatif, atau nol?
\item
  Mengapa satuan kovariansi sulit ditafsirkan langsung?
\item
  Bagaimana menghitungnya dari data?
\end{enumerate}

\subsection{A --- Apply}\label{a-apply-35}

\begin{Shaded}
\begin{Highlighting}[]
\NormalTok{cov\_xy }\OperatorTok{=}\NormalTok{ np.cov(X, Y, ddof}\OperatorTok{=}\DecValTok{0}\NormalTok{)[}\DecValTok{0}\NormalTok{,}\DecValTok{1}\NormalTok{]}
\BuiltInTok{print}\NormalTok{(}\StringTok{"Cov(X,Y) ="}\NormalTok{, cov\_xy)}
\end{Highlighting}
\end{Shaded}

\begin{verbatim}
Cov(X,Y) = 0.7693398942915717
\end{verbatim}

\subsubsection{Interpretasi}\label{interpretasi-19}

\begin{itemize}
\tightlist
\item
  Kovariansi positif: saat \(X\) di atas rata-rata, \(Y\) cenderung juga
  di atas rata-rata.
\item
  Kovariansi negatif: saat \(X\) di atas rata-rata, \(Y\) cenderung di
  bawah rata-rata.
\item
  Kovariansi nol: tidak ada hubungan linear yang jelas, tetapi belum
  tentu independen.
\end{itemize}

\section{5.9 Korelasi}\label{korelasi}

\subsection{K --- Konteks}\label{k-konteks-37}

Kovariansi bagus, tetapi nilainya dipengaruhi satuan. Kita sering butuh
ukuran yang lebih mudah dibandingkan antar-kasus.

\subsection{M --- Model}\label{m-model-37}

Korelasi adalah versi ternormalisasi dari kovariansi:

\[
\rho_{X,Y} = \frac{\mathrm{Cov}(X,Y)}{\sigma_X \sigma_Y}
\]

Nilainya berada antara -1 dan 1.

\subsection{Q --- Questions}\label{q-questions-37}

\begin{enumerate}
\def\labelenumi{\arabic{enumi}.}
\tightlist
\item
  Apa arti korelasi mendekati 1, 0, atau -1?
\item
  Apa bedanya korelasi dan kovariansi?
\item
  Mengapa korelasi tidak otomatis berarti hubungan sebab-akibat?
\end{enumerate}

\subsection{A --- Apply}\label{a-apply-36}

\begin{Shaded}
\begin{Highlighting}[]
\NormalTok{corr\_xy }\OperatorTok{=}\NormalTok{ np.corrcoef(X, Y)[}\DecValTok{0}\NormalTok{,}\DecValTok{1}\NormalTok{]}
\BuiltInTok{print}\NormalTok{(}\StringTok{"Corr(X,Y) ="}\NormalTok{, corr\_xy)}
\end{Highlighting}
\end{Shaded}

\begin{verbatim}
Corr(X,Y) = 0.7855402048546878
\end{verbatim}

\subsubsection{Interpretasi}\label{interpretasi-20}

\begin{itemize}
\tightlist
\item
  \(\rho \approx 1\): hubungan linear positif kuat,
\item
  \(\rho \approx -1\): hubungan linear negatif kuat,
\item
  \(\rho \approx 0\): hubungan linear lemah atau tidak ada.
\end{itemize}

\subsection{Peringatan penting}\label{peringatan-penting}

Korelasi: - tidak berarti sebab-akibat, - tidak menangkap semua hubungan
nonlinear, - bisa nol meskipun dua peubah sangat bergantung secara
nonlinear.

\section{5.10 Contoh: Korelasi Nol tetapi Tidak
Independen}\label{contoh-korelasi-nol-tetapi-tidak-independen}

Ambil: - \(X \sim \mathrm{Uniform}(-1,1)\) - \(Y = X^2\)

Maka jelas \(Y\) ditentukan sepenuhnya oleh \(X\), jadi mereka
\textbf{tidak independen}. Namun korelasinya bisa nol.

\begin{Shaded}
\begin{Highlighting}[]
\NormalTok{rng }\OperatorTok{=}\NormalTok{ np.random.default\_rng(}\DecValTok{123}\NormalTok{)}
\NormalTok{X }\OperatorTok{=}\NormalTok{ rng.uniform(}\OperatorTok{{-}}\DecValTok{1}\NormalTok{, }\DecValTok{1}\NormalTok{, size}\OperatorTok{=}\DecValTok{100000}\NormalTok{)}
\NormalTok{Y }\OperatorTok{=}\NormalTok{ X}\OperatorTok{**}\DecValTok{2}

\BuiltInTok{print}\NormalTok{(}\StringTok{"Corr(X,Y) \textasciitilde{}"}\NormalTok{, np.corrcoef(X, Y)[}\DecValTok{0}\NormalTok{,}\DecValTok{1}\NormalTok{])}
\end{Highlighting}
\end{Shaded}

\begin{verbatim}
Corr(X,Y) ~ -0.0005791433843902291
\end{verbatim}

\begin{Shaded}
\begin{Highlighting}[]
\NormalTok{plt.figure(figsize}\OperatorTok{=}\NormalTok{(}\DecValTok{6}\NormalTok{,}\DecValTok{5}\NormalTok{))}
\NormalTok{plt.scatter(X, Y, alpha}\OperatorTok{=}\FloatTok{0.15}\NormalTok{, s}\OperatorTok{=}\DecValTok{8}\NormalTok{)}
\NormalTok{plt.xlabel(}\StringTok{"X"}\NormalTok{)}
\NormalTok{plt.ylabel(}\StringTok{"Y = X\^{}2"}\NormalTok{)}
\NormalTok{plt.title(}\StringTok{"Korelasi hampir nol, tetapi jelas tidak independen"}\NormalTok{)}
\NormalTok{plt.grid(alpha}\OperatorTok{=}\FloatTok{0.3}\NormalTok{)}
\NormalTok{plt.show()}
\end{Highlighting}
\end{Shaded}

\pandocbounded{\includegraphics[keepaspectratio]{05-random-variable-multivariat-dan-fungsi_files/figure-pdf/cell-18-output-1.pdf}}

\subsubsection{Pelajaran penting}\label{pelajaran-penting}

Korelasi nol \textbf{bukan} bukti independensi.

\section{5.11 Fungsi dari Random
Variable}\label{fungsi-dari-random-variable}

\subsection{K --- Konteks}\label{k-konteks-38}

Sering kali yang kita minati bukan random variable awal, tetapi
\textbf{fungsi} darinya.

Contoh: - jika arus \(I\) acak, maka daya \(P = I^2R\), - jika harga
\(H\) dan jumlah terjual \(Q\) acak, maka revenue \(R = HQ\), - jika dua
tahap layanan acak, total waktu \(T = X_1 + X_2\).

\subsection{M --- Model}\label{m-model-38}

Jika \(X\) adalah random variable dan \(Y=g(X)\), maka \(Y\) juga random
variable.

Untuk dua peubah, jika: \[
Z = g(X,Y)
\] maka \(Z\) juga random variable.

\subsection{Q --- Questions}\label{q-questions-38}

\begin{enumerate}
\def\labelenumi{\arabic{enumi}.}
\tightlist
\item
  Bagaimana distribusi berubah setelah transformasi?
\item
  Mengapa fungsi nonlinear bisa mengubah bentuk distribusi secara
  drastis?
\item
  Bagaimana Python membantu melihat perubahan ini?
\end{enumerate}

\subsection{A --- Apply}\label{a-apply-37}

Misalkan: \[
I \sim \mathcal{N}(10, 0.8^2), \qquad P = I^2R, \quad R=5
\]

\begin{Shaded}
\begin{Highlighting}[]
\NormalTok{rng }\OperatorTok{=}\NormalTok{ np.random.default\_rng(}\DecValTok{7}\NormalTok{)}
\NormalTok{I }\OperatorTok{=}\NormalTok{ rng.normal(loc}\OperatorTok{=}\DecValTok{10}\NormalTok{, scale}\OperatorTok{=}\FloatTok{0.8}\NormalTok{, size}\OperatorTok{=}\DecValTok{100000}\NormalTok{)}
\NormalTok{R }\OperatorTok{=} \DecValTok{5}
\NormalTok{P }\OperatorTok{=}\NormalTok{ (I}\OperatorTok{**}\DecValTok{2}\NormalTok{) }\OperatorTok{*}\NormalTok{ R}

\BuiltInTok{print}\NormalTok{(}\StringTok{"Mean I ="}\NormalTok{, I.mean(), }\StringTok{"Var I ="}\NormalTok{, I.var())}
\BuiltInTok{print}\NormalTok{(}\StringTok{"Mean P ="}\NormalTok{, P.mean(), }\StringTok{"Var P ="}\NormalTok{, P.var())}
\end{Highlighting}
\end{Shaded}

\begin{verbatim}
Mean I = 9.998938944730092 Var I = 0.6378216295132455
Mean P = 503.0830082497668 Var P = 6392.509470652347
\end{verbatim}

\begin{Shaded}
\begin{Highlighting}[]
\NormalTok{plt.figure(figsize}\OperatorTok{=}\NormalTok{(}\DecValTok{8}\NormalTok{,}\DecValTok{4}\NormalTok{))}
\NormalTok{plt.hist(I, bins}\OperatorTok{=}\DecValTok{60}\NormalTok{, density}\OperatorTok{=}\VariableTok{True}\NormalTok{, alpha}\OperatorTok{=}\FloatTok{0.6}\NormalTok{, label}\OperatorTok{=}\StringTok{"I"}\NormalTok{)}
\NormalTok{plt.hist(P, bins}\OperatorTok{=}\DecValTok{60}\NormalTok{, density}\OperatorTok{=}\VariableTok{True}\NormalTok{, alpha}\OperatorTok{=}\FloatTok{0.6}\NormalTok{, label}\OperatorTok{=}\StringTok{"P = I\^{}2 R"}\NormalTok{)}
\NormalTok{plt.xlabel(}\StringTok{"Nilai"}\NormalTok{)}
\NormalTok{plt.ylabel(}\StringTok{"Density"}\NormalTok{)}
\NormalTok{plt.title(}\StringTok{"Fungsi nonlinear mengubah bentuk distribusi"}\NormalTok{)}
\NormalTok{plt.legend()}
\NormalTok{plt.grid(alpha}\OperatorTok{=}\FloatTok{0.3}\NormalTok{)}
\NormalTok{plt.show()}
\end{Highlighting}
\end{Shaded}

\pandocbounded{\includegraphics[keepaspectratio]{05-random-variable-multivariat-dan-fungsi_files/figure-pdf/cell-20-output-1.pdf}}

\subsubsection{Interpretasi}\label{interpretasi-21}

Transformasi nonlinear dapat: - mengubah pusat, - mengubah varians, -
mengubah skewness, - mengubah peluang melewati batas aman.

\section{5.12 Fungsi dari Dua Random
Variable}\label{fungsi-dari-dua-random-variable}

\subsection{K --- Konteks}\label{k-konteks-39}

Banyak keputusan melibatkan gabungan dua random variable.

Contoh: - revenue \(R = P \times Q\), - total waktu \(T = X + Y\), -
keuntungan \(G = \text{revenue} - \text{cost}\).

\subsection{M --- Model}\label{m-model-39}

Jika \(X\) dan \(Y\) acak, maka: - \(X+Y\), - \(XY\), - \(aX+bY\), -
\(\max(X,Y)\), - dan fungsi lain dari keduanya juga random variable.

\subsection{Q --- Questions}\label{q-questions-39}

\begin{enumerate}
\def\labelenumi{\arabic{enumi}.}
\tightlist
\item
  Bagaimana rata-rata jumlah dua random variable?
\item
  Bagaimana varians jumlah dua random variable?
\item
  Kapan covariance ikut berperan?
\end{enumerate}

\subsection{A --- Apply}\label{a-apply-38}

Untuk dua random variable: \[
E[X+Y] = E[X] + E[Y]
\]

dan: \[
\mathrm{Var}(X+Y)=\mathrm{Var}(X)+\mathrm{Var}(Y)+2\mathrm{Cov}(X,Y)
\]

Mari verifikasi dengan simulasi.

\begin{Shaded}
\begin{Highlighting}[]
\NormalTok{rng }\OperatorTok{=}\NormalTok{ np.random.default\_rng(}\DecValTok{11}\NormalTok{)}
\NormalTok{mean }\OperatorTok{=}\NormalTok{ [}\DecValTok{10}\NormalTok{, }\DecValTok{20}\NormalTok{]}
\NormalTok{cov }\OperatorTok{=}\NormalTok{ [[}\DecValTok{4}\NormalTok{, }\DecValTok{3}\NormalTok{],}
\NormalTok{       [}\DecValTok{3}\NormalTok{, }\DecValTok{9}\NormalTok{]]}

\NormalTok{data }\OperatorTok{=}\NormalTok{ rng.multivariate\_normal(mean, cov, size}\OperatorTok{=}\DecValTok{200000}\NormalTok{)}
\NormalTok{X }\OperatorTok{=}\NormalTok{ data[:,}\DecValTok{0}\NormalTok{]}
\NormalTok{Y }\OperatorTok{=}\NormalTok{ data[:,}\DecValTok{1}\NormalTok{]}
\NormalTok{T }\OperatorTok{=}\NormalTok{ X }\OperatorTok{+}\NormalTok{ Y}

\NormalTok{lhs\_mean }\OperatorTok{=}\NormalTok{ T.mean()}
\NormalTok{rhs\_mean }\OperatorTok{=}\NormalTok{ X.mean() }\OperatorTok{+}\NormalTok{ Y.mean()}

\NormalTok{lhs\_var }\OperatorTok{=}\NormalTok{ T.var()}
\NormalTok{rhs\_var }\OperatorTok{=}\NormalTok{ X.var() }\OperatorTok{+}\NormalTok{ Y.var() }\OperatorTok{+} \DecValTok{2}\OperatorTok{*}\NormalTok{np.cov(X, Y, ddof}\OperatorTok{=}\DecValTok{0}\NormalTok{)[}\DecValTok{0}\NormalTok{,}\DecValTok{1}\NormalTok{]}

\BuiltInTok{print}\NormalTok{(}\StringTok{"Mean left  ="}\NormalTok{, lhs\_mean)}
\BuiltInTok{print}\NormalTok{(}\StringTok{"Mean right ="}\NormalTok{, rhs\_mean)}
\BuiltInTok{print}\NormalTok{(}\StringTok{"Var left   ="}\NormalTok{, lhs\_var)}
\BuiltInTok{print}\NormalTok{(}\StringTok{"Var right  ="}\NormalTok{, rhs\_var)}
\end{Highlighting}
\end{Shaded}

\begin{verbatim}
Mean left  = 29.99708921283331
Mean right = 29.99708921283331
Var left   = 18.908340698332005
Var right  = 18.908340698332005
\end{verbatim}

\subsubsection{Interpretasi}\label{interpretasi-22}

Jika dua peubah berkorelasi positif, variance jumlahnya lebih besar
daripada jika mereka independen. Ini sangat penting dalam: - risk
aggregation, - portfolio, - total demand, - total delay.

\section{5.13 KMQA Mini-Case 1 --- Tinggi dan Berat
Badan}\label{kmqa-mini-case-1-tinggi-dan-berat-badan}

\subsection{K --- Konteks}\label{k-konteks-40}

Kita ingin memahami hubungan tinggi badan dan berat badan.

\subsection{M --- Model}\label{m-model-40}

\((X,Y)\) = (tinggi, berat) dimodelkan sebagai random vector kontinu.

\subsection{Q --- Questions}\label{q-questions-40}

\begin{enumerate}
\def\labelenumi{\arabic{enumi}.}
\tightlist
\item
  Apakah keduanya berkorelasi?
\item
  Apakah korelasi tinggi berarti sebab-akibat?
\item
  Apa arti marginal masing-masing?
\end{enumerate}

\subsection{A --- Apply}\label{a-apply-39}

Gunakan scatterplot, covariance, dan correlation untuk melihat pola
awal. Interpretasi harus hati-hati: hubungan statistik tidak otomatis
berarti hubungan sebab-akibat.

\section{5.14 KMQA Mini-Case 2 --- Jumlah Pelanggan dan
Revenue}\label{kmqa-mini-case-2-jumlah-pelanggan-dan-revenue}

\subsection{K --- Konteks}\label{k-konteks-41}

Sebuah toko ingin memodelkan: - \(X\) = jumlah pelanggan, - \(Y\) =
revenue harian.

\subsection{M --- Model}\label{m-model-41}

Secara sederhana, \(Y\) sering meningkat saat \(X\) meningkat. Bisa ada
hubungan linear kasar atau model fungsi: \[
Y = aX + \epsilon
\]

\subsection{Q --- Questions}\label{q-questions-41}

\begin{enumerate}
\def\labelenumi{\arabic{enumi}.}
\tightlist
\item
  Seberapa kuat hubungan keduanya?
\item
  Apakah variance revenue hanya ditentukan oleh variance pelanggan?
\item
  Apa implikasinya terhadap staffing dan stok?
\end{enumerate}

\subsection{A --- Apply}\label{a-apply-40}

\begin{Shaded}
\begin{Highlighting}[]
\NormalTok{rng }\OperatorTok{=}\NormalTok{ np.random.default\_rng(}\DecValTok{2026}\NormalTok{)}
\NormalTok{X }\OperatorTok{=}\NormalTok{ rng.poisson(lam}\OperatorTok{=}\DecValTok{100}\NormalTok{, size}\OperatorTok{=}\DecValTok{5000}\NormalTok{)}
\NormalTok{eps }\OperatorTok{=}\NormalTok{ rng.normal(}\DecValTok{0}\NormalTok{, }\DecValTok{50}\NormalTok{, size}\OperatorTok{=}\DecValTok{5000}\NormalTok{)}
\NormalTok{Y }\OperatorTok{=} \DecValTok{20}\OperatorTok{*}\NormalTok{X }\OperatorTok{+}\NormalTok{ eps}

\BuiltInTok{print}\NormalTok{(}\StringTok{"Cov(X,Y) ="}\NormalTok{, np.cov(X,Y, ddof}\OperatorTok{=}\DecValTok{0}\NormalTok{)[}\DecValTok{0}\NormalTok{,}\DecValTok{1}\NormalTok{])}
\BuiltInTok{print}\NormalTok{(}\StringTok{"Corr(X,Y) ="}\NormalTok{, np.corrcoef(X,Y)[}\DecValTok{0}\NormalTok{,}\DecValTok{1}\NormalTok{])}

\NormalTok{plt.figure(figsize}\OperatorTok{=}\NormalTok{(}\DecValTok{6}\NormalTok{,}\DecValTok{5}\NormalTok{))}
\NormalTok{plt.scatter(X, Y, alpha}\OperatorTok{=}\FloatTok{0.2}\NormalTok{, s}\OperatorTok{=}\DecValTok{8}\NormalTok{)}
\NormalTok{plt.xlabel(}\StringTok{"Jumlah pelanggan"}\NormalTok{)}
\NormalTok{plt.ylabel(}\StringTok{"Revenue"}\NormalTok{)}
\NormalTok{plt.title(}\StringTok{"Pelanggan dan revenue"}\NormalTok{)}
\NormalTok{plt.grid(alpha}\OperatorTok{=}\FloatTok{0.3}\NormalTok{)}
\NormalTok{plt.show()}
\end{Highlighting}
\end{Shaded}

\begin{verbatim}
Cov(X,Y) = 1968.1527526251139
Corr(X,Y) = 0.9695696157524835
\end{verbatim}

\pandocbounded{\includegraphics[keepaspectratio]{05-random-variable-multivariat-dan-fungsi_files/figure-pdf/cell-22-output-2.pdf}}

\section{5.15 KMQA Mini-Case 3 --- Arus dan
Daya}\label{kmqa-mini-case-3-arus-dan-daya}

\subsection{K --- Konteks}\label{k-konteks-42}

Arus listrik \(I\) acak, daya: \[
P = I^2 R
\]

\subsection{M --- Model}\label{m-model-42}

Daya adalah fungsi nonlinear dari arus.

\subsection{Q --- Questions}\label{q-questions-42}

\begin{enumerate}
\def\labelenumi{\arabic{enumi}.}
\tightlist
\item
  Mengapa variasi kecil di arus bisa diperbesar di daya?
\item
  Bagaimana distribusi daya dibanding distribusi arus?
\item
  Seberapa sering daya melewati batas aman?
\end{enumerate}

\subsection{A --- Apply}\label{a-apply-41}

\begin{Shaded}
\begin{Highlighting}[]
\NormalTok{rng }\OperatorTok{=}\NormalTok{ np.random.default\_rng(}\DecValTok{77}\NormalTok{)}
\NormalTok{I }\OperatorTok{=}\NormalTok{ rng.normal(}\DecValTok{10}\NormalTok{, }\FloatTok{1.0}\NormalTok{, size}\OperatorTok{=}\DecValTok{100000}\NormalTok{)}
\NormalTok{R }\OperatorTok{=} \DecValTok{4}
\NormalTok{P }\OperatorTok{=}\NormalTok{ I}\OperatorTok{**}\DecValTok{2} \OperatorTok{*}\NormalTok{ R}

\NormalTok{threshold }\OperatorTok{=} \DecValTok{500}
\BuiltInTok{print}\NormalTok{(}\StringTok{"P(P \textgreater{} 500) ="}\NormalTok{, np.mean(P }\OperatorTok{\textgreater{}}\NormalTok{ threshold))}
\end{Highlighting}
\end{Shaded}

\begin{verbatim}
P(P > 500) = 0.11884
\end{verbatim}

\begin{Shaded}
\begin{Highlighting}[]
\NormalTok{plt.figure(figsize}\OperatorTok{=}\NormalTok{(}\DecValTok{8}\NormalTok{,}\DecValTok{4}\NormalTok{))}
\NormalTok{plt.hist(P, bins}\OperatorTok{=}\DecValTok{60}\NormalTok{, density}\OperatorTok{=}\VariableTok{True}\NormalTok{, alpha}\OperatorTok{=}\FloatTok{0.75}\NormalTok{)}
\NormalTok{plt.axvline(threshold, color}\OperatorTok{=}\StringTok{\textquotesingle{}red\textquotesingle{}}\NormalTok{, linestyle}\OperatorTok{=}\StringTok{\textquotesingle{}{-}{-}\textquotesingle{}}\NormalTok{, label}\OperatorTok{=}\StringTok{\textquotesingle{}Batas aman\textquotesingle{}}\NormalTok{)}
\NormalTok{plt.xlabel(}\StringTok{"Daya"}\NormalTok{)}
\NormalTok{plt.ylabel(}\StringTok{"Density"}\NormalTok{)}
\NormalTok{plt.title(}\StringTok{"Distribusi daya sebagai fungsi dari arus acak"}\NormalTok{)}
\NormalTok{plt.legend()}
\NormalTok{plt.grid(alpha}\OperatorTok{=}\FloatTok{0.3}\NormalTok{)}
\NormalTok{plt.show()}
\end{Highlighting}
\end{Shaded}

\pandocbounded{\includegraphics[keepaspectratio]{05-random-variable-multivariat-dan-fungsi_files/figure-pdf/cell-24-output-1.pdf}}

\section{5.16 Python Toolbox untuk Bab
Ini}\label{python-toolbox-untuk-bab-ini-1}

\begin{Shaded}
\begin{Highlighting}[]
\ImportTok{import}\NormalTok{ numpy }\ImportTok{as}\NormalTok{ np}
\ImportTok{import}\NormalTok{ matplotlib.pyplot }\ImportTok{as}\NormalTok{ plt}
\ImportTok{from}\NormalTok{ scipy }\ImportTok{import}\NormalTok{ stats}
\end{Highlighting}
\end{Shaded}

\subsection{Kovariansi dan korelasi}\label{kovariansi-dan-korelasi}

\begin{Shaded}
\begin{Highlighting}[]
\NormalTok{X }\OperatorTok{=}\NormalTok{ np.array([}\DecValTok{1}\NormalTok{,}\DecValTok{2}\NormalTok{,}\DecValTok{3}\NormalTok{,}\DecValTok{4}\NormalTok{])}
\NormalTok{Y }\OperatorTok{=}\NormalTok{ np.array([}\DecValTok{2}\NormalTok{,}\DecValTok{3}\NormalTok{,}\DecValTok{5}\NormalTok{,}\DecValTok{7}\NormalTok{])}

\NormalTok{np.cov(X, Y, ddof}\OperatorTok{=}\DecValTok{0}\NormalTok{), np.corrcoef(X, Y)}
\end{Highlighting}
\end{Shaded}

\begin{verbatim}
(array([[1.25  , 2.125 ],
        [2.125 , 3.6875]]),
 array([[1.        , 0.98977827],
        [0.98977827, 1.        ]]))
\end{verbatim}

\subsection{Simulasi multivariate
normal}\label{simulasi-multivariate-normal}

\begin{Shaded}
\begin{Highlighting}[]
\NormalTok{mean }\OperatorTok{=}\NormalTok{ [}\DecValTok{0}\NormalTok{,}\DecValTok{0}\NormalTok{]}
\NormalTok{cov }\OperatorTok{=}\NormalTok{ [[}\DecValTok{1}\NormalTok{,}\FloatTok{0.6}\NormalTok{],[}\FloatTok{0.6}\NormalTok{,}\DecValTok{1}\NormalTok{]]}
\NormalTok{sample }\OperatorTok{=}\NormalTok{ np.random.default\_rng(}\DecValTok{0}\NormalTok{).multivariate\_normal(mean, cov, size}\OperatorTok{=}\DecValTok{5}\NormalTok{)}
\NormalTok{sample}
\end{Highlighting}
\end{Shaded}

\begin{verbatim}
array([[-0.05337744, -0.17153562],
       [-0.61972419, -0.52589867],
       [ 0.31740703,  0.64082748],
       [-1.58988058, -0.74278561],
       [ 1.19535362,  0.06352624]])
\end{verbatim}

\subsection{Heatmap joint distribution
diskrit}\label{heatmap-joint-distribution-diskrit}

\begin{Shaded}
\begin{Highlighting}[]
\NormalTok{joint }\OperatorTok{=}\NormalTok{ np.array([[}\FloatTok{0.1}\NormalTok{, }\FloatTok{0.2}\NormalTok{],[}\FloatTok{0.3}\NormalTok{, }\FloatTok{0.4}\NormalTok{]])}
\NormalTok{plt.imshow(joint, cmap}\OperatorTok{=}\StringTok{\textquotesingle{}Blues\textquotesingle{}}\NormalTok{, origin}\OperatorTok{=}\StringTok{\textquotesingle{}lower\textquotesingle{}}\NormalTok{)}
\NormalTok{plt.colorbar()}
\NormalTok{plt.show()}
\end{Highlighting}
\end{Shaded}

\pandocbounded{\includegraphics[keepaspectratio]{05-random-variable-multivariat-dan-fungsi_files/figure-pdf/cell-28-output-1.pdf}}

\section{5.17 Kesalahan Umum yang Harus
Dihindari}\label{kesalahan-umum-yang-harus-dihindari-2}

\begin{enumerate}
\def\labelenumi{\arabic{enumi}.}
\item
  \textbf{Menganggap joint distribution hanya dua histogram ditempel
  bersama}\\
  Tidak. Joint distribution memuat struktur hubungan keduanya.
\item
  \textbf{Menganggap marginal sudah cukup untuk semua keputusan}\\
  Tidak. Kadang hubungan antar peubah justru inti persoalannya.
\item
  \textbf{Menganggap korelasi berarti sebab-akibat}\\
  Tidak. Korelasi hanya ukuran asosiasi linear.
\item
  \textbf{Menganggap korelasi nol berarti independen}\\
  Tidak selalu.
\item
  \textbf{Mengabaikan covariance saat menghitung variance jumlah}\\
  Ini kesalahan penting dalam banyak aplikasi risiko.
\item
  \textbf{Lupa bahwa fungsi dari random variable juga random variable}\\
  Dalam banyak desain sistem, justru peubah turunan itulah yang paling
  penting.
\end{enumerate}

\section{5.18 Menyimpulkan Bab Ini}\label{menyimpulkan-bab-ini-4}

Bab ini memperluas perspektif kita dari satu random variable menjadi
dunia yang lebih realistis: dunia dengan beberapa peubah acak yang bisa
saling terkait.

Kita telah melihat: - joint distribution, - marginal distribution, -
conditional distribution, - independensi, - kovariansi, - korelasi, -
fungsi dari satu random variable, - dan fungsi dari beberapa random
variable.

Ini adalah langkah penting menuju cara berpikir yang lebih matang. Dunia
nyata jarang sesederhana satu peubah acak. Justru keputusan yang baik
sering muncul ketika kita mampu melihat hubungan antarpeubah dan
memahami bagaimana ketidakpastian merambat dari input ke output.

\section{5.19 Ringkasan Poin Inti}\label{ringkasan-poin-inti-4}

\begin{enumerate}
\def\labelenumi{\arabic{enumi}.}
\tightlist
\item
  Banyak masalah nyata melibatkan \textbf{lebih dari satu random
  variable}.
\item
  \textbf{Joint distribution} menggambarkan perilaku pasangan peubah
  secara simultan.
\item
  \textbf{Marginal distribution} diperoleh dengan ``menjumlahkan'' atau
  ``mengintegralkan keluar'' peubah lain.
\item
  \textbf{Conditional distribution} adalah distribusi yang relevan
  ketika informasi tambahan tersedia.
\item
  \textbf{Independensi} adalah syarat kuat: joint = hasil kali
  marginals.
\item
  \textbf{Covariance} mengukur kecenderungan bergerak bersama.
\item
  \textbf{Correlation} menormalkan covariance agar lebih mudah
  ditafsirkan.
\item
  Korelasi nol \textbf{tidak selalu} berarti independen.
\item
  Fungsi dari random variable juga merupakan random variable.
\item
  Dalam banyak keputusan nyata, justru random variable turunan seperti
  revenue, total delay, daya, dan profit yang paling penting.
\end{enumerate}

\section{5.20 Latihan Bab 5}\label{latihan-bab-5}

\subsection{A. Konseptual}\label{a.-konseptual-4}

\begin{enumerate}
\def\labelenumi{\arabic{enumi}.}
\tightlist
\item
  Jelaskan perbedaan joint, marginal, dan conditional distribution.
\item
  Apa arti dua random variable independen?
\item
  Apa bedanya covariance dan correlation?
\item
  Mengapa korelasi nol tidak menjamin independensi?
\end{enumerate}

\subsection{B. Hitungan}\label{b.-hitungan-3}

\begin{enumerate}
\def\labelenumi{\arabic{enumi}.}
\tightlist
\item
  Diberi joint PMF tabel kecil, hitung marginal \(p_X\) dan \(p_Y\).
\item
  Hitung conditional distribution \(P(X=x \mid Y=y)\).
\item
  Jika diketahui
  \(E[X], E[Y], \mathrm{Var}(X), \mathrm{Var}(Y), \mathrm{Cov}(X,Y)\),
  hitung \(E[X+Y]\) dan \(\mathrm{Var}(X+Y)\).
\end{enumerate}

\subsection{C. Python}\label{c.-python-3}

\begin{enumerate}
\def\labelenumi{\arabic{enumi}.}
\tightlist
\item
  Simulasikan dua peubah acak yang berkorelasi positif lalu hitung
  covariance dan correlation.
\item
  Simulasikan contoh \(Y=X^2\) untuk menunjukkan korelasi nol tetapi
  tidak independen.
\item
  Simulasikan model revenue sebagai fungsi dari jumlah pelanggan.
\end{enumerate}

\subsection{D. Aplikatif}\label{d.-aplikatif-3}

\begin{enumerate}
\def\labelenumi{\arabic{enumi}.}
\tightlist
\item
  Dalam konteks klinik, peubah apa saja yang masuk akal dimodelkan
  bersama?
\item
  Dalam konteks manufaktur, mengapa jumlah cacat dan biaya garansi
  sebaiknya dipikirkan bersama?
\item
  Dalam konteks listrik, mengapa distribusi daya lebih penting daripada
  sekadar distribusi arus?
\end{enumerate}

\section{5.21 Penutup Kecil}\label{penutup-kecil-4}

Dengan bab ini, Anda sudah punya salah satu kemampuan yang sangat
penting dalam probabilitas: melihat bahwa ketidakpastian jarang berdiri
sendiri. Ia biasanya datang berpasangan, berkelompok, saling terkait,
dan sering kali menghasilkan peubah baru yang lebih relevan bagi
keputusan.

Di bab terakhir, kita akan menutup perjalanan ini dengan merangkum
prinsip-prinsip utamanya: bagaimana random variable, distribusi, Python,
dan pengambilan keputusan menyatu menjadi satu cara berpikir yang kuat.

\bookmarksetup{startatroot}

\chapter{Bab 6. Penutup}\label{bab-6.-penutup}

\section{Tujuan Bab}\label{tujuan-bab-5}

Bab ini bertujuan untuk:

\begin{enumerate}
\def\labelenumi{\arabic{enumi}.}
\tightlist
\item
  merangkum gagasan utama yang telah dipelajari sepanjang buku,
\item
  menegaskan kembali hubungan antara \textbf{random variable},
  \textbf{probabilitas}, \textbf{Python}, dan \textbf{pengambilan
  keputusan},
\item
  membantu mahasiswa melihat struktur besar dari seluruh materi,
\item
  menutup perjalanan belajar ini dengan nada yang menguatkan, realistis,
  dan inspiratif.
\end{enumerate}

\section{Pembuka}\label{pembuka-5}

Kalau Anda sudah sampai pada bab ini, berarti Anda telah menempuh sebuah
perjalanan intelektual yang cukup penting.

Di awal, mungkin probabilitas dan statistika tampak seperti kumpulan
rumus yang asing. Ada simbol, distribusi, integral, penjumlahan, dan
notasi yang mudah terasa jauh dari kehidupan nyata. Tetapi sepanjang
buku ini, kita berusaha menempuh jalan yang berbeda. Kita tidak memulai
dari rumus demi rumus. Kita memulai dari \textbf{pertanyaan nyata}:

\begin{itemize}
\tightlist
\item
  bagaimana memilih keputusan ketika masa depan belum pasti?
\item
  bagaimana membandingkan pilihan yang sama-sama mungkin untung?
\item
  bagaimana menilai risiko?
\item
  bagaimana memodelkan jumlah, waktu, umur hidup, profit, dan kegagalan?
\item
  bagaimana menggunakan Python untuk melihat, mencoba, dan memahami?
\end{itemize}

Bab ini adalah tempat kita berhenti sejenak, melihat ke belakang, lalu
menyadari apa yang sebenarnya sudah kita pelajari.

\section{6.1 Ide Besar Buku Ini}\label{ide-besar-buku-ini}

Ada satu ide besar yang mengikat seluruh buku ini:

\begin{quote}
\textbf{Probabilitas dan statistika adalah alat berpikir untuk membuat
keputusan di bawah ketidakpastian.}
\end{quote}

Di balik semua contoh, distribusi, simulasi, dan rumus, ada satu
kebutuhan yang sama: kita ingin bertindak dengan lebih bijaksana
meskipun kita tidak tahu masa depan secara pasti.

Kita belajar bahwa: - masa depan sering kali tidak pasti, - tetapi
ketidakpastian itu bukan alasan untuk menyerah, - justru dengan model
yang tepat, kita bisa berpikir lebih jernih.

Dan salah satu model terpenting yang kita pelajari adalah \textbf{random
variable}.

\section{6.2 Apa yang Sebenarnya Sudah Kita
Kuasai?}\label{apa-yang-sebenarnya-sudah-kita-kuasai}

Mari kita lihat ulang perjalanan kita.

\subsection{Dari ketidakpastian ke
angka}\label{dari-ketidakpastian-ke-angka}

Kita belajar bahwa random variable adalah cara memetakan dunia acak ke
bilangan. Ini adalah langkah yang sangat besar. Begitu hasil acak
dipetakan menjadi angka, kita bisa: - menghitung peluang, - mencari
rata-rata, - mengukur variasi, - membandingkan pilihan, - membangun
simulasi, - dan merancang keputusan.

\subsection{Dari angka ke distribusi}\label{dari-angka-ke-distribusi}

Kita belajar bahwa satu angka tidak cukup. Yang lebih penting adalah
\textbf{bagaimana peluang tersebar} pada nilai-nilai yang mungkin.
Itulah distribusi.

Distribusi membuat kita bertanya: - apa nilai yang mungkin? - seberapa
mungkin? - seberapa besar peluang berada di bawah batas tertentu? -
seberapa liar penyimpangannya?

\subsection{Dari mean ke risk}\label{dari-mean-ke-risk}

Kita belajar bahwa mean atau ekspektasi sangat penting, tetapi tidak
cukup. Dua pilihan dengan rata-rata sama belum tentu sama baik.
Variance, simpangan baku, dan bentuk distribusi ikut menentukan kualitas
keputusan.

\subsection{Dari satu peubah ke beberapa
peubah}\label{dari-satu-peubah-ke-beberapa-peubah}

Kita belajar bahwa dunia nyata jarang hanya punya satu peubah acak.
Sering kali kita perlu memikirkan pasangan atau kelompok peubah acak: -
pelanggan dan revenue, - tinggi dan berat, - arus dan daya, - jumlah
cacat dan biaya garansi.

Dari situ kita belajar joint distribution, marginal, conditional,
covariance, correlation, dan transformasi fungsi.

\subsection{Dari teori ke eksperimen}\label{dari-teori-ke-eksperimen}

Kita juga belajar bahwa Python bukan sekadar kalkulator. Python adalah
laboratorium kecil yang memungkinkan kita: - mencoba ide, - melihat
pola, - menjalankan simulasi Monte Carlo, - memverifikasi rumus, - dan
membangun intuisi.

\section{6.3 Ringkasan Per Bab}\label{ringkasan-per-bab}

\subsection{Bab 1 --- Pendahuluan: Pengambilan
Keputusan}\label{bab-1-pendahuluan-pengambilan-keputusan}

Bab pertama menanamkan pesan utama buku ini: probabilitas bukan sekadar
kumpulan rumus, tetapi cara berpikir saat masa depan belum pasti.
Melalui contoh investasi, produk cacat, garansi, fungsi random variable,
risiko bangkrut, dan klinik gigi, kita melihat bahwa probabilitas hidup
dalam keputusan nyata.

\subsection{Bab 2 --- Random Variable
Umum}\label{bab-2-random-variable-umum}

Bab kedua memberi fondasi formal: - random variable sebagai fungsi, -
range, - diskrit dan kontinu, - PMF, CDF, PDF, - ekspektasi, - varians,
- simpangan baku.

Inilah kerangka dasar yang menopang semua bab sesudahnya.

\subsection{Bab 3 --- Distribusi
Diskrit}\label{bab-3-distribusi-diskrit}

Bab ketiga mengenalkan keluarga distribusi diskrit: - uniform diskrit, -
Bernoulli, - Binomial, - Geometric, - Poisson, - serta distribusi kustom
dari histogram.

Di sini kita belajar memilih model untuk hitungan kejadian.

\subsection{Bab 4 --- Distribusi
Kontinu}\label{bab-4-distribusi-kontinu}

Bab keempat membawa kita ke dunia waktu, umur hidup, dan pengukuran: -
uniform kontinu, - normal, - gamma, - exponential, - erlang, - weibull,
- pareto, - chi-square, - serta hubungan antardistribusi.

Di sini kita melihat bahwa banyak fenomena dunia nyata lebih nyaman
dipandang kontinu.

\subsection{Bab 5 --- Multivariat dan Fungsi Random
Variable}\label{bab-5-multivariat-dan-fungsi-random-variable}

Bab kelima memperluas perspektif: - joint distribution, - marginal, -
conditional, - independensi, - covariance, - correlation, - fungsi
random variable.

Di sini kita melihat bahwa ketidakpastian sering datang bersama-sama dan
sering menghasilkan peubah turunan yang justru paling penting bagi
keputusan.

\section{6.4 Prinsip-Prinsip Inti yang Perlu
Diingat}\label{prinsip-prinsip-inti-yang-perlu-diingat}

Berikut adalah beberapa prinsip inti yang layak dibawa terus, bahkan
setelah buku ini selesai dibaca.

\subsection{1. Mulai dari konteks, bukan dari
rumus}\label{mulai-dari-konteks-bukan-dari-rumus}

Ketika menghadapi masalah baru, jangan buru-buru bertanya ``rumus apa
yang dipakai?''\\
Mulailah dari: - apa keputusannya, - apa yang tidak pasti, - apa random
variable-nya, - apa ukuran yang perlu dihitung.

\subsection{2. Model yang baik lebih penting daripada rumus yang
banyak}\label{model-yang-baik-lebih-penting-daripada-rumus-yang-banyak}

Banyak kesalahan datang bukan karena salah menghitung, tetapi karena
salah memodelkan. Distribusi yang salah akan menghasilkan keputusan yang
salah.

\subsection{3. Ekspektasi penting, tetapi variance juga
penting}\label{ekspektasi-penting-tetapi-variance-juga-penting}

Mean memberi arah, variance memberi ukuran ombaknya. Dalam banyak
masalah nyata, survival, safety, dan reliability sangat dipengaruhi
variance.

\subsection{4. PDF bukan peluang}\label{pdf-bukan-peluang}

Untuk random variable kontinu, peluang datang dari area, bukan dari
tinggi kurva.

\subsection{5. Korelasi bukan
sebab-akibat}\label{korelasi-bukan-sebab-akibat}

Dua peubah bisa berkorelasi tanpa salah satunya menyebabkan yang lain.
Statistik membantu melihat pola, tetapi interpretasi tetap membutuhkan
akal sehat dan konteks.

\subsection{6. Python membantu belajar, bukan menggantikan
berpikir}\label{python-membantu-belajar-bukan-menggantikan-berpikir}

Kode yang baik bukan jalan pintas untuk menghindari pemahaman.
Sebaliknya, Python memberi kita kesempatan untuk menguji pemahaman
secara aktif.

\subsection{7. Ketidakpastian bisa dimodelkan, meskipun tidak bisa
dihilangkan}\label{ketidakpastian-bisa-dimodelkan-meskipun-tidak-bisa-dihilangkan}

Kita tidak akan pernah menghapus ketidakpastian dari hidup dan sistem
nyata. Tetapi kita bisa menghadapinya dengan lebih terstruktur.

\section{6.5 Checklist Kompetensi}\label{checklist-kompetensi}

Berikut checklist sederhana. Jika Anda sudah mampu melakukan sebagian
besar hal di bawah ini, berarti Anda sudah menempuh langkah yang sangat
baik.

\subsection{Konseptual}\label{konseptual}

\begin{itemize}
\tightlist
\item[$\square$]
  Saya bisa menjelaskan apa itu random variable.
\item[$\square$]
  Saya bisa membedakan random variable diskrit dan kontinu.
\item[$\square$]
  Saya tahu perbedaan PMF, PDF, dan CDF.
\item[$\square$]
  Saya bisa menjelaskan ekspektasi dan variance dengan kata-kata
  sendiri.
\item[$\square$]
  Saya tahu arti joint, marginal, conditional, covariance, dan
  correlation.
\end{itemize}

\subsection{Pemodelan}\label{pemodelan}

\begin{itemize}
\tightlist
\item[$\square$]
  Saya bisa memilih distribusi diskrit yang tepat untuk konteks count
  data.
\item[$\square$]
  Saya bisa memilih distribusi kontinu yang masuk akal untuk waktu
  tunggu atau lifetime.
\item[$\square$]
  Saya bisa menjelaskan arti parameter suatu distribusi.
\item[$\square$]
  Saya bisa memodelkan output sebagai fungsi dari input acak.
\end{itemize}

\subsection{Komputasional}\label{komputasional}

\begin{itemize}
\tightlist
\item[$\square$]
  Saya bisa mensimulasikan random variable dengan Python.
\item[$\square$]
  Saya bisa menggambar histogram, PMF, PDF, atau CDF.
\item[$\square$]
  Saya bisa memakai \texttt{numpy}, \texttt{scipy.stats}, dan
  \texttt{matplotlib}.
\item[$\square$]
  Saya bisa membandingkan hasil simulasi dengan model teoritis.
\end{itemize}

\subsection{Keputusan}\label{keputusan}

\begin{itemize}
\tightlist
\item[$\square$]
  Saya bisa menggunakan peluang untuk menilai risiko.
\item[$\square$]
  Saya bisa menggunakan ekspektasi dan variance untuk membandingkan
  pilihan.
\item[$\square$]
  Saya bisa menjelaskan implikasi keputusan berdasarkan model
  probabilistik.
\end{itemize}

\section{6.6 Python sebagai Partner
Belajar}\label{python-sebagai-partner-belajar}

Sepanjang buku ini, Python hadir bukan sekadar sebagai aksesoris. Python
dipakai karena ia memberi beberapa keuntungan penting.

\subsection{1. Python memberi quick
wins}\label{python-memberi-quick-wins}

Mahasiswa tidak perlu menunggu lama untuk merasakan bahwa konsep
probabilitas benar-benar bekerja. Dengan beberapa baris kode, kita bisa:
- mensimulasikan eksperimen, - membuat histogram, - menghitung mean dan
variance, - membandingkan Monte Carlo dan teori.

\subsection{2. Python membantu membangun
intuisi}\label{python-membantu-membangun-intuisi}

Sebelum semua definisi terasa nyaman, visualisasi dan simulasi sudah
lebih dulu memberi ``rasa'' terhadap konsep.

\subsection{3. Python membantu mengecek
asumsi}\label{python-membantu-mengecek-asumsi}

Dalam praktik nyata, kita jarang tahu jawaban dengan pasti. Simulasi
membantu memeriksa apakah intuisi kita masuk akal.

\subsection{4. Python membuat belajar lebih
aktif}\label{python-membuat-belajar-lebih-aktif}

Daripada hanya membaca rumus, mahasiswa bisa mengubah parameter,
menjalankan ulang simulasi, dan melihat apa yang berubah. Ini membuat
belajar menjadi proses eksplorasi, bukan hanya konsumsi.

\section{6.7 Dari Probabilitas ke Cara Berpikir
Engineer}\label{dari-probabilitas-ke-cara-berpikir-engineer}

Buku ini tidak hanya ingin mengajarkan isi mata kuliah. Buku ini juga
ingin menanamkan kebiasaan berpikir yang lebih matang. Seorang engineer
atau data scientist yang baik biasanya berpikir seperti ini:

\begin{enumerate}
\def\labelenumi{\arabic{enumi}.}
\tightlist
\item
  \textbf{Pahami konteksnya}
\item
  \textbf{Tentukan apa yang acak}
\item
  \textbf{Bangun random variable yang relevan}
\item
  \textbf{Pilih model distribusi yang masuk akal}
\item
  \textbf{Hitung ukuran yang penting}
\item
  \textbf{Gunakan simulasi bila perlu}
\item
  \textbf{Ambil keputusan dengan sadar akan asumsi dan risiko}
\end{enumerate}

Ini adalah pola berpikir yang jauh melampaui mata kuliah ini. Pola ini
akan sangat berguna dalam: - machine learning, - analitik data, -
reliability engineering, - keuangan, - quality control, - sistem
antrian, - dan banyak bidang lain.

\section{6.8 Sebuah Catatan Tentang
Ketidakpastian}\label{sebuah-catatan-tentang-ketidakpastian}

Ada satu pelajaran yang lebih filosofis di balik seluruh probabilitas.

Kita hidup di dunia yang tidak sepenuhnya pasti. Kita sering harus
bertindak tanpa mengetahui semua hal. Kadang hasilnya baik, kadang
tidak. Probabilitas tidak memberi kita kuasa untuk mengendalikan masa
depan sepenuhnya. Tetapi probabilitas memberi kita sesuatu yang sangat
penting: \textbf{kerendahan hati dan keberanian}.

\begin{itemize}
\tightlist
\item
  Rendah hati, karena kita tahu bahwa model kita tetap model, bukan
  kenyataan itu sendiri.
\item
  Berani, karena walaupun tidak tahu segalanya, kita tetap bisa membuat
  keputusan yang lebih baik daripada sekadar menebak.
\end{itemize}

Dalam arti itu, probabilitas adalah latihan intelektual sekaligus
latihan karakter.

\section{6.9 Arah Belajar Lanjut}\label{arah-belajar-lanjut}

Setelah random variable, ada banyak topik lanjutan yang menunggu.
Misalnya:

\subsection{1. Sampling distribution}\label{sampling-distribution}

Bagaimana statistik dari sampel juga menjadi random variable.

\subsection{2. Estimasi}\label{estimasi}

Bagaimana menebak parameter populasi dari data.

\subsection{3. Uji hipotesis}\label{uji-hipotesis}

Bagaimana membuat keputusan formal dari data.

\subsection{4. Inferensi Bayesian}\label{inferensi-bayesian}

Bagaimana memperbarui keyakinan saat data baru datang.

\subsection{5. Proses stokastik}\label{proses-stokastik}

Bagaimana memodelkan ketidakpastian yang berkembang terhadap waktu.

\subsection{6. Machine learning
probabilistik}\label{machine-learning-probabilistik}

Bagaimana ide probabilitas dipakai dalam klasifikasi, prediksi, dan
generative modeling.

Dengan kata lain, random variable bukanlah akhir. Ia adalah pintu masuk
ke wilayah yang jauh lebih luas.

\section{6.10 Sebuah Simulasi Penutup
Kecil}\label{sebuah-simulasi-penutup-kecil}

Untuk menutup, mari kembali ke ide paling dasar: jika dunia acak, kita
tidak harus menyerah. Kita bisa mensimulasikan, melihat pola, lalu
memahami.

Misalkan kita bandingkan rata-rata dari sampel acak Normal saat ukuran
sampel membesar.

\begin{Shaded}
\begin{Highlighting}[]
\ImportTok{import}\NormalTok{ numpy }\ImportTok{as}\NormalTok{ np}
\ImportTok{import}\NormalTok{ matplotlib.pyplot }\ImportTok{as}\NormalTok{ plt}

\NormalTok{rng }\OperatorTok{=}\NormalTok{ np.random.default\_rng(}\DecValTok{2026}\NormalTok{)}

\NormalTok{true\_mu }\OperatorTok{=} \DecValTok{100}
\NormalTok{true\_sigma }\OperatorTok{=} \DecValTok{15}

\NormalTok{sample\_sizes }\OperatorTok{=}\NormalTok{ [}\DecValTok{5}\NormalTok{, }\DecValTok{10}\NormalTok{, }\DecValTok{30}\NormalTok{, }\DecValTok{100}\NormalTok{, }\DecValTok{500}\NormalTok{]}
\NormalTok{means }\OperatorTok{=}\NormalTok{ []}

\ControlFlowTok{for}\NormalTok{ n }\KeywordTok{in}\NormalTok{ sample\_sizes:}
\NormalTok{    sample }\OperatorTok{=}\NormalTok{ rng.normal(loc}\OperatorTok{=}\NormalTok{true\_mu, scale}\OperatorTok{=}\NormalTok{true\_sigma, size}\OperatorTok{=}\NormalTok{n)}
\NormalTok{    means.append(sample.mean())}

\BuiltInTok{print}\NormalTok{(}\BuiltInTok{list}\NormalTok{(}\BuiltInTok{zip}\NormalTok{(sample\_sizes, means)))}
\end{Highlighting}
\end{Shaded}

\begin{verbatim}
[(5, np.float64(98.75556672945466)), (10, np.float64(100.67853136094531)), (30, np.float64(101.5093331179185)), (100, np.float64(100.4519454787968)), (500, np.float64(100.66502613594935))]
\end{verbatim}

\begin{Shaded}
\begin{Highlighting}[]
\NormalTok{plt.figure(figsize}\OperatorTok{=}\NormalTok{(}\DecValTok{8}\NormalTok{,}\DecValTok{4}\NormalTok{))}
\NormalTok{plt.plot(sample\_sizes, means, marker}\OperatorTok{=}\StringTok{\textquotesingle{}o\textquotesingle{}}\NormalTok{, label}\OperatorTok{=}\StringTok{\textquotesingle{}Sample mean\textquotesingle{}}\NormalTok{)}
\NormalTok{plt.axhline(true\_mu, color}\OperatorTok{=}\StringTok{\textquotesingle{}red\textquotesingle{}}\NormalTok{, linestyle}\OperatorTok{=}\StringTok{\textquotesingle{}{-}{-}\textquotesingle{}}\NormalTok{, label}\OperatorTok{=}\StringTok{\textquotesingle{}True mean\textquotesingle{}}\NormalTok{)}
\NormalTok{plt.xlabel(}\StringTok{"Ukuran sampel"}\NormalTok{)}
\NormalTok{plt.ylabel(}\StringTok{"Rata{-}rata sampel"}\NormalTok{)}
\NormalTok{plt.title(}\StringTok{"Semakin banyak data, estimasi rata{-}rata cenderung stabil"}\NormalTok{)}
\NormalTok{plt.legend()}
\NormalTok{plt.grid(alpha}\OperatorTok{=}\FloatTok{0.3}\NormalTok{)}
\NormalTok{plt.show()}
\end{Highlighting}
\end{Shaded}

\pandocbounded{\includegraphics[keepaspectratio]{06-penutup_files/figure-pdf/cell-3-output-1.pdf}}

Simulasi kecil ini mengingatkan kita pada sesuatu yang sederhana tetapi
penting: - ketidakpastian tidak hilang, - tetapi dengan model, data, dan
perhitungan, kita bisa bergerak menuju pemahaman yang lebih stabil.

\section{6.11 Ringkasan Besar Buku}\label{ringkasan-besar-buku}

Kalau seluruh buku ini harus diringkas dalam beberapa kalimat, mungkin
bunyinya seperti ini:

\begin{enumerate}
\def\labelenumi{\arabic{enumi}.}
\tightlist
\item
  Dunia nyata penuh ketidakpastian.
\item
  Random variable adalah cara memetakan ketidakpastian itu ke bilangan.
\item
  Distribusi menjelaskan bagaimana peluang tersebar.
\item
  Ekspektasi, variance, dan ukuran lain membantu kita menilai pilihan.
\item
  Banyak keputusan nyata memerlukan lebih dari satu peubah acak.
\item
  Python membuat konsep-konsep ini bisa dilihat, dicoba, dan diuji.
\item
  Tujuan akhir probabilitas bukan sekadar menjawab soal, tetapi membantu
  kita \textbf{berpikir lebih jernih saat harus memilih}.
\end{enumerate}

\section{6.12 Pesan untuk Mahasiswa}\label{pesan-untuk-mahasiswa}

Kalau Anda pernah merasa probabilitas itu sulit, itu wajar. Banyak orang
merasakannya. Tetapi sulit bukan berarti mustahil. Sering kali yang
dibutuhkan bukan hanya lebih banyak rumus, melainkan cara masuk yang
tepat: konteks yang hidup, contoh yang relevan, visualisasi yang jelas,
dan keberanian untuk mencoba.

Semoga buku ini sudah memberi jalan masuk itu.

Yang lebih penting lagi: semoga setelah mempelajari buku ini, Anda tidak
hanya menjadi mahasiswa yang bisa menghitung peluang, tetapi juga
seseorang yang lebih cermat dalam berpikir, lebih jujur terhadap asumsi,
dan lebih bertanggung jawab saat mengambil keputusan.

\section{6.13 Ringkasan Poin Inti}\label{ringkasan-poin-inti-5}

\begin{enumerate}
\def\labelenumi{\arabic{enumi}.}
\tightlist
\item
  Probabilitas adalah alat berpikir di bawah ketidakpastian.
\item
  Random variable adalah fondasi utama untuk memodelkan dunia acak.
\item
  Distribusi, mean, variance, dan CDF membantu kita membuat keputusan
  yang lebih baik.
\item
  Banyak persoalan nyata melibatkan joint behavior beberapa peubah acak.
\item
  Fungsi dari peubah acak sering kali justru menjadi besaran yang paling
  relevan.
\item
  Python membantu mengubah konsep menjadi pengalaman belajar yang
  konkret.
\item
  Tujuan akhir belajar probabilitas bukan sekadar lulus ujian, tetapi
  bertumbuh dalam cara berpikir.
\end{enumerate}

\section{6.14 Latihan Penutup}\label{latihan-penutup}

\subsection{Reflektif}\label{reflektif}

\begin{enumerate}
\def\labelenumi{\arabic{enumi}.}
\tightlist
\item
  Konsep apa yang paling mengubah cara Anda melihat probabilitas?
\item
  Kasus mana dalam buku ini yang paling terasa relevan bagi Anda?
\item
  Apa hubungan antara probabilitas dan tanggung jawab dalam mengambil
  keputusan?
\end{enumerate}

\subsection{Konseptual}\label{konseptual-1}

\begin{enumerate}
\def\labelenumi{\arabic{enumi}.}
\tightlist
\item
  Mengapa random variable menjadi jembatan antara dunia acak dan
  matematika?
\item
  Mengapa mean saja sering tidak cukup?
\item
  Apa perbedaan penting antara korelasi dan independensi?
\end{enumerate}

\subsection{Komputasional}\label{komputasional-1}

\begin{enumerate}
\def\labelenumi{\arabic{enumi}.}
\tightlist
\item
  Pilih satu distribusi diskrit dan satu distribusi kontinu dari buku
  ini, lalu buat simulasi dan visualisasi ulang dengan parameter pilihan
  Anda sendiri.
\item
  Buat satu mini-case baru dan modelkan dengan random variable yang
  sesuai.
\end{enumerate}

\section{6.15 Penutup Akhir}\label{penutup-akhir}

Buku ini sendiri lahir dari sebuah keputusan: keputusan untuk
mengajarkan probabilitas dan statistika bukan sebagai tumpukan simbol
yang dingin, tetapi sebagai ilmu yang dekat dengan pilihan nyata, logika
yang hidup, dan eksplorasi yang menyenangkan.

Kalau setelah perjalanan ini Anda mulai melihat dunia dengan pertanyaan
seperti: - ``apa random variable-nya?'' - ``apa distribusinya?'' -
``berapa mean dan variance-nya?'' - ``apa risiko dari keputusan ini?'' -
``bisakah saya simulasikan dulu di Python?''

maka buku ini sudah mencapai tujuannya.

Selamat melanjutkan perjalanan belajar. Ketidakpastian tidak akan pernah
hilang dari dunia nyata. Tetapi sekarang Anda sudah memiliki bahasa,
alat, dan cara berpikir yang lebih kuat untuk menghadapinya.




\end{document}
